% La reconstruction/renaissance de l’Europe et les révolutions industrielles — Histoire 1re Tech — Cours signé OnBuch+
% Programme officiel MINESEC. Compiler avec : tectonic N01-reconstruction-renaissance-europe-revolutions-industrie.tex
\newcommand{\DOCMATIERE}{Histoire}
\newcommand{\DOCNIVEAU}{1re Tech}
\newcommand{\DOCMODULE}{Histoire – classe de Première technique}
\newcommand{\DOCLECON}{N01}
\newcommand{\DOCTITRE}{La reconstruction/renaissance de l’Europe et les révolutions industrielles}
% OnBuch+ — préambule « cours signés » ENRICHI (compilé avec tectonic / XeLaTeX)
% Système visuel « Pop » d'OnBuch+ : fond crème (jamais blanc), bordures encre
% épaisses, ombres « dures » décalées, titres Archivo Black en capitales.
% Version MATHÉMATIQUES : graphes (pgfplots/tikz), boîtes pédagogiques
% (définition, propriété, méthode, attention, le savais-tu, exercice résolu,
% exercice, TP, bilan), page de garde et signature OnBuch+ ancrée en bas.
%
% Macros attendues dans main.tex AVANT \input{preamble} :
%   \DOCMATIERE  \DOCNIVEAU  \DOCMODULE  \DOCLECON (numéro)  \DOCTITRE
\documentclass[11pt]{article}

\usepackage{fontspec}
\usepackage[french]{babel}
\usepackage[a4paper,margin=2cm,top=3.4cm,bottom=2.8cm,headheight=1.4cm,footskip=1.3cm]{geometry}
\usepackage{amsmath,amssymb,amsfonts}
\usepackage{mathtools} % \xmapsto, \coloneqq... souvent utilisés par les modèles
\usepackage{cancel} % simplifications barrées (bilans de forces)
% \abs{z} (module, valeur absolue) et \norm{u} (norme), utilisés par les modèles
\providecommand{\abs}{}\let\abs\relax\DeclarePairedDelimiter{\abs}{\lvert}{\rvert}
\providecommand{\norm}{}\let\norm\relax\DeclarePairedDelimiter{\norm}{\lVert}{\rVert}
\usepackage[table]{xcolor} % [table] : fournit aussi \rowcolors (alternance de lignes)
\usepackage{pagecolor}
\usepackage{tikz}
\usetikzlibrary{arrows.meta,positioning,calc,shapes.geometric,shapes.misc,shapes.arrows,decorations.pathmorphing,decorations.pathreplacing,decorations.markings,patterns,shadows,mindmap,trees,fit,backgrounds,matrix,intersections,angles,quotes}
\usepackage{pgfplots}
\usepackage[version=4]{mhchem} % équations nucléaires \ce{^{235}_{92}U}
\usepackage[european,siunitx]{circuitikz} % schémas électriques
\pgfplotsset{compat=1.18}
\usepgfplotslibrary{fillbetween,groupplots} % aires entre courbes (calcul intégral)
\usepackage{enumitem}
\usepackage{fancyhdr}
% [most] charge toutes les bibliothèques tcolorbox (listings, minted, raster,
% documentation...) et gonfle inutilement la mémoire TeX sur un document long
% (~300 boîtes) : on ne charge que ce qui est réellement utilisé.
\usepackage[skins,breakable]{tcolorbox}
\usepackage{graphicx}
\usepackage{colortbl}
\usepackage{booktabs}
\usepackage{tabularx}
\usepackage{multirow}
\usepackage{diagbox} % case d'en-tête diagonale des tableaux à double entrée
% Colonnes extensibles centrée / gauche / droite (C, L, R), souvent utilisées par les modèles
\newcolumntype{C}{>{\centering\arraybackslash}X}
\newcolumntype{L}{>{\raggedright\arraybackslash}X}
\newcolumntype{R}{>{\raggedleft\arraybackslash}X}
\usepackage{array}
\usepackage{multicol}
\usepackage{siunitx}
% parse-numbers=false : accepte aussi \SI{2,5 \times 10^{-3}}{...}
\sisetup{locale=FR,output-decimal-marker={,},group-minimum-digits=4,parse-numbers=false}
\DeclareSIUnit\ohm{\ensuremath{\symup{\Omega}}} % U+2126 absent de la police
\DeclareSIUnit\division{div} % réglages d'oscilloscope (ms/div, V/div)
\DeclareSIUnit\tr{tr} % tours (vitesse de rotation, ex. \SI{1500}{\tr\per\minute})
\DeclareSIUnit\molar{M} % concentration molaire (ex. \SI{3}{\molar}, \SI{1}{\micro\molar})
\DeclareSIUnit\an{an} % année (le modèle confond parfois avec \year, un compteur TeX) — ex. \SI{5}{\kilo\meter\per\an}
\DeclareSIUnit\FCFA{FCFA} % franc CFA (ex. \SI{500}{\FCFA\per\kilo\gram})
\DeclareSIUnit\degreeBrix{\textdegree Brix} % teneur en sucre d'un sirop/jus (réfractométrie)
\DeclareSIUnit\month{mois} % le modèle utilise parfois \month, un compteur TeX, comme unité de durée
\DeclareSIUnit\microliter{\micro\liter} % le modèle écrit parfois \microliter en un seul token
\DeclareSIUnit\million{millions} % dénombrement (ex. \SI{15}{\million\per\mL} spermatozoïdes)
\usepackage{qrcode}
% \degree : commande fréquemment utilisée par les modèles pour un angle ou une
% température en dehors de \SI{}{} (non fournie par les paquets chargés).
\providecommand{\degree}{^{\circ}}
% \qed (fin de démonstration), utilisé par les modèles sans amsthm
\providecommand{\qed}{\hfill$\square$}
% \barre : double barre des valeurs interdites dans un tableau de variations (tkz-tab)
\providecommand{\barre}{\ensuremath{\|}}
% environnement proof (amsthm non chargé) utilisé par les modèles
\ifdefined\proof\else\newenvironment{proof}[1][Démonstration]{\par\noindent\textit{#1.}\ }{\qed\par}\fi
\usepackage{lastpage}
\usepackage{hyperref}
\hypersetup{hidelinks}

% ============================================================
% Palette « Pop » OnBuch+ — mêmes hex que la classe P (plus_ui.dart)
% ============================================================
\definecolor{poporange}{HTML}{F59321}
\definecolor{popdark}{HTML}{E07A0C}
\definecolor{popcream}{HTML}{FFF6E8}
\definecolor{popink}{HTML}{1C1714}
\definecolor{poppurple}{HTML}{7A5AE0}
\definecolor{popgreen}{HTML}{2E9E6A}
\definecolor{poppink}{HTML}{E0457B}
\definecolor{popblue}{HTML}{2F7DE1}
\definecolor{popgold}{HTML}{C98A12}
\definecolor{popmuted}{HTML}{8A7F76}
\definecolor{popline}{HTML}{EADFD2}
\definecolor{popgray}{HTML}{8A7F76}
\colorlet{popgrey}{popgray}
% Alias tolérants (le modèle nomme parfois les couleurs différemment)
\colorlet{popwhite}{white}
\colorlet{popblack}{popink}
\colorlet{popred}{poppink}
\colorlet{popyellow}{popgold}
% Teintes claires (fonds de tableaux/encadrés) — le modèle nomme parfois
% les couleurs avec un suffixe L (ex. popblueL) sans les avoir définies.
\colorlet{poporangeL}{poporange!20}
\colorlet{popdarkL}{popdark!20}
\colorlet{poppurpleL}{poppurple!20}
\colorlet{popgreenL}{popgreen!20}
\colorlet{poppinkL}{poppink!20}
\colorlet{popblueL}{popblue!20}
\colorlet{popgoldL}{popgold!20}
\colorlet{popredL}{popred!20}
\colorlet{popgrayL}{popgray!20}
% Teintes claires pour fonds de tableaux / zones
\colorlet{popblueL}{popblue!10!white}
\colorlet{poporangeL}{poporange!14!white}
\colorlet{popgreenL}{popgreen!12!white}
\colorlet{poppurpleL}{poppurple!12!white}
\colorlet{poppinkL}{poppink!12!white}
\colorlet{popgoldL}{popgold!14!white}

\pagecolor{popcream}

% ============================================================
% Typographie
% ============================================================
\defaultfontfeatures{Ligatures=TeX}
\setmainfont{PlusJakartaSans-Medium.ttf}[Path=../../../fonts/,BoldFont=PlusJakartaSans-Bold.ttf]
% Maths en sans-serif (Fira Math) avec chiffres et lettres droites de Plus
% Jakarta, pour que les formules se fondent dans le texte.
\usepackage[math-style=ISO,bold-style=ISO]{unicode-math}
\setmathfont{FiraMath-Regular.otf}[Path=../../../fonts/]
\setmathfont{PlusJakartaSans-Medium.ttf}[Path=../../../fonts/,range={up/num,up/{Latin,latin},\mathup}]
\usepackage{icomma} % virgule décimale sans espace en mode math
% Caractères mathématiques Unicode absents des polices de texte (Plus Jakarta,
% Archivo Black) : rendus par leur équivalent mathématique, en texte comme en
% formule — sinon ils s'affichent en carré vide (ex. « Arithmétique dans ℤ »).
\usepackage{newunicodechar}
\newunicodechar{ℝ}{\ensuremath{\mathbb{R}}}
\newunicodechar{ℤ}{\ensuremath{\mathbb{Z}}}
\newunicodechar{ℂ}{\ensuremath{\mathbb{C}}}
\newunicodechar{ℕ}{\ensuremath{\mathbb{N}}}
\newunicodechar{ℚ}{\ensuremath{\mathbb{Q}}}
\newunicodechar{𝔻}{\ensuremath{\mathbb{D}}}
\newunicodechar{α}{\ensuremath{\alpha}}
\newunicodechar{β}{\ensuremath{\beta}}
\newunicodechar{γ}{\ensuremath{\gamma}}
\newunicodechar{δ}{\ensuremath{\delta}}
\newunicodechar{ε}{\ensuremath{\varepsilon}}
\newunicodechar{θ}{\ensuremath{\theta}}
\newunicodechar{λ}{\ensuremath{\lambda}}
\newunicodechar{μ}{\ensuremath{\mu}}
\newunicodechar{σ}{\ensuremath{\sigma}}
\newunicodechar{τ}{\ensuremath{\tau}}
\newunicodechar{φ}{\ensuremath{\varphi}}
\newunicodechar{ψ}{\ensuremath{\psi}}
\newunicodechar{ω}{\ensuremath{\omega}}
\newunicodechar{Σ}{\ensuremath{\Sigma}}
% majuscules produites par \MakeUppercase dans les titres (α-aminés -> Α)
\newunicodechar{Α}{\ensuremath{\alpha}}
\newunicodechar{Β}{\ensuremath{\beta}}
\newunicodechar{Γ}{\ensuremath{\Gamma}}
\newunicodechar{Δ}{\ensuremath{\Delta}}
\newunicodechar{Ε}{\ensuremath{\varepsilon}}
\newunicodechar{Θ}{\ensuremath{\Theta}}
\newunicodechar{Λ}{\ensuremath{\Lambda}}
\newunicodechar{Μ}{\ensuremath{\mu}}
\newunicodechar{Τ}{\ensuremath{\tau}}
\newunicodechar{Φ}{\ensuremath{\Phi}}
\newunicodechar{Ψ}{\ensuremath{\Psi}}
\newunicodechar{Ω}{\ensuremath{\Omega}}
\newunicodechar{↦}{\ensuremath{\mapsto}}
\newunicodechar{∈}{\ensuremath{\in}}
\newunicodechar{∉}{\ensuremath{\notin}}
\newunicodechar{∧}{\ensuremath{\wedge}}
\newunicodechar{⇔}{\ensuremath{\Leftrightarrow}}
\newunicodechar{⟺}{\ensuremath{\Longleftrightarrow}}
\newunicodechar{⇒}{\ensuremath{\Rightarrow}}
\newunicodechar{⟹}{\ensuremath{\Longrightarrow}}
\newunicodechar{∘}{\ensuremath{\circ}}
\newunicodechar{∪}{\ensuremath{\cup}}
\newunicodechar{∩}{\ensuremath{\cap}}
\newunicodechar{≡}{\ensuremath{\equiv}}
\newunicodechar{⊂}{\ensuremath{\subset}}
\newunicodechar{⊥}{\ensuremath{\perp}}
\newunicodechar{∃}{\ensuremath{\exists}}
\newunicodechar{∀}{\ensuremath{\forall}}
\newunicodechar{∛}{\ensuremath{\sqrt[3]{\,}}}
\newunicodechar{∜}{\ensuremath{\sqrt[4]{\,}}}
\newunicodechar{⃗}{\ensuremath{{}^{\to}}}
\newunicodechar{ⁿ}{\ifmmode^{n}\else\textsuperscript{\ensuremath{n}}\fi}
\newunicodechar{ˣ}{\ifmmode^{x}\else\textsuperscript{\ensuremath{x}}\fi}
\newunicodechar{ᵇ}{\ifmmode^{b}\else\textsuperscript{\ensuremath{b}}\fi}
\newunicodechar{ʳ}{\ifmmode^{r}\else\textsuperscript{\ensuremath{r}}\fi}
\newunicodechar{ᵘ}{\ifmmode^{u}\else\textsuperscript{\ensuremath{u}}\fi}
\newunicodechar{ʸ}{\ifmmode^{y}\else\textsuperscript{\ensuremath{y}}\fi}
\newunicodechar{ᵃ}{\ifmmode^{a}\else\textsuperscript{\ensuremath{a}}\fi}
\newunicodechar{ᵉ}{\ifmmode^{e}\else\textsuperscript{\ensuremath{e}}\fi}
\newunicodechar{ᵏ}{\ifmmode^{k}\else\textsuperscript{\ensuremath{k}}\fi}
\newunicodechar{ᵖ}{\ifmmode^{p}\else\textsuperscript{\ensuremath{p}}\fi}
\newunicodechar{ᴺ}{\ifmmode^{N}\else\textsuperscript{\ensuremath{N}}\fi}
\newunicodechar{ᶜ}{\ifmmode^{c}\else\textsuperscript{\ensuremath{c}}\fi}
\newunicodechar{ʰ}{\ifmmode^{h}\else\textsuperscript{\ensuremath{h}}\fi}
\newunicodechar{ᵠ}{\ifmmode^{q}\else\textsuperscript{\ensuremath{q}}\fi}
\newunicodechar{ₙ}{\ifmmode_{n}\else\textsubscript{\ensuremath{n}}\fi}
\newunicodechar{ₐ}{\ifmmode_{a}\else\textsubscript{\ensuremath{a}}\fi}
\newunicodechar{ᵢ}{\ifmmode_{i}\else\textsubscript{\ensuremath{i}}\fi}
\newunicodechar{ᵣ}{\ifmmode_{r}\else\textsubscript{\ensuremath{r}}\fi}
\newunicodechar{ₖ}{\ifmmode_{k}\else\textsubscript{\ensuremath{k}}\fi}
\newunicodechar{ᵦ}{\ifmmode_{\beta}\else\textsubscript{\ensuremath{\beta}}\fi}
\newunicodechar{ₓ}{\ifmmode_{x}\else\textsubscript{\ensuremath{x}}\fi}
\newfontfamily\popdisplay{ArchivoBlack-Regular.ttf}[Path=../../../fonts/]
\newfontfamily\poplogo{SpaceGrotesk-Bold.ttf}[Path=../../../fonts/]
\newfontfamily\popsemi{PlusJakartaSans-SemiBold.ttf}[Path=../../../fonts/]
\newfontfamily\popxbold{PlusJakartaSans-ExtraBold.ttf}[Path=../../../fonts/]
\newfontfamily\popxboldls{PlusJakartaSans-ExtraBold.ttf}[Path=../../../fonts/,LetterSpace=3.0]

\setlist[enumerate]{leftmargin=1.7em,itemsep=5pt,topsep=4pt}
\setlist[itemize]{leftmargin=1.7em,itemsep=4pt,topsep=4pt,label={\color{poporange}\textbullet}}
\setlength{\parindent}{0pt}
\setlength{\parskip}{5pt}
\renewcommand{\arraystretch}{1.25}

% pgfplots : style Pop par défaut
\pgfplotsset{
  every axis/.append style={axis line style={popink,line width=1pt},tick style={popink},
    grid style={popline,line width=0.5pt},label style={font=\small},tick label style={font=\footnotesize}},
  popcurve/.style={poporange,line width=1.8pt,no markers,smooth},
}
% Même style disponible dans un \draw TikZ simple (hors pgfplots)
\tikzset{popcurve/.style={poporange,line width=1.8pt}}
\tikzset{popgrid/.style={popline,very thin}}
\colorlet{popcurve}{poporange} % le nom du style est parfois utilisé comme couleur

% ============================================================
% Pill / squiggle
% ============================================================
\newcommand{\poppill}[3]{%
  \tikz[baseline=-0.55ex]\node[draw=popink,line width=1.2pt,rounded corners=2.6pt,%
    fill=#2,inner xsep=6.5pt,inner ysep=3pt]{%
    \popxboldls\fontsize{7}{7}\selectfont\color{#3}\MakeUppercase{#1}};%
}
\newcommand{\popsquiggle}[1]{%
  \tikz[baseline=-0.4ex]{
    \draw[#1,line width=2.6pt,line cap=round]
      plot [smooth] coordinates {(0,0.09) (0.34,0.01) (0.68,0.21) (1.02,0.01) (1.36,0.10)};
  }%
}
% Mot-clé surligné (marqueur orange sous le texte)
\newcommand{\cle}[1]{{\popxbold\color{popdark}#1}}

% ============================================================
% En-tête / pied
% ============================================================
\pagestyle{fancy}
\fancyhf{}
\renewcommand{\headrulewidth}{0pt}
\renewcommand{\footrulewidth}{0pt}
\lhead{\raisebox{-0.15ex}{\poplogo\fontsize{13}{13}\selectfont\color{popink}OnBuch\color{poporange}+}}
\rhead{\poppill{\DOCMATIERE\ \textbullet\ \DOCNIVEAU\ \textbullet\ Leçon \DOCLECON}{white}{popink}}
\lfoot{\popxbold\fontsize{7.6}{7.6}\selectfont\color{popmuted}\MakeUppercase{Cours signé OnBuch+}\ \textbullet\ {\popsemi\DOCTITRE}}
\rfoot{\tikz[baseline=-0.6ex]\node[rounded corners=8.5pt,fill=popink,minimum height=17pt,minimum width=17pt,inner xsep=5pt,inner sep=0]{\popxbold\fontsize{8}{8}\selectfont\color{popcream}\thepage\,/\,\pageref*{LastPage}};}

% ============================================================
% Titres
% ============================================================
\newcommand{\chaptitre}[2]{%
  \begin{tcolorbox}[enhanced,colback=poporange,colframe=popink,boxrule=2.6pt,arc=11pt,
      drop shadow=popink,left=15pt,right=15pt,top=12pt,bottom=11pt,before skip=3pt,after skip=18pt]
    \poppill{#2}{popink}{popcream}\par\vspace{9pt}
    {\raggedright\hyphenpenalty=10000\exhyphenpenalty=10000\popdisplay\fontsize{22}{24}\selectfont\color{popcream}\MakeUppercase{#1}\par}
  \end{tcolorbox}
}
% Section numérotée : pastille orange avec le numéro + titre Archivo
\newcounter{popsec}
\newcommand{\coursec}[1]{%
  \stepcounter{popsec}\setcounter{popsub}{0}%
  \par\vspace{16pt}\noindent
  \begin{minipage}{\linewidth}
  \tikz[baseline=-0.7ex]\node[draw=popink,line width=1.6pt,fill=poporange,rounded corners=4pt,minimum size=22pt,inner sep=0]{\popdisplay\fontsize{12}{12}\selectfont\color{popcream}\thepopsec};%
  \hspace{8pt}{\popdisplay\fontsize{15}{17}\selectfont\color{popink}\MakeUppercase{#1}}\par
  \vspace{4pt}\noindent{\color{poporange}\rule{36pt}{3.6pt}}
  \end{minipage}\par\nopagebreak\vspace{8pt}
}
\newcounter{popsub}
\newcommand{\courssub}[1]{%
  \stepcounter{popsub}%
  \par\vspace{9pt}\noindent{\popxbold\fontsize{11.5}{13}\selectfont\color{popink}{\color{poporange}\thepopsec.\thepopsub}\hspace{6pt}#1}\par\nopagebreak\vspace{4pt}
}

% ============================================================
% Boîtes pédagogiques — même recette : fond blanc, bordure épaisse colorée,
% ombre dure encre, badge plein. Titre optionnel [#1].
% ============================================================
\tcbset{popbase/.style={breakable,enhanced,colback=white,boxrule=2.3pt,arc=9pt,
  shadow={3pt}{-3pt}{0pt}{popink},left=14pt,right=14pt,top=11pt,bottom=11pt,before skip=11pt,after skip=15pt,
  fontupper=\popsemi\fontsize{10.6}{14.5}\selectfont\color{popink},
  fontlower=\popsemi\fontsize{10.6}{14.5}\selectfont\color{popink}}}
% Coche dessinée (le glyphe ✓ n'existe pas dans les polices du document)
\newcommand{\popcheck}[1]{\tikz[baseline=-0.5ex]\draw[#1,line width=1.6pt,line cap=round,line join=round](-0.13,0.0)--(-0.04,-0.09)--(0.13,0.1);}
\providecommand{\coche}{\popcheck{popgreen}}
\providecommand{\resultat}[1]{\par\smallskip\noindent\fcolorbox{popgreen}{popgreenL}{\parbox{\dimexpr\linewidth-2\fboxsep-2\fboxrule\relax}{\textbf{Résultat.} #1}}\par\smallskip}
\newcommand{\popboxhead}[3]{\poppill{#1}{#2}{white}\ifx\relax#3\relax\else\hspace{6pt}{\popxbold\fontsize{10.6}{12}\selectfont\color{popink}#3}\fi\par\vspace{7pt}}

\newtcolorbox{aretenir}[1][]{popbase,colframe=popgreen,before upper={\popboxhead{À retenir}{popgreen}{#1}}}
\newtcolorbox{exemplebox}[1][]{popbase,colframe=poporange,before upper={\popboxhead{Exemple}{poporange}{#1}}}
\newtcolorbox{definition}[1][]{popbase,colframe=popblue,before upper={\popboxhead{Définition}{popblue}{#1}}}
\newtcolorbox{propriete}[1][]{popbase,colframe=poppurple,before upper={\popboxhead{Propriété}{poppurple}{#1}}}
\newtcolorbox{methode}[1][]{popbase,colframe=popgold,before upper={\popboxhead{Méthode}{popgold}{#1}}}
\newtcolorbox{attention}[1][]{popbase,colframe=poppink,before upper={\popboxhead{Attention}{poppink}{#1}}}
\newtcolorbox{experience}[1][]{popbase,colframe=popblue,colback=popblueL,before upper={\popboxhead{Expérience}{popblue}{#1}}}
% « Le savais-tu ? » : carte violette pleine (contexte, culture, Cameroun)
\newtcolorbox{savaistu}[1][]{popbase,colback=poppurple,colframe=popink,
  fontupper=\popsemi\fontsize{10.4}{14.2}\selectfont\color{white},
  before upper={\poppill{Le savais-tu\,?}{white}{poppurple}\ifx\relax#1\relax\else\hspace{6pt}{\popxbold\color{white}#1}\fi\par\vspace{7pt}}}
% Exercice résolu : énoncé en haut, \tcblower puis la solution sur fond crème
\newcounter{popexo}\newcounter{popexr}
\newtcolorbox{exoresolu}[1][]{popbase,colframe=poporange,colbacklower=popcream,
  segmentation style={popink,line width=1.2pt,dashed},
  before upper={\stepcounter{popexr}\popboxhead{Exercice résolu \thepopexr}{poporange}{#1}},
  before lower={\poppill{Solution}{popink}{popcream}\par\vspace{6pt}}}
% Exercice (non résolu en place) — la correction est dans la partie Corrigés
\newtcolorbox{exercice}[1][]{popbase,colframe=popink,
  before upper={\stepcounter{popexo}\popboxhead{Exercice \thepopexo}{popink}{#1}}}
% Corrigé
\newtcolorbox{corrige}[1][]{popbase,colframe=popgreen,colback=popgreenL,
  before upper={\popboxhead{Corrigé}{popgreen}{#1}}}
% Objectifs / prérequis / bilan
\newtcolorbox{objectifs}{popbase,colframe=popink,colback=poporangeL,
  before upper={\popboxhead{Objectifs de la leçon}{popink}{}}}
\newtcolorbox{prerequis}{popbase,colframe=popink,colback=popblueL,
  before upper={\popboxhead{Prérequis}{popblue}{}}}
\newtcolorbox{bilan}{popbase,colframe=popink,colback=white,boxrule=2.8pt,
  before upper={\popboxhead{Fiche bilan}{poporange}{L'essentiel à connaître pour l'examen}}}
% Figure encadrée avec légende
\newtcolorbox{popfigure}[1][]{enhanced,breakable=false,colback=white,colframe=popline,boxrule=1.4pt,arc=8pt,
  left=10pt,right=10pt,top=10pt,bottom=8pt,before skip=10pt,after skip=12pt,halign=center,
  lower separated=false,halign lower=center,
  fontlower=\popsemi\fontsize{9}{11.5}\selectfont\color{popmuted},
  #1}
\newcommand{\legende}[1]{\par\vspace{6pt}{\popsemi\fontsize{9}{11.5}\selectfont\color{popmuted}#1}}

% ============================================================
% Signature OnBuch+
% ============================================================
\newcommand{\poplogoinline}[1]{{\poplogo\fontsize{#1}{#1}\selectfont\color{popink}OnBuch\color{poporange}+}}
\newcommand{\signatureonbuch}{%
  \begin{tcolorbox}[enhanced,colback=popink,colframe=popink,boxrule=2.2pt,arc=10pt,
      drop shadow=poporange,left=14pt,right=14pt,top=10pt,bottom=10pt,before skip=0pt,after skip=0pt]
    \tikz[baseline=-0.7ex]\node[circle,fill=popgreen,draw=popcream,line width=1.3pt,minimum size=22pt,inner sep=0]{\popcheck{white}};%
    \hspace{9pt}\begin{minipage}[c]{0.62\linewidth}
      {\popxbold\fontsize{12}{13}\selectfont\color{popcream}Signé\ \textbullet\ {\poplogo OnBuch\color{poporange}+}}\par\vspace{2pt}
      {\raggedright\popsemi\fontsize{8}{10.5}\selectfont\color{popline}\MakeUppercase{Équipe pédagogique OnBuch+}\\\MakeUppercase{Conforme au programme officiel MINESEC}\par}
    \end{minipage}\hfill
    {\poplogo\fontsize{22}{22}\selectfont\color{popcream}OnBuch\color{poporange}+}
  \end{tcolorbox}%
}

% ============================================================
% Page de garde
% #1 titre · #2 matière · #3 niveau · #4 module · #5 n° leçon · #6 durée ·
% #7 description / objectifs (texte ou itemize)
% ============================================================
\newcommand{\pagegarde}[7]{%
  \thispagestyle{empty}
  \newgeometry{a4paper,margin=1.7cm,top=1.6cm,bottom=1.4cm}%
  \begin{tikzpicture}[remember picture,overlay]
    \fill[poporange!18] ([xshift=-3.2cm,yshift=-3.6cm]current page.north east) circle (4.6cm);
    \fill[poppurple!14] ([xshift=2.4cm,yshift=7.6cm]current page.south west) circle (3.8cm);
  \end{tikzpicture}%
  {\poplogo\fontsize{30}{30}\selectfont\color{popink}OnBuch\color{poporange}+}\hfill
  \raisebox{0.6ex}{\poppill{Cours signé \textbullet\ Premium}{popink}{popcream}}\par
  \vspace{1pt}\popsquiggle{poppurple}\par
  \vspace{8pt}
  \begin{tcolorbox}[enhanced,colback=poporange,colframe=popink,boxrule=2.8pt,arc=13pt,
      drop shadow=popink,left=18pt,right=18pt,top=12pt,bottom=13pt,after skip=10pt]
    \poppill{#2 \textbullet\ #3}{popink}{popcream}\hspace{5pt}\poppill{#4}{white}{popink}\par\vspace{8pt}
    {\popdisplay\fontsize{14}{14}\selectfont\color{popink}LEÇON #5}\par\vspace{4pt}
    {\raggedright\hyphenpenalty=10000\exhyphenpenalty=10000\popdisplay\fontsize{25}{27}\selectfont\color{popcream}\MakeUppercase{#1}\par}
    \vspace{8pt}
    {\popxbold\fontsize{9.5}{9.5}\selectfont\color{popink}\MakeUppercase{Durée conseillée : #6}}
  \end{tcolorbox}
  \begin{tcolorbox}[enhanced,colback=white,colframe=popink,boxrule=2.2pt,arc=10pt,
      drop shadow=popink,left=16pt,right=16pt,top=10pt,bottom=10pt,after skip=8pt]
    \poppill{Dans cette leçon}{poppurple}{white}\par\vspace{5pt}
    \popsemi\fontsize{9.3}{12.6}\selectfont\color{popink}#7
  \end{tcolorbox}
  \noindent\begin{minipage}[t]{0.487\linewidth}
    \begin{tcolorbox}[enhanced,colback=popgreen,colframe=popink,boxrule=2.2pt,arc=10pt,
        drop shadow=popink,left=11pt,right=11pt,top=10pt,bottom=10pt,height=3.1cm,valign=center]
      \tcbox[colback=white,colframe=popink,boxrule=1.3pt,arc=5pt,boxsep=0pt,nobeforeafter,
        left=3pt,right=3pt,top=3pt,bottom=3pt,valign=center]{\qrcode[height=1.9cm]{https://play.google.com/store/apps/details?id=com.luvvix.onbuch&hl=fr}}%
      \hspace{8pt}\begin{minipage}[c]{0.5\linewidth}
        {\popdisplay\fontsize{11}{12.5}\selectfont\color{white}\MakeUppercase{Télécharge OnBuch}}\par\vspace{4pt}
        {\raggedright\popsemi\fontsize{8.4}{11}\selectfont\color{white}Gratuit \textbullet\ Google Play\\Tuteur IA, annales, concours, résultats\par}
      \end{minipage}
    \end{tcolorbox}
  \end{minipage}\hfill
  \begin{minipage}[t]{0.487\linewidth}
    \begin{tcolorbox}[enhanced,colback=poppurple,colframe=popink,boxrule=2.2pt,arc=10pt,
        drop shadow=popink,left=11pt,right=11pt,top=10pt,bottom=10pt,height=3.1cm,valign=center]
      \tikz[baseline=-0.7ex]\node[circle,fill=white,draw=popink,line width=1.3pt,minimum size=1.6cm,inner sep=0]{\popdisplay\fontsize{24}{24}\selectfont\color{poppurple}+};%
      \hspace{8pt}\begin{minipage}[c]{0.6\linewidth}
        {\popdisplay\fontsize{11}{12.5}\selectfont\color{white}\MakeUppercase{Passe à OnBuch+}}\par\vspace{4pt}
        {\raggedright\popsemi\fontsize{8.4}{11}\selectfont\color{white}Tous les cours signés, Tuteur IA illimité, fiches PDF — onglet OnBuch+ de l'app\par}
      \end{minipage}
    \end{tcolorbox}
  \end{minipage}
  \vspace{8pt}
  \popcontenu
  \vfill
  \signatureonbuch
  \restoregeometry
  \newpage
}

% Rangée « Ce cours contient » (page de garde)
\newcommand{\poptile}[3]{% #1 couleur #2 gros texte #3 légende
  \begin{tcolorbox}[enhanced,colback=#1,colframe=popink,boxrule=2pt,arc=9pt,shadow={2.5pt}{-2.5pt}{0pt}{popink},
      left=6pt,right=6pt,top=9pt,bottom=9pt,halign=center,valign=center,height=1.75cm,width=\linewidth,nobeforeafter]
    {\popdisplay\fontsize{12.5}{13}\selectfont\color{white}\MakeUppercase{#2}}\par\vspace{3pt}
    {\centering\popsemi\fontsize{7.8}{9.5}\selectfont\color{white}#3\par}
  \end{tcolorbox}}
\newcommand{\popcontenu}{%
  \par\noindent{\popxboldls\fontsize{7.5}{7.5}\selectfont\color{popmuted}CE COURS SIGNÉ CONTIENT}\par\vspace{7pt}
  \noindent\begin{minipage}[t]{0.235\linewidth}\poptile{popblue}{Cours}{complet, illustré, pas à pas}\end{minipage}\hfill
  \begin{minipage}[t]{0.235\linewidth}\poptile{poporange}{Méthodes}{et exercices résolus}\end{minipage}\hfill
  \begin{minipage}[t]{0.235\linewidth}\poptile{poppink}{Exercices}{type examen + corrigés}\end{minipage}\hfill
  \begin{minipage}[t]{0.235\linewidth}\poptile{popgreen}{Fiche bilan}{l'essentiel à retenir}\end{minipage}\par}

% Sommaire
\newtcolorbox{sommaire}{enhanced,colback=white,colframe=popink,boxrule=2.2pt,arc=10pt,
  drop shadow=popink,left=18pt,right=18pt,top=13pt,bottom=13pt,before skip=6pt,after skip=20pt,
  before upper={\poppill{Sommaire}{poppurple}{white}\par\vspace{9pt}\popsemi\fontsize{10.6}{14.5}\selectfont\color{popink}}}

% Signature de fin de cours — poussée tout en bas de la dernière page
\newcommand{\findecours}{%
  \par\vfill\nopagebreak
  \begin{center}\popsquiggle{poporange}\end{center}\vspace{4pt}
  \signatureonbuch
}
\AtBeginDocument{\def\llbracket{\mathopen{[\mkern-3.5mu[}}\def\rrbracket{\mathclose{]\mkern-3.5mu]}}} % intervalles d'entiers (glyphe ⟦ absent de la police)
% Glyphes absents de Fira Math : redéfinis à partir de glyphes disponibles
\AtBeginDocument{%
  \def\setminus{\mathbin{\backslash}}\let\smallsetminus\setminus
  \def\vdots{\mathord{\vbox{\baselineskip3.2pt\lineskiplimit0pt\kern4pt\hbox{.}\hbox{.}\hbox{.}}}}%
  \def\ddots{\mathinner{\mkern1mu\raise6.5pt\hbox{.}\mkern2mu\raise3.6pt\hbox{.}\mkern2mu\raise0.7pt\hbox{.}\mkern1mu}}%
  \def\triangle{\mathord{\scalebox{0.9}{$\symup{\Delta}$}}}%
  \def\nabla{\mathord{\raisebox{\depth}{\scalebox{1}[-1]{$\symup{\Delta}$}}}}%
  \def\checkmark{\popcheck{popdark}}%
  \def\star{\mathbin{\ast}}\def\bigstar{\mathord{\ast}}%
  \def\diamond{\mathbin{\scalebox{0.6}{$\Diamond$}}}%
  \def\bigcup{\mathop{\mathchoice{\raisebox{-0.3em}{\scalebox{1.7}{$\cup$}}}{\scalebox{1.25}{$\cup$}}{\cup}{\cup}}\displaylimits}%
  \def\bigcap{\mathop{\mathchoice{\raisebox{-0.3em}{\scalebox{1.7}{$\cap$}}}{\scalebox{1.25}{$\cap$}}{\cap}{\cap}}\displaylimits}%
  \def\lesseqqgtr{\mathrel{\lessgtr}}\def\bigoplus{\mathop{\oplus}\displaylimits}%
}
\providecommand{\ssi}{si et seulement si} % abréviation fréquente
\AtBeginDocument{\providecommand{\wideparen}{\overparen}} % arc de cercle

\begin{document}

\pagegarde{La reconstruction/renaissance de l’Europe et les révolutions industrielles}{Histoire}{1re Tech}{Histoire – classe de Première technique}{N01}{2 séances (fiche de progression harmonisée)}{Cette leçon analyse le double processus de reconstruction politique de l'Europe post-napoléonienne (Congrès de Vienne, Printemps des peuples) et de sa mutation économique par les révolutions industrielles, fondements de l'impérialisme qui touchera le Cameroun.
\vspace{4pt}
\begin{multicols}{2}\small
\begin{itemize}
\item À la fin de cette leçon, je sais expliquer les décisions du Congrès de Vienne (1815) et leur impact sur la carte de l'Europe.
\item Je sais analyser les causes et les étapes des deux révolutions industrielles (charbon/fer/vapeur puis électricité/pétrole/acier).
\item Je sais établir le lien entre les besoins industriels (matières premières, débouchés, capitaux) et le partage de l'Afrique (Conférence de Berlin).
\item Je sais construire un croquis synthétique de l'Europe industrielle vers 1850 et 1900 à partir de données statistiques.
\item Je sais rédiger une argumentation structurée justifiant l'entrée de l'Europe dans l'ère impériale.
\item Je sais mobiliser la notion de 'révolution des transports' pour expliquer l'intégration des marchés camerounais à l'économie mondiale.
\end{itemize}\end{multicols}}

\begin{sommaire}
\begin{multicols}{2}
\begin{enumerate}
\item I. L'Europe sort des guerres : reconstruction politique et ordre nouveau (1815-1848)
\item II. La Première Révolution Industrielle : le charbon, le fer et la vapeur (vers 1780-1850)
\item III. La Deuxième Révolution Industrielle : acier, électricité, pétrole, chimie (1850-1914)
\item IV. Les fondements de l'impérialisme européen : pourquoi l'Afrique ? (1870-1914)
\item V. Synthèse : L'Europe, atelier du monde et gendarme colonial (Bilan 1914)
\item Activité d'intégration
\item Méthodes et exercices résolus
\item Exercices
\item Corrigés des exercices
\item Fiche bilan
\end{enumerate}
\end{multicols}
\end{sommaire}

\begin{objectifs}
\begin{itemize}
\item À la fin de cette leçon, je sais expliquer les décisions du Congrès de Vienne (1815) et leur impact sur la carte de l'Europe.
\item Je sais analyser les causes et les étapes des deux révolutions industrielles (charbon/fer/vapeur puis électricité/pétrole/acier).
\item Je sais établir le lien entre les besoins industriels (matières premières, débouchés, capitaux) et le partage de l'Afrique (Conférence de Berlin).
\item Je sais construire un croquis synthétique de l'Europe industrielle vers 1850 et 1900 à partir de données statistiques.
\item Je sais rédiger une argumentation structurée justifiant l'entrée de l'Europe dans l'ère impériale.
\item Je sais mobiliser la notion de 'révolution des transports' pour expliquer l'intégration des marchés camerounais à l'économie mondiale.
\item Je sais comparer les modèles d'industrialisation britannique, belge, française et allemande à l'aide de graphiques de production.
\end{itemize}
\end{objectifs}

\begin{prerequis}
\begin{itemize}
\item La Révolution française et l'Empire napoléonien (classe de 3ème / début 1ère).
\item La traite négrière et le commerce légitime en Afrique précoloniale (Histoire 4ème / 3ème).
\item Notions de base en économie : capital, travail, marché, matière première (Enseignement technologique).
\item Lecture de cartes thématiques et de graphiques statistiques (compétences transversales).
\item Le vocabulaire spécifique : industrialisation, impérialisme, protectionnisme, libre-échange.
\end{itemize}
\end{prerequis}

\newpage

\coursec{I. L'Europe sort des guerres : reconstruction politique et ordre nouveau (1815-1848)}

\textbf{Situation d'accroche.} En observant les usines de la zone industrielle de Bassa à Douala ou les chantiers du barrage de Nachtigal sur la Sanaga, on voit des machines, de l'acier et de l'énergie électrique à l'œuvre. Savez-vous que ces technologies, qui transforment aujourd'hui le Cameroun, sont nées en Europe il y a deux siècles lors d'un bouleversement sans précédent ? Mais avant de devenir « l'atelier du monde », l'Europe a d'abord dû panser ses plaies : entre 1792 et 1815, les guerres de la Révolution et de l'Empire napoléonien ont ravagé le continent, déplacé les frontières et renversé des trônes.

\textbf{Problématique.} Comment l'Europe, ravagée par les guerres napoléoniennes, a-t-elle réussi à renaître, à organiser la paix, puis à affronter la vague révolutionnaire de 1848, avant de devenir l'atelier du monde et de conquérir l'Afrique ?

\courssub{Le Congrès de Vienne (septembre 1814 -- juin 1815) : refaire la carte de l'Europe}

Après la première abdication de Napoléon Bonaparte, exilé à l'île d'Elbe en avril 1814, l'Europe est en ruines. Les frontières tracées par les conquêtes françaises ne valent plus rien, les dynasties détrônées réclament leurs couronnes, et les vainqueurs veulent des compensations. Pour tout régler, les grandes puissances se réunissent à Vienne, capitale de l'Autriche. Le Congrès est interrompu par le retour éclair de Napoléon (les « Cent-Jours », mars-juin 1815), définitivement battu à Waterloo le 18 juin 1815 et exilé à Sainte-Hélène.

\begin{definition}[Le Congrès de Vienne]
Le \cle{Congrès de Vienne} est une réunion diplomatique majeure (septembre 1814 -- juin 1815) qui redessine la carte de l'Europe après la chute de Napoléon, afin de garantir une paix durable par l'\cle{équilibre des puissances}. Son document final est l'\cle{Acte final} du 9 juin 1815.
\end{definition}

\textbf{Les acteurs.} Quatre grandes puissances victorieuses dominent les débats :
\begin{itemize}
\item l'\textbf{Autriche}, représentée par son habile ministre des Affaires étrangères, le prince de \cle{Metternich}, véritable « maître d'orchestre » du Congrès ;
\item la \textbf{Russie} du tsar \cle{Alexandre Ier}, qui veut étendre son influence vers l'ouest (Pologne) ;
\item la \textbf{Prusse}, qui convoite la Rhénanie et la Saxe ;
\item la \textbf{Grande-Bretagne}, représentée par Castlereagh, qui veut empêcher toute domination du continent par une seule puissance et garder ses conquêtes coloniales.
\end{itemize}
Fait remarquable : la \textbf{France} vaincue n'est pas exclue. Son diplomate \cle{Talleyrand} réussit, par son talent, à faire admettre la France dans le cercle des négociateurs et à limiter les pertes françaises.

\textbf{Les trois principes directeurs.} L'intuition de Metternich est simple : pour éviter qu'une nouvelle guerre générale n'éclate, il faut que personne ne soit ni trop fort ni trop humilié. De cette intuition découlent trois principes rigoureux :
\begin{enumerate}
\item le principe de \cle{légitimité} : les dynasties renversées par Napoléon sont restaurées. En France, les Bourbon reviennent avec Louis XVIII ; en Espagne, les Bourbon également ;
\item le principe d'\cle{équilibre} : aucune puissance ne doit dominer le continent. On renforce donc les États voisins de la France pour l'encadrer : la Prusse reçoit la Rhénanie, les Pays-Bas reçoivent la Belgique ;
\item le principe de \cle{compensation} : les vainqueurs reçoivent des territoires en récompense de leurs efforts de guerre (la Russie obtient l'essentiel de la Pologne, la Grande-Bretagne garde Malte, le Cap et Ceylan).
\end{enumerate}

\begin{exemplebox}[Un exemple chiffré : le prix de la défaite pour la France]
Par le second traité de Paris (20 novembre 1815), la France paie \textbf{700 millions de francs-or} d'indemnité de guerre, voit son territoire ramené aux frontières de 1790 (elle perd notamment la Sarre et la Savoie) et supporte une occupation militaire alliée de plusieurs années. Malgré cela, grâce à Talleyrand, elle reste une grande puissance reconnue.
\end{exemplebox}

\begin{methode}[Analyser un traité diplomatique : l'Acte final de Vienne]
Pour analyser un document diplomatique comme l'Acte final de Vienne, procédez en quatre étapes :
\begin{enumerate}
\item \textbf{Identifier les signataires} : qui négocie ? Qui est exclu ? (Ici : les grandes puissances décident, les petits États et les peuples ne sont pas consultés.)
\item \textbf{Repérer les clauses territoriales} : quels territoires changent de mains ? (Rhénanie à la Prusse, Belgique aux Pays-Bas, Pologne à la Russie.)
\item \textbf{Dégager les principes directeurs} : quelle logique organise le texte ? (Légitimité, équilibre, compensation.)
\item \textbf{Chercher les faiblesses} : que voit-on à l'usage ? Le traité ignore totalement les \cle{nationalités} : les Polonais sont partagés, les Belges sont rattachés aux Pays-Bas contre leur gré, les Allemands et les Italiens restent morcelés. Cette faiblesse explique les révolutions de 1830 et de 1848.
\end{enumerate}
\end{methode}

\courssub{La Sainte-Alliance, la Quadruple-Alliance et le Concert européen}

\textbf{Les alliances du maintien de l'ordre.} Pour défendre l'ordre établi à Vienne, deux instruments sont créés :
\begin{itemize}
\item la \cle{Sainte-Alliance} (septembre 1815), voulue par le tsar Alexandre Ier : c'est un pacte mystique et religieux unissant les monarques d'Autriche, de Russie et de Prusse, qui promettent de gouverner selon les principes chrétiens et de s'entraider contre les révolutions ;
\item la \cle{Quadruple-Alliance} (novembre 1815) : Autriche, Russie, Prusse et Grande-Bretagne s'engagent concrètement à maintenir l'ordre monarchique et à intervenir militairement contre les libéraux et les nationalistes. C'est elle qui écrase les soulèvements libéraux à Naples (1821) et en Espagne (1823).
\end{itemize}

\begin{attention}[Erreur fréquente à éviter]
Ne confondez pas la \textbf{Sainte-Alliance} et le \textbf{Concert européen}. La Sainte-Alliance est un \emph{pacte mystique et monarchique} (1815), une déclaration de principes religieux entre souverains. Le \cle{Concert européen} est une \emph{pratique diplomatique de concertation} : les grandes puissances se réunissent régulièrement en conférences pour gérer les crises ensemble, sans idéologie religieuse. L'un est un serment, l'autre est une méthode de travail.
\end{attention}

\textbf{Le Concert européen en action.} Plutôt que de se faire la guerre, les grandes puissances se concertent. Cette diplomatie des conférences gère plusieurs crises :
\begin{itemize}
\item la \textbf{question grecque} (1821-1830) : les Grecs se soulèvent contre l'Empire ottoman ; la Grande-Bretagne, la France et la Russie interviennent (bataille navale de Navarin, 1827) et imposent l'indépendance de la Grèce en 1830 ;
\item la \textbf{question belge} (1830-1831) : les Belges se révoltent contre les Pays-Bas ; les puissances reconnaissent l'indépendance de la Belgique et garantissent sa neutralité ;
\item la \textbf{question d'Orient} : les puissances surveillent de près l'affaiblissement de l'Empire ottoman, « l'homme malade de l'Europe », chacune cherchant à profiter de son déclin sans laisser les autres s'agrandir.
\end{itemize}

\begin{aretenir}[L'essentiel sur l'ordre de Vienne]
Entre 1815 et 1848, l'Europe vit sous le « système de Metternich » : restauration des monarchies légitimes, équilibre des puissances, répression des idées libérales et nationales, et gestion collective des crises par le Concert européen. Ce système assure près de quarante ans sans guerre générale, mais il étouffe deux aspirations profondes des peuples : la \textbf{liberté} (constitutions, suffrage) et la \textbf{nation} (unité, indépendance).
\end{aretenir}

\begin{popfigure}
\begin{center}
\begin{tikzpicture}[scale=0.92, every node/.style={font=\small}]
% Fond Europe schématique
\fill[popcream] (-0.5,-0.5) rectangle (12.5,7.5);
\node[font=\bfseries\large, popdark] at (6,7.1) {L'Europe de 1815 à 1848 : ordre de Vienne et foyers révolutionnaires};
% Blocs simplifiés
\fill[popblueL, draw=popblue, thick] (0.5,3.2) rectangle (2.6,5.4); % France
\fill[popgreenL, draw=popgreen, thick] (2.7,4.4) rectangle (4.4,5.6); % Pays-Bas/Belgique
\fill[popgoldL, draw=popgold, thick] (2.8,2.2) rectangle (4.6,4.3); % Confédération germanique
\fill[popgoldL, draw=popgold, thick] (4.7,3.4) rectangle (6.6,5.2); % Prusse
\fill[poporangeL, draw=poporange, thick] (4.7,1.6) rectangle (6.6,3.3); % Autriche
\fill[poppinkL, draw=poppink, thick] (2.9,0.2) rectangle (4.3,2.0); % Italie (botte simplifiée)
\fill[popgreenL, draw=popgreen, thick] (6.7,2.0) rectangle (10.5,5.8); % Russie
\fill[popmuted!30, draw=popmuted, thick] (6.8,0.2) rectangle (9.3,1.8); % Empire ottoman
% Noms
\node at (1.55,4.3) {\textbf{France}};
\node[font=\scriptsize] at (3.55,5.0) {Belgique};
\node[font=\scriptsize, align=center] at (3.7,3.2) {Confédération\\germanique (39 États)};
\node at (5.65,4.3) {\textbf{Prusse}};
\node at (5.65,2.45) {\textbf{Autriche}};
\node[font=\scriptsize, align=center] at (3.6,1.1) {Italie\\morcelée};
\node at (8.6,3.9) {\textbf{Russie}};
\node[font=\scriptsize] at (8.05,1.0) {Empire ottoman};
% Foyers révolutionnaires 1848 (étoiles)
\node[popink, font=\large] at (1.0,5.0) {$\bigstar$};
\node[popink, font=\large] at (3.2,3.9) {$\bigstar$};
\node[popink, font=\large] at (5.2,1.9) {$\bigstar$};
\node[popink, font=\large] at (3.4,1.7) {$\bigstar$};
\node[popink, font=\large] at (6.2,2.9) {$\bigstar$};
% Flèches des mouvements nationaux
\draw[-{Stealth[length=3mm]}, very thick, poppurple] (4.6,3.0) -- (5.6,3.9);
\draw[-{Stealth[length=3mm]}, very thick, poppurple] (3.6,0.5) -- (3.6,1.6);
\draw[-{Stealth[length=3mm]}, very thick, poppurple] (6.6,2.2) -- (6.0,2.5);
% Légende
\node[anchor=west, font=\scriptsize] at (10.7,5.4) {\textcolor{popink}{$\bigstar$} Foyers du};
\node[anchor=west, font=\scriptsize] at (10.7,5.0) {« Printemps des peuples » (1848)};
\draw[-{Stealth[length=2mm]}, thick, poppurple] (10.7,4.4) -- (11.3,4.4);
\node[anchor=west, font=\scriptsize] at (11.4,4.4) {Mouvements nationaux};
% Nord
\draw[-{Stealth[length=2.5mm]}, thick, popdark] (11.8,6.0) -- (11.8,6.8);
\node[font=\scriptsize\bfseries] at (11.8,5.8) {N};
\end{tikzpicture}
\end{center}
\legende{Croquis schématique de l'Europe après le Congrès de Vienne (1815). Les grandes monarchies (Russie, Autriche, Prusse) encadrent la France ; l'Allemagne et l'Italie restent morcelées. En 1848, les foyers révolutionnaires (étoiles) éclatent de Paris à Vienne, et les flèches violettes indiquent les mouvements nationaux vers l'unité allemande et italienne. Croquis simplifié, sans précision cartographique.}
\end{popfigure}

\courssub{L'ébranlement de 1848 : le « Printemps des peuples »}

En février 1848, une révolution éclate à Paris et chasse le roi Louis-Philippe. Comme une traînée de poudre, la révolte gagne en quelques semaines Vienne, Berlin, Milan, Venise, Prague et Budapest : c'est le \cle{Printemps des peuples}. Les causes sont à la fois économiques (crise agricole et chômage en 1846-1847) et politiques.

Les révolutionnaires portent deux types de revendications :
\begin{itemize}
\item des revendications \textbf{libérales} : des constitutions limitant le pouvoir des rois, la liberté de la presse, l'élargissement du suffrage (le droit de vote) ;
\item des revendications \textbf{nationales} : l'unité de l'Allemagne (les députés se réunissent au Parlement de Francfort pour rédiger une constitution allemande), l'unité de l'Italie contre l'occupant autrichien, l'indépendance de la Hongrie contre Vienne.
\end{itemize}

\textbf{Bilan mitigé.} À court terme, c'est un échec : les armées des monarques reprennent le contrôle partout en 1849, le Parlement de Francfort se dissout sans avoir réalisé l'unité allemande. Mais à moyen terme, 1848 porte des fruits décisifs :
\begin{itemize}
\item en \textbf{France}, la II\textsuperscript{e} République proclame le \cle{suffrage universel masculin} (1848) : tous les hommes adultes votent, une conquête démocratique majeure ;
\item l'échec des révolutions convainc les patriotes allemands et italiens que l'unité ne se fera pas par les discours, mais par la force : ce sera l'unification « par le \cle{fer et le sang} » menée par \cle{Bismarck} en Prusse (unité allemande achevée en 1871) et par \cle{Cavour} au Piémont (unité italienne en 1861-1870).
\end{itemize}

\begin{popfigure}
\begin{center}
\begin{tikzpicture}
\begin{axis}[
  ybar, bar width=22pt,
  width=0.8\linewidth, height=6cm,
  ymin=0, ymax=45,
  ylabel={Nombre d'États allemands},
  symbolic x coords={1815 (Confédération germanique), 1871 (Empire allemand)},
  xtick=data,
  x tick label style={font=\small, align=center},
  nodes near coords,
  every node near coord/.append style={font=\bfseries\small, popdark},
  axis lines=left,
  ymajorgrids=true, grid style={popline!40},
]
\addplot[fill=popblue, draw=popdark] coordinates {(1815 (Confédération germanique),39) (1871 (Empire allemand),1)};
\end{axis}
\end{tikzpicture}
\end{center}
\legende{Du morcellement à l'unité : en 1815, le Congrès de Vienne maintient 39 États allemands au sein de la Confédération germanique ; en 1871, l'Empire allemand proclamé à Versailles les réunit en un seul État. Source : données historiques du traité de Vienne et de la proclamation de l'Empire.}
\end{popfigure}

\courssub{Conséquence pour l'Afrique : l'Europe stabilisée tourne son regard vers l'extérieur}

L'ordre de Vienne, puis la stabilisation progressive du continent après 1848, ont une conséquence capitale pour l'Afrique : une Europe apaisée à l'intérieur dispose de forces, de capitaux et de curiosité pour l'extérieur. Dès la première moitié du XIX\textsuperscript{e} siècle, les explorations se multiplient (Mungo Park sur le Niger, Livingstone en Afrique australe), et les comptoirs européens s'installent durablement sur les côtes africaines, notamment dans le golfe de Guinée. C'est le prélude à la colonisation que nous étudierons dans la suite de la leçon.

\begin{savaistu}[De Vienne (1815) à Berlin (1885) : un lien direct avec le Cameroun]
Le Cameroun n'est pas concerné directement par le Congrès de Vienne. Mais un principe posé en 1815 sera décisif pour son histoire : la \textbf{liberté de navigation sur les fleuves internationaux} (Rhin, Danube). Ce principe sera réactivé à la Conférence de Berlin (1884-1885) pour le Congo et le Niger, ouvrant juridiquement la voie à la pénétration européenne en Afrique centrale — et donc au Cameroun, devenu protectorat allemand le 12 juillet 1884.
\end{savaistu}

\begin{exoresolu}[Exercice d'application : distinguer les principes de Vienne]
\textbf{Énoncé.} Voici trois décisions prises en 1815 : (a) Louis XVIII remonte sur le trône de France ; (b) la Prusse reçoit la Rhénanie pour surveiller la France ; (c) la Russie reçoit l'essentiel de la Pologne en récompense de sa participation à la guerre. Attribuez à chaque décision le principe du Congrès de Vienne qu'elle illustre, en justifiant.
\tcblower
\textbf{Solution détaillée.}
\begin{enumerate}
\item (a) illustre le principe de \textbf{légitimité} : on restaure la dynastie des Bourbon, considérée comme la famille régnante légitime, chassée par la Révolution et Napoléon.
\item (b) illustre le principe d'\textbf{équilibre} : en donnant la Rhénanie (frontalière de la France) à la Prusse, on crée un État fort capable de contenir une éventuelle nouvelle expansion française, afin qu'aucune puissance ne domine le continent.
\item (c) illustre le principe de \textbf{compensation} : la Russie, victorieuse, reçoit un territoire en récompense de ses efforts de guerre.
\end{enumerate}
\textbf{Conclusion.} Un même congrès peut appliquer plusieurs principes à la fois ; l'examen demande de les identifier précisément et de les justifier par les faits.
\end{exoresolu}

\begin{exercice}[Pour s'entraîner]
Rédigez un paragraphe argumenté (10 à 15 lignes) répondant à la question : « Le Congrès de Vienne a-t-il réussi à garantir la paix en Europe ? » Vous mobiliserez au moins trois arguments (équilibre des puissances, Concert européen, ignorance des nationalités) et vous conclurez par un bilan nuancé.
\end{exercice}

\begin{corrige}[Exercice]
\textbf{Éléments de corrigé.} Le candidat doit d'abord reconnaître les réussites : l'équilibre des puissances empêche toute guerre générale pendant près de quarante ans ; le Concert européen règle pacifiquement les crises grecque et belge ; la France vaincue n'est pas humiliée, ce qui évite un esprit de revanche immédiat. Ensuite, il doit montrer les limites : le système repose sur la répression (Sainte-Alliance, interventions à Naples et en Espagne) et ignore les aspirations nationales et libérales, ce qui provoque les révolutions de 1830 et de 1848. \textbf{Conclusion nuancée} : Vienne garantit la paix entre États, mais pas la paix intérieure des peuples ; l'ordre de 1815 s'effondre sous le « Printemps des peuples » de 1848.
\end{corrige}

\coursec{II. La Première Révolution Industrielle : le charbon, le fer et la vapeur (vers 1780-1850)}

Dans la section précédente, nous avons vu comment l'Europe, sortie des guerres napoléoniennes, a reconstruit son ordre politique au Congrès de Vienne (1815). Mais pendant que les diplomates redessinaient les frontières, une transformation bien plus profonde était déjà en marche, silencieuse et irréversible : dans les vallées anglaises, des usines crachaient leur fumée, des machines filaient le coton jour et nuit, et bientôt des locomotives siffleraient sur des rails de fer. C'est cette mutation économique et sociale, la \cle{Première Révolution industrielle}, que nous étudions maintenant. Elle fait de l'Europe, et d'abord de la Grande-Bretagne, l'\emph{atelier du monde}.

\courssub{Qu'est-ce que la Révolution industrielle ?}

Commençons par l'intuition. Imagine un village anglais vers 1750 : on y file la laine à la main, au rouet, dans les fermes ; on y forge le fer au marteau, avec du charbon de bois ; on se déplace à pied ou à cheval. Un siècle plus tard, dans ce même pays, des usines de plusieurs étages emploient des milliers d'ouvriers, des machines actionnées par la vapeur produisent en une heure ce qu'un artisan faisait en une semaine, et le train relie les villes à près de 50 km/h. Ce bouleversement, c'est la Révolution industrielle.

\begin{definition}[Révolution industrielle]
La \cle{Révolution industrielle} est le passage d'une économie \textbf{agraire et artisanale} (production manuelle, dispersée, dépendante des forces naturelles : muscles, vent, eau, bois) à une économie \textbf{industrielle et mécanisée}, fondée sur trois piliers : l'\textbf{énergie fossile} (le charbon), la \textbf{machine} (mue par la vapeur) et l'\textbf{usine} (concentration des ouvriers et du capital). Elle commence en Grande-Bretagne vers 1780 et se diffuse en Europe continentale au XIX\textsuperscript{e} siècle.
\end{definition}

\begin{attention}[Le mot « révolution » peut tromper]
Contrairement à une révolution politique (quelques jours ou quelques années), la Révolution industrielle est un processus \textbf{long} : elle s'étale sur plusieurs générations. On parle de « révolution » parce que ses \emph{conséquences} sont radicales : elles transforment le travail, la ville, la famille, le paysage et les rapports sociaux.
\end{attention}

\courssub{Les préconditions : pourquoi la Grande-Bretagne en premier ?}

Pourquoi ce bouleversement naît-il en Angleterre, et non en France ou en Chine, pourtant riches et peuplées ? Trois conditions se sont combinées sur l'île britannique au XVIII\textsuperscript{e} siècle.

\textbf{a) La révolution agraire.} Dès le XVII\textsuperscript{e} siècle, les \emph{enclosures} regroupent les terres en grandes exploitations clôturées. On pratique l'assolement amélioré (rotation avec fourrages, suppression de la jachère), on sélectionne le bétail, on utilise de nouveaux outils. Résultat : les \textbf{rendements augmentent fortement}. Deux conséquences décisives :
\begin{itemize}
\item des \textbf{excédents alimentaires} permettent de nourrir une population urbaine croissante sans importer massivement de blé ;
\item les petits paysans expropriés deviennent une \textbf{main-d'œuvre disponible} qui migre vers les villes et les usines.
\end{itemize}

\textbf{b) L'accumulation capitaliste.} Le commerce triangulaire, la traite négrière et les plantations coloniales (sucre, coton) ont rapporté d'immenses profits aux ports européens : Liverpool et Bristol en Angleterre, Nantes et Bordeaux en France. Ces capitaux, accumulés pendant deux siècles, cherchent à \emph{fructifier} : ils seront investis dans les filatures, les mines et les chemins de fer. Les banques anglaises, développées depuis la création de la Banque d'Angleterre (1694), facilitent le crédit.

\textbf{c) Les progrès techniques.} L'invention décisive est la \textbf{machine à vapeur}. Thomas Newcomen (1712) en avait conçu une première version, très gourmande en charbon, utilisée pour épuiser l'eau des mines. En \textbf{1769}, l'Écossais \cle{James Watt} dépose le brevet d'une machine à \emph{condenseur séparé}, beaucoup plus économique, et surtout capable de transformer la vapeur en \textbf{mouvement rotatif} : la vapeur peut désormais actionner n'importe quelle machine, n'importe où, sans dépendre d'une rivière. L'industrie est libérée des contraintes de la nature.

\begin{propriete}[Les conditions du « décollage » (selon W. Rostow)]
L'économiste Walt Rostow a théorisé le \emph{take-off} (décollage) industriel. Il nécessite trois conditions réunies :
\begin{enumerate}
\item un \textbf{taux d'investissement supérieur à 10 \%} du produit national (contre 4 à 5 \% dans une économie traditionnelle) ;
\item un \textbf{secteur manufacturier leader} qui entraîne les autres (chez les Anglais : le textile coton) ;
\item un \textbf{cadre institutionnel favorable} : propriété privée garantie, liberté du commerce, marché unifié.
\end{enumerate}
La Grande-Bretagne est le premier pays à réunir ces trois conditions, vers 1780-1830.
\end{propriete}

\courssub{Le modèle britannique : « l'atelier du monde »}

\textbf{a) Le textile, secteur moteur.} Le coton importé des Indes et d'Amérique est filé puis tissé par des machines : la \emph{spinning jenny} de Hargreaves (1765), le \emph{water frame} d'Arkwright (1769), la \emph{mule-jenny} de Crompton (1779), puis le métier à tisser mécanique de Cartwright (1785). Le prix du tissu s'effondre tandis que le volume de production explose : le coton britannique inonde le monde entier.

\textbf{b) La sidérurgie au coke.} Le bois devenant rare, Abraham Darby utilise dès 1709 le \textbf{coke} (charbon purifié) pour fondre le minerai de fer. Henry Cort perfectionne le procédé du \emph{puddlage} (1784), qui donne un fer de meilleure qualité. Le fer devient abondant et bon marché : on peut construire machines, rails, ponts et navires.

\textbf{c) Le charbon, sang de la révolution.} Tout repose sur le charbon : il chauffe les hauts fourneaux, alimente les machines à vapeur, éclaire les villes (gaz d'éclairage). Les bassins de Newcastle, du Pays de Galles et des Midlands deviennent des paysages noirs de puits et de terrils.

\textbf{d) Le rôle de l'État.} L'État britannique accompagne le mouvement par le \textbf{libre-échange} (abrogation des \emph{Corn Laws} en 1846, qui taxaient le blé étranger) et par les \emph{Poor Laws} (lois sur les pauvres, réformées en 1834), qui poussent les indigents vers le travail salarié en usine. La Royal Navy, première flotte du monde, protège les routes commerciales.

\begin{exemplebox}[La suprématie britannique en chiffres (vers 1850)]
En 1850, le Royaume-Uni, avec environ 2 \% de la population mondiale, produit environ :
\begin{itemize}
\item \textbf{les deux tiers du charbon mondial} ;
\item \textbf{la moitié du fer mondial} ;
\item \textbf{la moitié du coton filé mondial}.
\end{itemize}
L'Exposition universelle de Londres (1851), sous le \emph{Crystal Palace} de fer et de verre, est une démonstration triomphale de cette puissance. La Belgique, elle, est déjà le premier producteur continental de fer par habitant.
\end{exemplebox}

\begin{popfigure}
\begin{center}
\begin{tikzpicture}[scale=0.95, every node/.style={font=\small}]
% === Usine (gauche) ===
\begin{scope}
\node[font=\bfseries, popblue] at (2.6,4.6) {Le système d'usine};
% chaudière
\draw[fill=poporangeL] (0.2,0.4) rectangle (1.6,1.6);
\node[align=center] at (0.9,1.0) {Chaudière\\(charbon)};
\draw[fill=popmuted!40] (0.5,1.6) rectangle (0.7,2.3);
\draw[fill=popmuted!40] (1.1,1.6) rectangle (1.3,2.3);
\node[popmuted] at (0.9,2.7) {\tiny vapeur};
% machine à vapeur
\draw[fill=popblueL] (2.4,0.4) circle (0.65);
\node[align=center] at (2.4,0.4) {\tiny Machine\\[-2pt]\tiny à vapeur};
\draw[-{Stealth}, thick, poporange] (1.6,1.0) -- (1.75,0.6);
% arbre de transmission
\draw[very thick, popdark] (3.05,0.4) -- (5.6,0.4);
\node[below, popdark] at (4.3,0.25) {\tiny arbre de transmission};
% poulies et courroies vers machines
\foreach \x in {3.5,4.4,5.3}{
  \draw[fill=popgreenL] (\x-0.35,2.2) rectangle (\x+0.35,3.2);
  \draw[popdark] (\x,0.4) -- (\x,2.2);
  \draw[popdark] (\x,0.4) circle (0.12);
}
\node at (3.5,2.7) {\tiny métier};
\node at (4.4,2.7) {\tiny métier};
\node at (5.3,2.7) {\tiny métier};
\node[align=center, popblue] at (4.4,3.7) {\tiny dizaines de machines\\[-2pt]\tiny sous un même toit};
% ouvriers
\node[align=center] at (2.6,-0.5) {\tiny ouvriers salariés, horaires fixes, division du travail};
\end{scope}
% === Artisanat (droite) ===
\begin{scope}[xshift=7.2cm]
\node[font=\bfseries, popgreen] at (2.2,4.6) {L'artisanat (avant)};
\draw[fill=popcream] (0.4,0.4) rectangle (2.6,1.8);
\node[align=center] at (1.5,1.1) {Atelier\\familial};
\draw[fill=popgreenL] (3.2,0.9) circle (0.35);
\node[below] at (3.2,0.45) {\tiny artisan};
\draw[-{Stealth}, popgreen] (2.6,1.1) -- (2.85,1.1);
\node[align=center] at (1.8,2.6) {\tiny force humaine, eau, vent\\[-2pt]\tiny production lente, sur commande};
\draw[fill=popgoldL] (3.0,2.1) rectangle (3.5,2.5);
\node[right] at (3.55,2.3) {\tiny produit fini};
\draw[-{Stealth}, popgold] (3.2,1.3) -- (3.2,2.05);
\end{scope}
% flèche de transition
\draw[-{Stealth}, very thick, poppink] (6.0,2.2) -- (7.0,2.2);
\node[above, poppink] at (6.5,2.25) {\tiny 1780-1850};
\end{tikzpicture}
\end{center}
\legende{Schéma comparatif : dans l'usine, une seule source d'énergie (la vapeur, 1) actionne via un arbre de transmission (2) des dizaines de machines (3) concentrées sous un même toit, avec des ouvriers salariés ; dans l'artisanat, l'artisan (4) produit seul, lentement, avec la force musculaire ou hydraulique.}
\end{popfigure}

\courssub{La diffusion en Europe continentale}

La révolution industrielle franchit la Manche, mais à des rythmes très différents.

\textbf{a) La Belgique, premier pays continental industrialisé.} Dès les années 1820-1830, le bassin de la \textbf{Sambre-et-Meuse} (Charleroi, Liège) exploite son charbon et développe la sidérurgie, grâce à des techniciens et des capitaux anglais (John Cockerill et sa famille à Seraing). Indépendante en 1830, la Belgique devient le pays le plus industrialisé du continent : densité de voies ferrées record, exportation de rails et de machines dans toute l'Europe.

\textbf{b) La France, une industrialisation « douce ».} La France s'industrialise plus \emph{lentement} : paysannerie nombreuse et attachée à sa terre, charbon rare et mal réparti (Nord, Loire), capitalisme prudent qui préfère le placement à l'investissement risqué. Il en résulte un \textbf{dualisme} persistant : d'un côté une industrie lourde (Le Creusot, mines du Nord), de l'autre un artisanat de qualité et de luxe (soieries de Lyon, modes parisiennes) qui survit longtemps. La France ne connaît pas de rupture brutale mais une croissance régulière : c'est un \emph{déclin relatif}, non un retard absolu.

\textbf{c) Les États allemands : le Zollverein puis la Ruhr.} L'Allemagne n'existe pas encore : c'est une mosaïque d'une trentaine d'États. En \textbf{1834}, la Prusse fonde le \cle{Zollverein}, union douanière qui supprime les barrières intérieures et crée un grand marché unifié de plus de 25 millions de consommateurs. L'industrialisation démarre alors dans la \textbf{Ruhr} et la \textbf{Sarre} (charbon), en Silésie et en Saxe. Après 1850, le rythme allemand devient fulgurant : l'Allemagne rattrapera puis dépassera la Grande-Bretagne lors de la deuxième révolution industrielle.

\begin{popfigure}
\begin{center}
\begin{tikzpicture}
\begin{axis}[
  width=0.85\linewidth, height=6.5cm,
  xlabel={Année}, ylabel={Production de charbon (Mt)},
  xmin=1800, xmax=1850, ymin=0, ymax=60,
  xtick={1800,1810,1820,1830,1840,1850},
  ytick={0,10,20,30,40,50,60},
  legend pos=north west, legend style={font=\scriptsize},
  grid=both, grid style={popline!40},
  axis line style={popdark},
]
\addplot[popcurve, color=popblue, very thick] coordinates {(1800,15) (1810,17) (1820,21) (1830,29) (1840,42) (1850,57)};
\addlegendentry{Grande-Bretagne}
\addplot[popcurve, color=poporange, thick] coordinates {(1800,2.5) (1810,3.5) (1820,4) (1830,5) (1840,6.5) (1850,8.3)};
\addlegendentry{États allemands}
\addplot[popcurve, color=popgreen, thick] coordinates {(1800,0.8) (1810,1) (1820,1.4) (1830,2.5) (1840,4) (1850,5.8)};
\addlegendentry{France}
\addplot[popcurve, color=poppink, thick] coordinates {(1800,0.5) (1810,1) (1820,2) (1830,3.5) (1840,5) (1850,6.5)};
\addlegendentry{Belgique}
\end{axis}
\end{tikzpicture}
\end{center}
\legende{Production de charbon (en millions de tonnes, ordres de grandeur), 1800-1850. La courbe britannique (bleue) écrase le continent : en 1850, la Grande-Bretagne produit environ 7 fois plus que la France. Noter le démarrage précoce de la Belgique (rose), premier pays continental industrialisé.}
\end{popfigure}

\courssub{La révolution des transports : des canaux au chemin de fer}

L'industrie a besoin de transporter des masses énormes : charbon, minerai, tissus. D'abord, on creuse des \textbf{canaux} (le réseau anglais atteint environ 6 500 km vers 1830 ; en France, canal du Centre, canal de Saint-Quentin). Mais la vraie rupture, c'est le \textbf{chemin de fer}.

\begin{itemize}
\item \textbf{1825} : première ligne publique du monde, \textbf{Stockton-Darlington} (Angleterre), avec la locomotive de George Stephenson.
\item \textbf{1829} : la \emph{Rocket} de Stephenson atteint environ 47 km/h au concours de Rainhill.
\item \textbf{1830} : ligne Liverpool-Manchester, premier succès commercial.
\item \textbf{1835} : première ligne continentale, \textbf{Bruxelles-Malines} (Belgique) ; en Prusse, Nuremberg-Fürth (1835) ; en France, Paris-Saint-Germain (1837).
\end{itemize}

Le chemin de fer provoque un puissant \textbf{effet d'entraînement} : chaque kilomètre de voie exige des rails en fer puis en acier, des traverses, des locomotives, des wagons. Il stimule donc la \textbf{sidérurgie}, les \textbf{mines} et la \textbf{machine-outil} (tours, fraiseuses, machines à fabriquer des machines). Il crée aussi des emplois, unit les marchés nationaux et transforme la perception de l'espace et du temps (adoption de l'heure uniforme).

\begin{popfigure}
\begin{center}
\begin{tikzpicture}[scale=1.05, every node/.style={font=\scriptsize}]
% Europe schématique
\draw[fill=popcream, draw=popline, thick]
  (0.2,1.2) -- (0.6,2.6) -- (0.4,3.8) -- (1.2,4.6) -- (2.6,5.0) -- (4.2,4.8)
  -- (5.6,4.4) -- (6.4,3.6) -- (6.8,2.6) -- (6.2,1.6) -- (4.8,1.0) -- (3.0,0.8)
  -- (1.4,0.9) -- cycle;
\node[popmuted] at (5.9,1.2) {Europe (croquis)};
% nœuds
\coordinate (LON) at (1.1,4.0);
\coordinate (PAR) at (2.2,3.0);
\coordinate (BRU) at (2.7,3.7);
\coordinate (BER) at (4.6,3.8);
\coordinate (VIE) at (5.2,2.4);
\coordinate (MAN) at (1.0,4.4);
\coordinate (RUH) at (3.4,3.8);
% lignes
\draw[very thick, popblue] (LON) -- (BRU);
\draw[very thick, popblue] (BRU) -- (PAR);
\draw[very thick, popblue] (BRU) -- (RUH);
\draw[very thick, popblue] (RUH) -- (BER);
\draw[very thick, popblue] (BER) -- (VIE);
\draw[very thick, popblue, dashed] (PAR) -- (VIE);
\draw[very thick, popblue] (LON) -- (MAN);
% villes
\foreach \p/\n/\pos in {LON/Londres/left, PAR/Paris/below, BRU/Bruxelles/above, BER/Berlin/above, VIE/Vienne/right, MAN/Manchester/above left, RUH/Ruhr/above right}{
  \fill[poporange] (\p) circle (0.09);
  \node[\pos, popdark] at (\p) {\n};
}
% nord
\draw[-{Stealth}, thick, popdark] (0.5,0.4) -- (0.5,1.0);
\node[above] at (0.5,1.0) {N};
\node[popmuted] at (3.4,0.35) {Échelle indicative : 1 cm $\approx$ 300 km};
\end{tikzpicture}
\end{center}
\legende{Croquis du réseau ferré européen vers 1850 (schématique) : les premières lignes relient les bassins industriels et les capitales. Nœuds principaux : Londres, Bruxelles (premier réseau continental planifié par l'État), Paris, la Ruhr, Berlin, Vienne. En pointillés : liaisons encore incomplètes en 1850.}
\end{popfigure}

\begin{methode}[Construire le croquis d'une région industrielle (ex. : bassin du Nord-Pas-de-Calais ou de la Ruhr)]
\begin{enumerate}
\item \textbf{Tracer le fond} : contour simplifié de la région, flèche du nord, échelle indicative.
\item \textbf{Placer les ressources} d'abord : symbole charbon (losange noir) sur le gisement, minerai de fer (triangle) s'il existe.
\item \textbf{Placer les villes-usines} (carrés ou points nommés : Lille, Douai, Essen, Dortmund) \emph{sur ou près} du gisement.
\item \textbf{Dessiner les transports} : voies ferrées (trait noir double ou fléché), canaux (trait bleu), ports.
\item \textbf{Organiser la légende} dans un cadre, dans le même ordre que le plan de l'exposé : ressources $\rightarrow$ industries $\rightarrow$ transports $\rightarrow$ villes.
\item \textbf{Titrer} le croquis (lieu + date) : « Le bassin houiller de la Ruhr vers 1850 ».
\end{enumerate}
\end{methode}

\courssub{Les nouvelles organisations du travail et la question sociale}

L'usine bouleverse la condition des travailleurs. Le \textbf{prolétariat} naît : des ouvriers qui ne possèdent rien d'autre que leur force de travail, vendue pour un salaire. La \textbf{division du travail} parcellise les gestes : l'ouvrier répète le même mouvement 12 à 15 heures par jour, femmes et enfants compris (les enfants, dès 6-7 ans, dans les mines et les filatures). La \textbf{misère ouvrière} est documentée par les enquêtes officielles et par des observateurs comme Friedrich Engels (\emph{La Situation de la classe laborieuse en Angleterre}, 1845) : taudis insalubres, épidémies de choléra, espérance de vie très faible dans les quartiers ouvriers de Manchester.

Face à cette misère, les ouvriers ripostent :
\begin{itemize}
\item le \cle{luddisme} (1811-1816) : les ouvriers brisent les machines, accusées de voler le travail (mouvement attribué au mythique « général Ludd ») ;
\item les premières \textbf{grèves} et \emph{trade-unions} (syndicats), longtemps interdites puis tolérées en Angleterre après 1824 ;
\item le \cle{chartisme} (1838-1848) : grand mouvement politique ouvrier anglais qui réclame, par une « Charte du peuple », le suffrage universel masculin et la réforme du Parlement. Il échoue sur le moment, mais annonce la démocratie moderne.
\end{itemize}

L'État finit par légiférer : les \cle{Factory Acts} britanniques interdisent le travail des enfants de moins de 9 ans dans les filatures (1833) et limitent la journée des femmes et des adolescents à 10 heures (1847). En France, une loi de 1841 limite le travail des enfants. Ce sont les toutes premières lois sociales.

\begin{exoresolu}[Analyser un document : la condition ouvrière]
\textbf{Énoncé.} Document : « Dans les filatures de coton de Manchester, les enfants commencent à 6 ans. La journée dure 14 heures. Les ouvriers habitent des caves humides, à plusieurs familles par pièce. Le choléra de 1832 y a fait des ravages. » (Enquête parlementaire britannique, années 1830.) Questions : 1) Relevez deux éléments montrant la misère ouvrière. 2) Expliquez pourquoi l'usine concentre ces problèmes. 3) Citez une réponse des ouvriers et une réponse de l'État.
\tcblower
\textbf{Solution détaillée.}
1) La misère apparaît dans le \textbf{travail des enfants dès 6 ans} et la durée de la journée (\textbf{14 heures}), ainsi que dans l'habitat : « caves humides », surpeuplement (« plusieurs familles par pièce »), épidémies (choléra, 1832).
2) L'usine concentre des milliers de salariés dans des villes qui grossissent sans planification : les logements construits à la hâte, sans eau ni égouts, favorisent les épidémies. De plus, la concurrence entre ouvriers nombreux et sans qualification permet au patron d'imposer bas salaires et longues journées.
3) Réponse ouvrière : le \textbf{chartisme} (1838-1848), qui réclame le suffrage universel, ou le bris de machines (luddisme). Réponse de l'État : les \textbf{Factory Acts} (1833 : interdiction du travail avant 9 ans ; 1847 : journée de 10 heures pour femmes et adolescents).
\end{exoresolu}

\courssub{L'impact démographique et urbain}

L'industrialisation s'accompagne d'une \textbf{explosion démographique et urbaine}. La \cle{transition démographique} commence : la mortalité recule (meilleure alimentation, hygiène, progrès médicaux) tandis que la natalité reste d'abord élevée, d'où une forte croissance de la population. Les villes industrielles grandissent à une vitesse vertigineuse :

\begin{center}
\begin{tabularx}{\linewidth}{|l|X|X|}
\hline
\rowcolor{popblueL} \textbf{Ville} & \textbf{Vers 1800} & \textbf{Vers 1850} \\
\hline
Manchester & environ 75 000 hab. & environ 400 000 hab. \\
\hline
Londres & environ 1 million & plus de 2,3 millions \\
\hline
Lille & environ 55 000 hab. & environ 100 000 hab. \\
\hline
Essen (Ruhr) & bourg de 4 000 hab. & ville industrielle en pleine croissance (Krupp) \\
\hline
\end{tabularx}
\end{center}

Cette urbanisation sauvage crée des villes noires de suie, sans égouts ni eau potable, où les épidémies de choléra (1832, 1849) font des milliers de morts. Les grands travaux d'assainissement ne viendront qu'après 1850.

\begin{attention}[Erreur fréquente à l'examen]
Ne dis jamais que « la France n'a pas fait sa révolution industrielle ». Elle l'a faite, mais \emph{doucement} : croissance régulière sans rupture brutale, avec un \textbf{dualisme persistant} entre un artisanat de luxe dynamique (soieries lyonnaises, mode parisienne) et une industrie lourde concentrée en quelques points (Nord, Le Creusot). Il s'agit d'un \textbf{déclin relatif} par rapport à la Grande-Bretagne, non d'une absence d'industrialisation.
\end{attention}

\begin{savaistu}[Et le Cameroun dans tout cela ?]
Les premières locomotives qui ont roulé au Cameroun, sur la ligne \textbf{Douala-Nkongsamba} (construite à partir de 1906 sous administration allemande), étaient des machines fabriquées par les firmes allemandes \emph{Krauss} et \emph{Henschel} : des fruits directs de la Première Révolution industrielle (acier, vapeur, machine-outil) que nous venons d'étudier. Le chemin de fer, né à Stockton-Darlington en 1825, atteint ainsi l'Afrique centrale moins d'un siècle plus tard — preuve que la révolution industrielle européenne est aussi le point de départ de l'expansion coloniale que nous étudierons dans la section IV.
\end{savaistu}

\begin{aretenir}[L'essentiel de la section]
\begin{itemize}
\item La \textbf{Première Révolution industrielle} (vers 1780-1850) repose sur la triade \textbf{charbon -- fer -- vapeur} et sur le système d'\textbf{usine}.
\item Elle naît en \textbf{Grande-Bretagne} grâce à trois préconditions : révolution agraire, accumulation capitaliste (commerce colonial et traite), progrès techniques (machine de Watt, 1769).
\item Elle se diffuse inégalement : \textbf{Belgique} précoce (Sambre-et-Meuse), \textbf{France} lente et dualiste, \textbf{États allemands} boostés par le \textbf{Zollverein} (1834).
\item Le \textbf{chemin de fer} (Stockton-Darlington, 1825) entraîne sidérurgie et machine-outil.
\item Conséquences sociales majeures : naissance du \textbf{prolétariat}, misère ouvrière, premières luttes (\textbf{luddisme}, \textbf{chartisme}) et premières lois sociales (\textbf{Factory Acts}).
\end{itemize}
\end{aretenir}

\coursec{La Deuxième Révolution Industrielle : acier, électricité, pétrole, chimie (1850-1914)}

\courssub{Transition : de la vapeur à l'énergie électrique et au pétrole}

La Première Révolution industrielle (vers 1780-1850) a reposé sur le couple \textbf{charbon -- vapeur -- fer}. Elle a transformé l'Angleterre en « atelier du monde » et essaimé en Belgique, dans le Nord de la France, en Rhénanie. Mais dès les années 1850, ses limites apparaissent : le rendement de la machine à vapeur plafonne (10 à 15 \%), le fer puddlé reste fragile pour les grands ouvrages (rails, coques, canons), et le transport par canaux puis chemins de fer à voie unique sature.

La Deuxième Révolution industrielle (1850-1914) change de paradigme : \textbf{nouvelles sources d'énergie} (électricité, pétrole, gaz), \textbf{nouvelles matières} (acier, aluminium, produits chimiques de synthèse), \textbf{nouvelles industries} (chimie, électrotechnique, automobile, construction navale moderne), \textbf{nouvelle organisation du travail} (taylorisme, fordisme) et \textbf{nouvelle concentration du capital} (cartels, trusts, banques d'affaires). L'Allemagne et les États-Unis deviennent les leaders ; la France se spécialise ; la Russie décolle grâce aux capitaux français. Cette mutation économique alimente directement l'impérialisme (section IV) : il faut des débouchés, des matières premières, des champs d'investissement.

\begin{popfigure}
\begin{tikzpicture}[scale=0.9, every node/.style={font=\small}]
\begin{axis}[
    width=\linewidth,
    height=0.55\linewidth,
    xlabel={Année},
    ylabel={Production d'acier (millions de tonnes)},
    xmin=1870, xmax=1913,
    ymin=0, ymax=35,
    xtick={1870,1880,1890,1900,1910,1913},
    ytick={0,5,10,15,20,25,30,35},
    legend style={at={(0.02,0.98)}, anchor=north west, font=\tiny, fill=white, draw=popline},
    grid=major, grid style={popline!40},
    axis line style={popdark},
    tick style={popdark},
    label style={popdark},
]
\addplot[popblue, thick, mark=*] coordinates {
    (1870,0.7) (1880,1.3) (1890,3.6) (1900,6.7) (1910,13.5) (1913,17.6)
};
\addlegendentry{Allemagne}
\addplot[poporange, thick, mark=square*] coordinates {
    (1870,0.6) (1880,1.2) (1890,4.3) (1900,10.3) (1910,26.1) (1913,31.8)
};
\addlegendentry{États-Unis}
\addplot[popgreen, thick, mark=triangle*] coordinates {
    (1870,0.5) (1880,1.0) (1890,3.0) (1900,5.0) (1910,6.5) (1913,7.7)
};
\addlegendentry{Grande-Bretagne}
\addplot[poppurple, thick, mark=diamond*] coordinates {
    (1870,0.4) (1880,0.7) (1890,1.5) (1900,1.6) (1910,3.6) (1913,4.6)
};
\addlegendentry{France}
\addplot[poppink, thick, mark=pentagon*] coordinates {
    (1870,0.2) (1880,0.4) (1890,1.0) (1900,2.2) (1910,4.8) (1913,5.5)
};
\addlegendentry{Russie}
\end{axis}
\end{tikzpicture}
\legende{Figure 1 -- Production d'acier des grandes puissances (1870-1913). Source : Bairoch, \emph{Victoires et déboires}, adapté. Le croisement Allemagne / Grande-Bretagne intervient vers 1900 ; les États-Unis dominent sans partage dès 1890.}
\end{popfigure}

\courssub{Nouvelles sources d'énergie : électricité, pétrole, gaz}

\textbf{Électricité : de la dynamo au réseau.}\par
La pile (Volta, 1800) donne un courant continu mais faible. La \textbf{dynamo} (Gramme, 1871) produit un courant industriel continu. Le \textbf{transformateur} (Gaulard-Gibbs, 1883) et le \textbf{courant alternatif} (Tesla, Westinghouse) permettent le transport à haute tension sur de longues distances avec peu de pertes. Edison (1879) met au point l'ampoule à incandescence et crée le premier réseau de distribution (Pearl Street, New York, 1882). Siemens (Allemagne) et Thomson-Houston (France/USA) construisent les premiers réseaux urbains et industriels.

\begin{definition}[Électricité industrielle]
L'électricité devient une \textbf{énergie vectorielle} : elle se transporte, se distribue et se convertit en travail mécanique (moteurs), en chaleur (four électrique), en lumière, en signaux (télégraphe, téléphone). Elle libère l'usine de la contrainte de l'arbre de transmission central (vapeur) : chaque machine a son moteur (\emph{unit drive}), ce qui rend l'atelier flexible.
\end{definition}

\textbf{Pétrole : du kérosène au moteur à explosion.}\par
Le forage de Drake (Titusville, Pennsylvanie, 1859) lance l'exploitation industrielle. Le \textbf{kérosène} (éclairage) domine d'abord. L'invention du \textbf{moteur à explosion} (Otto, 1876 ; cycle 4 temps) puis du \textbf{moteur Diesel} (1892, allumage par compression, rendement ~30 \%) et du \textbf{moteur à essence} (Daimler, Benz, 1885-1886) ouvre l'ère de l'automobile, du camion, du sous-marin, de l'aviation (1903, frères Wright). Le pétrole devient énergie stratégique : Standard Oil (Rockefeller) contrôle 90 \% du raffinage US en 1880 ; Royal Dutch Shell (1907) et la Anglo-Persian Oil (1909, future BP) sécurisent l'approvisionnement britannique.

\textbf{Gaz : énergie urbaine.}\par
Le gaz de houille (cokéfaction) éclaire les villes dès 1810 (Londres, Paris). Le gaz naturel (gisements USA, Russie) reste marginal avant 1914. Le gaz alimente l'éclairage public, le chauffage domestique et certains fours industriels.

\begin{attention}[Déclin relatif ≠ disparition]
Le charbon reste la \textbf{première énergie mondiale} en 1913 (environ 70 \% de la consommation primaire). La vapeur propulse encore 95 \% de la marine marchande et 100 \% des chemins de fer. L'électricité et le pétrole \emph{complètent} puis \emph{supplantent} la vapeur seulement dans les usages à haute valeur ajoutée (moteurs individuels, éclairage, chimie, transport routier/automobile). Ne pas écrire « le charbon a disparu ».
\end{attention}

\courssub{Nouvelles matières : acier, aluminium, chimie de synthèse}

\textbf{L'acier : du luxe au matériau de masse.}\par
Le fer puddlé (Première révolution) est malléable mais hétérogène. L'acier (0,2 à 2 \% carbone) allie dureté, élasticité et ténacité. Trois procédés le rendent abordable :
\begin{itemize}
    \item \textbf{Bessemer (1855)} : soufflage d'air dans la fonte en convertisseur basculant (acide). Rapide (20 min), mais ne traite que les fontes pauvres en phosphore (Angleterre, Suède, USA).
    \item \textbf{Martin-Siemens (1864)} : four réverbératoire à régénérateurs (chaleur récupérée). Plus lent (6-8 h), mais contrôle précis, accepte ferraille et fontes diverses. Dominant en France (Lorraine) et USA.
    \item \textbf{Thomas-Gilchrist (1878)} : convertisseur \textbf{basique} (revêtement dolomite + chaux) qui capte le phosphore (scories phosphatées = engrais). Rend exploitables les minerais \textbf{phosphoreux} de \textbf{Lorraine (Minette)} et de \textbf{Ruhr}. L'Allemagne devient premier producteur européen d'acier vers 1900.
\end{itemize}

\begin{popfigure}
\begin{tikzpicture}[scale=0.85, every node/.style={font=\small}]
% Convertisseur Thomas - vue de coupe simplifiée
\draw[popdark, thick] (0,0) -- (6,0) -- (6,8) -- (0,8) -- cycle; % enceinte
\draw[popdark, thick] (1,8) -- (1,10) -- (5,10) -- (5,8); % col
\draw[poporange, thick, ->] (3,10) -- (3,11.5) node[above] {Gaz de combustion (CO, CO$_2$, P$_2$O$_5$)};
\draw[popblue, thick, ->] (3,-1) -- (3,0) node[below] {Air comprimé (tuyères)};
\draw[popgreen, thick, ->] (7,4) -- (6,4) node[left] {Fonte phosphoreuse liquide};
\draw[poppurple, thick, ->] (3,4) -- (3,2) node[below] {Acier liquide (coulée)};
\draw[popmuted, thick, ->] (-1,4) -- (0,4) node[right] {Chaux + dolomite (fondants)};

% Tuyères
\foreach \x in {1.5,2.5,3.5,4.5} {
    \draw[popdark, thick] (\x,0) -- (\x,-0.5);
    \draw[popblue] (\x,-0.5) circle (0.15);
}

% Zones internes
\fill[poporange!20] (0.5,0.5) rectangle (5.5,3.5); % bain
\fill[popgreen!20] (0.5,3.5) rectangle (5.5,5.5); % réaction
\fill[poppurple!20] (0.5,5.5) rectangle (5.5,7.5); % scories

\node[popdark] at (3,2) {\textbf{Bain de fonte + acier}};
\node[popdark] at (3,4.5) {\textbf{Zone de décarburation / déphosphoration}};
\node[popdark] at (3,6.5) {\textbf{Scories basiques (CaO, P$_2$O$_5$) $\to$ engrais}};
\node[popdark, align=center] at (3,9) {\textbf{Convertisseur Thomas (basique)}\\\emph{Revêtement dolomitique + fondants basiques}};
\end{tikzpicture}
\legende{Figure 2 -- Schéma du convertisseur Thomas (1878). L'air comprimé traverse la fonte phosphoreuse ; le phosphore passe dans les scories basiques (engrais), l'acier est produit en 15-20 min à bas coût. Rails, coques de navires, canons deviennent abordables.}
\end{popfigure}

\textbf{Aluminium : le métal de l'avenir.}\par
Découvert (Wöhler 1827), il reste cher (or) jusqu'à l'électrolyse (Hall-Héroult 1886, indépendamment). L'électricité bon marché (barrages alpins, rhénans) permet la production de masse (Pechiney France, Alcoa USA, VAW Allemagne). Léger, inoxydable, bon conducteur : aviation, câbles haute tension, emballages.

\textbf{Chimie de synthèse : la science au service de l'industrie.}\par
L'Allemagne domine (BASF, Bayer, Hoechst, AGFA). Trois piliers :
\begin{itemize}
    \item \textbf{Teintures} (mauvéine Perkin 1856, alizarine synthétique 1869, indigo 1897) : ruinent les cultures d'indigo et de garance coloniales.
    \item \textbf{Engrais} : acide sulfurique (contact), soude (Solvay 1863), superphosphates (scories Thomas), puis \textbf{ammoniac de synthèse} (Haber-Bosch 1913, BASF) : fixent l'azote atmosphérique $\to$ engrais azotés $\to$ explosion des rendements agricoles.
    \item \textbf{Explosifs} : nitroglycérine (Sobrero 1847), dynamite (Nobel 1867), TNT (1902). Mines, travaux publics (Suez 1869, Panama 1914), armement.
\end{itemize}

\begin{savaistu}[Le procédé Thomas change la carte de l'Europe]
Sans Thomas (1878), la Minette lorraine (38 \% Fe, 1-2 \% P) et les minerais de Ruhr restent inexploités. L'Allemagne passe de 0,7 Mt acier (1870) à 17,6 Mt (1913) ; la France de 0,4 à 4,6 Mt. L'acier bon marché alimente la course aux armements (cuirassés Dreadnought, canons Krupp, rails stratégiques). En 1914, le front occidental se stabilise sur la ligne des gisements (Lorraine, Longwy-Briey). La géologie a fait la stratégie.
\end{savaistu}

\courssub{Nouvelles industries : construction mécanique, automobile, électrotechnique}

\textbf{Construction mécanique : naval et ferroviaire.}\par
L'acier + la vapeur (puis turbine à vapeur Parsons 1884, turbine à gaz) révolutionnent la marine : coques en acier, propulsion hélice, chaudières à tubes d'eau, turbines. Tonnage mondial x10 (1850-1914). Paquebots (Mauretania, Titanic), cuirassés (Dreadnought 1906). Chemins de fer : rails acier (durée x5), locomotives plus puissantes, réseau européen dense (200 000 km en 1914), Transsibérien (1891-1904), Berlin-Bagdad (projet).

\textbf{Automobile : du luxe à la grande série.}\par
Benz (1885, tricycle), Daimler (1886, quadricycle), Peugeot (1891, moteur Daimler), Renault (1898, voiturette). France leader jusqu'en 1906 (30 000 véhicules/an). Puis USA : Olds (1901, série), Ford (Modèle T 1908, 15 millions d'exemplaires en 1927). L'automobile tire l'acier, le caoutchouc (pneumatiques Michelin, Dunlop), le pétrole, le verre, la machine-outil.

\textbf{Électrotechnique : l'industrie de la fée électricité.}\par
Siemens (Allemagne, 1847 télégraphe, 1866 dynamo, 1879 tramway), AEG (Rathenau, 1883), Thomson-Houston (France/USA), Westinghouse, General Electric (Edison, 1892). Production : dynamos, moteurs, transformateurs, câbles, lampes, tramways, signalisation ferroviaire, électrolyse (aluminium, chlore, soude). L'Allemagne exporte 50 \% de la production mondiale d'équipement électrique en 1913.

\begin{exoresolu}[Lire un graphique de production industrielle]
\textbf{Énoncé :} À l'aide de la Figure 1, décrire et expliquer l'évolution de la production d'acier de l'Allemagne et de la Grande-Bretagne entre 1870 et 1913.

\textbf{Solution détaillée :}
\begin{enumerate}
    \item \textbf{Lecture des échelles :} Axe horizontal : années 1870-1913. Axe vertical : millions de tonnes (Mt). Échelle linéaire.
    \item \textbf{Identification des courbes :} Courbe bleue (Allemagne) : départ 0,7 Mt (1870) $\to$ 17,6 Mt (1913). Courbe verte (GB) : départ 0,5 Mt $\to$ 7,7 Mt. Croisement vers 1900.
    \item \textbf{Ruptures de pente :} Allemagne : accélération nette 1890-1900 (Thomas + unification + marché intérieur + protectionnisme Bismarck 1879). GB : croissance lente, quasi plate après 1900 (investissements à l'étranger, conservatisme technique, mines épuisées).
    \item \textbf{Hypothèse explicative :} L'Allemagne rattrape et dépasse la GB grâce au \textbf{procédé Thomas} (minerais phosphoreux locaux), à l'intégration verticale (mines-usines-banques), au marché protégé (Zollverein puis Empire), et à l'effort d'armement naval (lois de flotte 1898, 1900). La GB, pionnière, subit la \textbf{loi des rendements décroissants} sur ses gisements et sous-investit dans les nouveaux procédés (Martin, Thomas).
    \item \textbf{Conclusion :} En 1913, l'Allemagne est le premier producteur européen (17,6 Mt), la GB chute au 3\textsuperscript{e} rang mondial derrière USA (31,8 Mt) et Allemagne. Le graphique illustre le basculement de la leadership industriel.
\end{enumerate}
\end{exoresolu}

\courssub{Nouvelle organisation : taylorisme, fordisme, concentration capitaliste}

\textbf{Taylorisme : l'Organisation Scientifique du Travail (OST).}\par
Frederick Winslow Taylor (\emph{The Principles of Scientific Management}, 1911) observe l'inefficacité des ateliers : « soldiering » (ralentissement volontaire), gestes parasites, outillage inadapté. Il propose :
\begin{enumerate}
    \item \textbf{Décomposition} de chaque tâche en mouvements élémentaires (chronométrage).
    \item \textbf{Séparation conception / exécution} : le bureau d'études (ingénieurs, chronométreurs) définit \emph{la} « one best way » ; l'ouvrier exécute sans réfléchir.
    \item \textbf{Standardisation} : outillage, temps, conditions.
    \item \textbf{Paiement au rendement} (piece-rate) : salaire = taux de base + prime si temps < temps standard.
    \item \textbf{Planification} : carte d'instruction, fiche de poste, planning centralisé.
\end{enumerate}

\begin{definition}[Taylorisme (OST)]
Système d'organisation du travail fondé sur la \textbf{rationalisation scientifique} des gestes : analyse chronométrique, élimination des temps morts, séparation stricte entre travail manuel (exécution) et travail intellectuel (conception/planification), incitation financière individuelle au rendement. Objectif : productivité maximale, coût unitaire minimal.
\end{definition}

\textbf{Fordisme : la chaîne de montage.}\par
Henry Ford (Highland Park, 1913) applique le taylorisme à l'échelle de l'usine entière : \textbf{chaîne mobile} (tapis roulant) qui amène le châssis à l'ouvrier fixe. Temps de montage châssis : 12 h $\to$ 1 h 33. Standardisation absolue (Modèle T unique, noir). Salaires doublés (5 \$/jour, 1914) pour fidéliser et créer un marché de masse (ouvriers acheteurs). Productivité x8, prix divisé par 4 (850 \$ $\to$ 260 \$).

\begin{methode}[Analyser un modèle d'organisation industrielle]
\begin{enumerate}
    \item Identifier le \textbf{produit} (standardisé ? diversifié ?) et le \textbf{marché} (masse ? niche ?).
    \item Repérer la \textbf{division du travail} : tâches simples (taylorisme) ou poste fixe / chaîne (fordisme) ou îlots autonomes (toyotisme, hors programme).
    \item Étudier la \textbf{gestion des flux} : stocks (zéro stock ?), approvisionnement (juste-à-temps ?), transport interne.
    \item Analyser la \textbf{relation salariale} : temps, pièce, prime, participation, qualification requise.
    \item Conclure sur le \textbf{rapport capital/travail} : contrôle, conflictualité, productivité, partage des gains.
\end{enumerate}
\end{methode}

\textbf{Concentration capitaliste : cartels, trusts, holdings.}\par
La taille critique des usines (aciéries, chimie, électrotechnique) exige des capitaux massifs. La \textbf{Grande Dépression} (1873-1896, voir Propriété) pousse à limiter la concurrence destructrice.

\begin{center}
\begin{tabularx}{\linewidth}{|l|X|X|}
\hline
\rowcolor{popblueL} \textbf{Forme} & \textbf{Mécanisme} & \textbf{Exemples (vers 1914)} \\
\hline
\textbf{Cartel} & Entente contractuelle entre entreprises \emph{juridiquement indépendantes} : fixation des prix, partage des marchés, quotas. Réversible. & \emph{Rheinstahl} (acier Ruhr), \emph{Rhenish-Westphalian Coal Syndicate}, \emph{Comptoir des fers} (France). Privilégié en \textbf{Allemagne}. \\
\hline
\textbf{Trust} & \textbf{Fusion juridique} : les entreprises apportent leurs actifs à une société holding qui les contrôle totalement. Centralisation totale. & \emph{Standard Oil} (Rockefeller, 1882, démantelé 1911), \emph{U.S. Steel} (Morgan, 1901, 1 Md \$), \emph{General Electric}, \emph{Du Pont}. Modèle \textbf{USA}. \\
\hline
\textbf{Holding / Konzern} & Société mère (holding) détient la majorité du capital de filiales opérationnelles. Contrôle financier sans fusion opérationnelle. & \emph{Krupp} (Allemagne, famille + holding), \emph{Schneider} (France, Creusot), \emph{Stinnes} (charbon/électricité/presse). Modèle \textbf{européen}. \\
\hline
\end{tabularx}
\end{center}

\begin{attention}[Cartel vs Trust : confusion fréquente]
\textbf{Cartel} = accord \emph{horizontal} (concurrents s'entendent) ou \emph{vertical} (producteurs/distributeurs), \textbf{indépendance juridique maintenue}. \textbf{Trust} = \textbf{fusion}, disparition des personnes morales absorbées, \textbf{centralisation totale} sous une direction unique. L'Allemagne (droit des sociétés souple, banques universelles) favorise le cartel/Konzern ; les USA (droit des corporations, antitrust Sherman Act 1890) poussent au trust (fusion) puis le combattent. Ne pas écrire « cartel allemand = trust allemand ».
\end{attention}

\textbf{Banques d'affaires et banques universelles.}\par
Le financement de la croissance (usines, réseaux, trusts) dépasse les capacités des banques de dépôt. \textbf{Crédit Mobilier} (Pereire, France, 1852) : banque d'affaires, émission d'obligations, participation au capital (railways, immobilier, Suez). Faillite 1867, mais modèle imité. \textbf{Deutsche Bank} (1870), \textbf{Dresdner Bank}, \textbf{Commerzbank} : \textbf{banques universelles} (dépôts + crédit à long terme + prise de participation + émission). Siègent aux conseils d'administration (Aufsichtsrat) $\to$ contrôle stratégique, stabilité, financement de l'export (Bagdadbahn, Ottoman railways). En France, banques de dépôt (Crédit Lyonnais, Société Générale) + banques d'affaires (Paribas, Lazard) restent plus séparées.

\begin{popfigure}
\begin{tikzpicture}[scale=0.85, every node/.style={font=\small},
    box/.style={draw=popdark, thick, rounded corners, fill=popcream, minimum width=2.2cm, minimum height=1cm, align=center},
    arrow/.style={-Stealth, thick, popdark},
    doublearrow/.style={Stealth-Stealth, thick, poporange},
]
% Holding
\node[box, fill=popblue!20, align=center] (H) at (0,0){\textbf{HOLDING}\\(Krupp / Standard Oil)\\\emph{Contrôle financier \& stratégique}};

% Filiales - niveau 1
\node[box, align=center] (Mines) at (-5,-2.5){\textbf{MINES}\\(Charbon, fer, pétrole)};
\node[box, align=center] (Usines) at (0,-2.5){\textbf{USINES}\\(Acier, chimie, mécanique)};
\node[box, align=center] (Banques) at (5,-2.5){\textbf{BANQUES}\\(Crédit, émission, trésorerie)};
\node[box, align=center] (Transport) at (2.5,-4.5){\textbf{TRANSPORT}\\(Chemins de fer, flotte, ports)};

% Flèches holding -> filiales
\draw[arrow] (H.south west) -- (Mines.north);
\draw[arrow] (H.south) -- (Usines.north);
\draw[arrow] (H.south east) -- (Banques.north);

% Intégration verticale
\draw[doublearrow] (Mines.east) -- (Usines.west) node[midway, above, font=\tiny] {Matières premières};
\draw[doublearrow] (Usines.east) -- (Banques.west) node[midway, above, font=\tiny] {Financement / Trésorerie};
\draw[doublearrow] (Usines.south east) -- (Transport.north west) node[midway, right, font=\tiny] {Produits finis};

% Intégration horizontale (ex: plusieurs usines)
\node[box, dashed, fill=poporange!10, align=center] (Usines2) at (-2.5,-4.5){\textbf{USINES 2}\\(Région / Spécialité)};
\draw[doublearrow, dashed] (Usines.south west) -- (Usines2.north) node[midway, left, font=\tiny] {Spécialisation / Partage marché};

% Légende
\node[align=left, font=\tiny, popmuted] at (7.5,0) {
\textbf{Légende :}\\
Trait plein = Contrôle (Holding $\to$ Filiales)\\
Double flèche = Flux internes (intégration verticale)\\
Trait pointillé = Intégration horizontale (cartellisation)
};
\end{tikzpicture}
\legende{Figure 3 -- Organigramme type d'un trust / Konzern (ex. Krupp, Standard Oil, U.S. Steel). La holding centralise la stratégie, le financement et le contrôle. Les filiales assurent l'intégration verticale (mines $\to$ usines $\to$ transport $\to$ banques) et horizontale (partage de marchés, spécialisation).}
\end{popfigure}

\courssub{Nouveaux acteurs : Allemagne, USA, France, Russie}

\begin{center}
\begin{tabularx}{\linewidth}{|l|X|X|X|}
\hline
\rowcolor{popblueL} \textbf{Pays} & \textbf{Spécialisations / Leadership} & \textbf{Atouts structurels} & \textbf{Faiblesses / Particularismes} \\
\hline
\textbf{Allemagne} & \textbf{Chimie} (teintures, engrais, pharma, explosifs), \textbf{Électricité} (Siemens, AEG), \textbf{Acier Thomas} (Ruhr, Sarre), \textbf{Machine-outil}, \textbf{Optique} (Zeiss). & Unification (1871), marché intérieur protégé (Tarif Bismarck 1879), banques universelles, enseignement technique supérieur (TH), recherche universitaire-industrielle, main-d'œuvre qualifiée. & Dépendance importations (coton, caoutchouc, pétrole, minerais non phosphoreux), tensions sociales (SPD, syndicats), encerclement diplomatique. \\
\hline
\textbf{États-Unis} & \textbf{Pétrole} (Standard Oil), \textbf{Automobile} (Ford, GM), \textbf{Électricité} (GE, Westinghouse), \textbf{Acier} (Bessemer, puis Martin, US Steel), \textbf{Agro-industrie}, \textbf{Machine-outil} (système américain). & Ressources immenses (charbon, fer, pétrole, cuivre), marché intérieur continental unique, immigration (main-d'œuvre), esprit d'entreprise, brevets, protectionnisme élevé (tarifs McKinley 1890, Dingley 1897, Payne-Aldrich 1909). & Dualisme Nord industriel / Sud agricole, trusts contestés (Sherman Act 1890, Clayton Act 1914), grèves violentes (Homestead, Pullman). \\
\hline
\textbf{France} & \textbf{Automobile} (Peugeot, Renault, Citroën), \textbf{Aéronautique} (naissante), \textbf{Produits de luxe}, \textbf{Électricité} (Thomson-Houston, Alsthom), \textbf{Chimie fine} (Rhône-Poulenc). & Épargne abondante (placement en obligations russes, ottomanes), ingénieurs de haut niveau (Écoles), colonies (matières premières), protectionnisme modéré (Méline 1892). & Démographie stagnante, capitalisme familial / réticent à la concentration, structure PME, retard chimie lourde / électricité lourde vs Allemagne. \\
\hline
\textbf{Russie} & \textbf{Décollage tardif} (1890-1913 : +8 \%/an), \textbf{Charbon} (Donetz), \textbf{Acier} (Donetz, Urals), \textbf{Pétrole} (Bakou), \textbf{Chemins de fer} (Transsibérien). & Immensité des ressources, main-d'œuvre bon marché, capitaux français (obligations), État dirigiste (Witte, 1892-1903 : étalon-or, tarifs protecteurs, monopole vodka). & Structure agraire archaïque (mirs), dépendance technologique/financière extérieure, répression politique, révolution 1905, instabilité. \\
\hline
\end{tabularx}
\end{center}

\begin{exemplebox}[Repères chiffrés 1913 (année de référence avant-guerre)]
\begin{itemize}
    \item \textbf{Production d'acier (Mt) :} USA 31,8 -- Allemagne 17,6 -- GB 7,7 -- France 4,6 -- Russie 5,5. (Allemagne = 1\textsuperscript{er} Europe, 2\textsuperscript{e} monde).
    \item \textbf{Production d'électricité (TWh) :} USA 38 -- Allemagne 8,5 -- GB 6,5 -- France 3,5.
    \item \textbf{Consommation d'énergie par habitant (teq charbon) :} USA 7,0 -- GB 5,5 -- Allemagne 3,5 -- France 2,0 -- Russie 0,5.
    \item \textbf{Part dans l'industrie manufacturière mondiale (Bairoch) :} USA 32,0 \% -- Allemagne 14,8 \% -- GB 13,6 \% -- France 6,1 \% -- Russie 8,2 \%.
    \item \textbf{Exportations de capitaux (milliards \$ or) :} GB 18 -- France 12 -- Allemagne 5 -- USA (importateur net jusqu'en 1914).
\end{itemize}
\end{exemplebox}

\courssub{Crises cycliques, protectionnisme et course aux débouchés}

\textbf{La « Grande Dépression » (1873-1896) : déflation, pas récession.}\par
Le krach de la Bourse de Vienne (mai 1873) déclenche une crise bancaire internationale. Mais la production \emph{continue de croître} (nouvelles techniques, nouveaux pays). Ce sont les \textbf{prix} qui baissent (déflation ~1,5 \%/an) : gains de productivité (acier, chimie, transport), offre excédentaire (blé US/Argentine/Russie, acier, textile). Les profits se compriment, les salaires réels montent (déflation), les taux d'intérêt réels sont élevés.

\begin{propriete}[La « Grande Dépression » 1873-1896]
Ce n'est \textbf{pas} une récession continue (PIB en baisse), mais une \textbf{déflation des prix} structurelle due aux gains de productivité massifs (acier, chimie, transport, agriculture mécanisée) et à l'élargissement de l'offre mondiale. Conséquences majeures : (1) compression des profits $\to$ concentration (cartels, trusts) pour rétablir le pouvoir de prix ; (2) protectionnisme défensif (retour aux tarifs) ; (3) recherche de débouchés coloniaux « captifs » (impérialisme) ; (4) montée du mouvement ouvrier (revendications salariales, grèves).
\end{propriete}

\textbf{Protectionnisme : la fin du libre-échange (Cobden-Chevalier 1860).}\par
\begin{itemize}
    \item \textbf{Allemagne :} Tarif Bismarck (1879) : droits sur céréales + fer/acier. Alliance « fer et seigle » (Junkers + industriels Ruhr). Renforcé 1887, 1902 (Bülow).
    \item \textbf{France :} Loi Méline (1892) : tarifs protecteurs agriculture + industrie. Fin du traité de commerce avec GB. Préférence coloniale (1897, 1906).
    \item \textbf{USA :} Tarifs très élevés (McKinley 1890, Dingley 1897, Payne-Aldrich 1909). Marché intérieur protégé = base des trusts.
    \item \textbf{GB :} Reste libre-échangiste jusqu'en 1915 (guerre), mais subit l'érosion de sa part de marché.
    \item \textbf{Russie :} Tarifs Witte (1891, 1903) : protection forte + étalon-or (1897) pour attirer capitaux français.
\end{itemize}

\begin{methode}[Interpréter un graphique de production industrielle (révision)]
\begin{enumerate}
    \item \textbf{Identifier} : variables (axes), unités, période, pays/courbes.
    \item \textbf{Repérer} : tendances (croissance, stagnation, rupture), croisements, écarts.
    \item \textbf{Corréler} : ruptures $\leftrightarrow$ événements (guerre, innovation, loi tarifaire, krach, découverte gisement).
    \item \textbf{Comparer} : rythmes (pente), niveaux (ordre de grandeur), parts relatives.
    \item \textbf{Hypothétiser} : facteurs d'offre (technique, ressources, organisation) et de demande (démographie, armement, infrastructures, colonies).
    \item \textbf{Conclure} : hiérarchie des puissances, basculements, liens avec l'impérialisme (section IV).
\end{enumerate}
\end{methode}

\courssub{Bilan 1914 : l'Europe, atelier du monde et gendarme colonial}

En 1914, l'Europe (avec les USA) concentre \textbf{les 2/3 de la production industrielle mondiale}. L'Allemagne a dépassé la Grande-Bretagne en acier, chimie, électricité ; les USA dominent en volume absolu. La France tient son rang par la spécialisation (luxe, auto naissante, finance). La Russie émerge. Cette puissance industrielle se traduit par :
\begin{itemize}
    \item \textbf{Supériorité militaire :} artillerie Krupp, cuirassés Dreadnought, mitrailleuses, sous-marins, aviation naissante, logistique ferroviaire.
    \item \textbf{Domination financière :} Londres placeur mondial, Paris prêteur de la Russie et des Balkans, Berlin finance l'Empire ottoman (Bagdadbahn).
    \item \textbf{Partage de l'Afrique (Conférence de Berlin 1885) :} matières premières (caoutchouc, cuivre, coton, minerais), débouchés captifs, champs d'investissement (chemins de fer, mines, plantations).
    \item \textbf{Fragilités :} interdépendance économique (blé, coton, pétrole, finance) qui rend la guerre longue possible mais catastrophique ; tensions sociales (grèves, suffrage universel, socialisme) ; rivalités impériales non résolues (Maroc, Balkans, course navale).
\end{itemize}

\begin{aretenir}[Synthèse : Deuxième Révolution Industrielle (1850-1914)]
\begin{enumerate}
    \item \textbf{Énergies :} Électricité (réseaux, moteurs), Pétrole (moteur explosion, Diesel), Gaz. Charbon encore dominant mais relatif déclin.
    \item \textbf{Matériaux :} Acier (Bessemer, Martin, \textbf{Thomas} = minerais phosphoreux Lorraine/Ruhr), Aluminium (électrolyse), Chimie de synthèse (teintures, engrais Haber-Bosch, explosifs).
    \item \textbf{Industries :} Mécanique lourde (naval, rail), Automobile (Benz, Ford), Électrotechnique (Siemens, AEG, GE).
    \item \textbf{Organisation :} Taylorisme (OST, chronométrage, séparation conception/exécution), Fordisme (chaîne 1913, salaire 5 \$/jour), Concentration (Cartels Allemagne, Trusts USA, Holdings Europe).
    \item \textbf{Acteurs :} Allemagne (leader chimie/électricité/acier Thomas), USA (pétrole/auto/électricité/volume), France (luxe/auto/électricité/finance), Russie (décollage tardif capitaux français).
    \item \textbf{Crises :} Grande Dépression 1873-1896 (déflation par gains productivité) $\to$ Protectionnisme (Bismarck 1879, Méline 1892) $\to$ Impérialisme (débouchés, matières premières, capitaux).
\end{enumerate}
\end{aretenir}

\begin{exercice}[Entraînement : Document d'époque -- Extrait de Friedrich List, \emph{Le Système national d'économie politique} (1841), et statistiques Krupp 1913]
\textbf{Document 1 (List) :} « La puissance manufacturière est la source de la puissance commerciale, navale et militaire. [...] Un peuple qui ne possède pas l'industrie manufacturière ne peut pas conserver son indépendance politique. »
\textbf{Document 2 (Krupp 1913) :} Effectifs : 81 000 ouvriers. Production : 1,5 Mt acier, 300 000 obus/an, 50 cuirassés livrés. Filiales : mines (charbon, fer), aciéries, chantiers navals (Germania), usines d'armement, banque propre, chemins de fer privés, colonies agricoles (Allemagne de l'Est). Structure : Holding familiale + participation Deutsche Bank.

\textbf{Questions :}
\begin{enumerate}
    \item Situer List dans le débat libre-échange / protectionnisme en Allemagne (contexte 1841 vs 1879).
    \item À partir du Document 2, qualifier la forme de concentration de Krupp (cartel, trust, holding/Konzern) et justifier par deux éléments.
    \item Expliquer comment l'acier Thomas (Figure 2) a permis l'expansion de Krupp et la politique de flotte allemande (lois 1898, 1900).
    \item En quoi l'exemple Krupp illustre-t-il le lien entre Deuxième Révolution industrielle et impérialisme (section IV) ?
\end{enumerate}
\end{exercice}

\begin{corrige}[Exercice 1]
\begin{enumerate}
    \item \textbf{List (1841) :} Économiste allemand, théoricien du \textbf{protectionnisme éducateur} : l'industrie naissante a besoin de tarifs protecteurs pour rattraper le retard sur l'Angleterre. Précurseur du \textbf{Tarif Bismarck (1879)} (« fer et seigle »). En 1841, la Confédération germanique (Zollverein) commence l'unification douanière ; List plaide pour l'unification politique + protectionnisme.
    \item \textbf{Krupp = Holding / Konzern (familial + bancaire).} Justifications : (a) \textbf{Holding familiale} : la famille Krupp détient le contrôle via une société de tête (Fried. Krupp AG, 1903) qui possède les filiales ; (b) \textbf{Intégration verticale complète} : mines $\to$ aciéries $\to$ transformation (chantiers, armement) $\to$ transport $\to$ banque propre $\to$ agriculture (auto-suffisance vivrière). Pas de fusion avec concurrents (donc pas trust) ; pas simple entente sur prix (donc pas cartel). La Deutsche Bank siège au conseil (Aufsichtsrat) : financement + contrôle stratégique.
    \item \textbf{Acier Thomas + Flotte :} Le procédé Thomas (1878) rend exploitables les minerais phosphoreux de Ruhr (base de Krupp). Acier bon marché, homogène, en grandes quantités $\to$ plaques de blindage, canons, coques. Les \textbf{lois de flotte (Flottengesetze) 1898 et 1900} (Tirpitz) commandent 38 cuirassés, 20 grands croiseurs, 38 petits croiseurs en 20 ans. Krupp est le fournisseur quasi exclusif (blindages, artillerie). L'acier Thomas = condition technique de la \emph{Weltpolitik} navale.
    \item \textbf{Lien industrialisation / impérialisme :} (a) \textbf{Matières premières :} Krupp importe minerais non phosphoreux (Suède, Espagne, plus tard colonies) $\to$ besoin de contrôle des sources. (b) \textbf{Débouchés :} Surproduction aciérienne européenne $\to$ exportation rails, locomotives, ponts, matériel vers colonies/empire ottoman (Bagdadbahn financée par Deutsche Bank, acier Krupp). (c) \textbf{Capitaux :} Profits industriels réinvestis dans mines coloniales, chemins de fer, banques locales. (d) \textbf{Militaire :} Supériorité technique (canons, navires) = capacité de conquête et de maintien de l'ordre colonial. Krupp incarne la \textbf{complexe militaro-industriel} qui pousse à l'expansion (section IV).
\end{enumerate}
\end{corrige}

\coursec{IV. Les fondements de l'impérialisme européen : pourquoi l'Afrique ? (1870-1914)}

La \emph{Deuxième Révolution Industrielle} (acier, électricité, pétrole, chimie) que nous venons d'étudier a métamorphosé l'économie européenne : elle a décuplé la productivité, concentré les capitaux dans de gigantesques trusts et créé une interdépendance planétaire. Mais cette puissance industrielle a un revers : elle engendre des \emph{surcapacités} de production et une soif insatiable de matières premières nouvelles (caoutchouc, cuivre, pétrole, coton). L'Europe industrielle se heurte à ses propres frontières. La solution trouvée par les puissances ? L'expansion coloniale. L'Afrique, dernière terre « vierge » sur les cartes européennes, devient le terrain de jeu de cette rivalité exacerbée. Comprendre l'impérialisme, c'est donc articuler les \textbf{besoins économiques nouveaux}, les \textbf{rivalités politiques européennes} et les \textbf{justifications idéologiques} qui rendent l'entreprise socialement acceptable.

\courssub{Les causes profondes : un impérialisme à multiples moteurs}

L'historiographie distingue classiquement trois registres explicatifs qui s'entrelacent : l'économique, le politique et l'idéologique. Aucun ne suffit seul ; leur conjonction déclenche le « partage de l'Afrique ».

\begin{definition}[Impérialisme]
Politique d'expansion d'une puissance étendant sa domination (politique, économique, culturelle) sur des territoires extérieurs à ses frontières métropolitaines. Au sens strict (1870-1914), il se distingue du colonialisme ancien (comptoirs, traite) par la volonté d'\textbf{occupation effective}, de \textbf{mise en valeur économique intégrée} à la métropole et de \textbf{contrôle administratif direct}. Le \emph{néocolonialisme} désigne, après les indépendances (années 1960), la pérennisation de cette domination par des moyens indirects (économiques, monétaires, culturels, militaires) sans administration coloniale formelle.
\end{definition}

\paragraph{Causes économiques : le moteur structurel.}
Trois analyses classiques éclairent ce moteur :
\begin{itemize}
    \item \textbf{John A. Hobson (1902), \emph{L'Impérialisme} :} La suraccumulation de capitaux en Europe (épargne excédentaire, profits des trusts) ne trouve plus de placements rentables sur le marché intérieur saturé. L'exportation de capitaux vers des territoires « protégés » (chemins de fer, mines, plantations) assure des taux de profit supérieurs.
    \item \textbf{Vladimir Lénine (1916), \emph{L'Impérialisme, stade suprême du capitalisme} :} L'impérialisme est une nécessité vitale pour le capitalisme monopoliste. Cinq traits : concentration de la production (monopoles), fusion capital bancaire/industriel (capital financier), exportation de \emph{capitaux} (non plus seulement de marchandises), partage du monde entre trusts, partage \emph{territorial} du monde entre grandes puissances. La guerre de 1914 en est l'aboutissement logique.
    \item \textbf{Joseph Schumpeter (1919) :} Critique l'explication purement économique. L'impérialisme est un \emph{atavisme} : survie d'une classe guerrière aristocratique et d'une machine d'État expansionniste dans une société capitaliste par nature pacifique (libre-échange).
\end{itemize}

\begin{aretenir}[Synthèse des moteurs économiques (1870-1914)]
\begin{enumerate}
    \item \textbf{Recherche de débouchés (marchés captifs) :} Surproduction industrielle (textile, acier, ferroviaire) $\rightarrow$ besoin de marchés protégés par le tarif douanier colonial (ex: Empire colonial français « préférentiel »).
    \item \textbf{Recherche de matières premières stratégiques :} Caoutchouc (pneumatiques, isolation électrique), Huile de palme (savonnerie, graissage industriel), Coton (textile), Cuivre/Étain/Or (électricité, armement), Pétrole (marine, automobile naissante).
    \item \textbf{Placement des capitaux excédentaires :} Exportation de capitaux (prêts d'État, actions de compagnies coloniales, infrastructures) $\rightarrow$ rentes garanties par la puissance publique métropolitaine.
\end{enumerate}
\end{aretenir}

\paragraph{Causes politiques et stratégiques : l'échiquier européen.}
L'économie fournit le \emph{moteur}, la politique européenne fournit le \emph{volant}.
\begin{itemize}
    \item \textbf{La revanche française (1870) :} Perte de l'Alsace-Lorraine $\rightarrow$ compensation coloniale (Tunisie 1881, Congo, Afrique Occidentale) pour restaurer le prestige et l'encerclement de l'Allemagne.
    \item \textbf{La \emph{Weltpolitik} allemande (Guillaume II, 1897) :} « Place au soleil » (\emph{Platz an der Sonne}). Flotte de guerre (lois navales Tirpitz) $\rightarrow$ besoin de bases charbonnières et de stations radio (Douala, Dar es Salam, Tsingtao).
    \item \textbf{L'hégémonie britannique menacée :} Contrôle de la route des Indes (Suez 1875, Égypte 1882), Cape-to-Cairo (Rhodes) vs rêve français Dakar-Djibouti. Crise de Fachoda (1898) : apogée de la rivalité franco-britannique, résolue par l'Entente cordiale (1904).
    \item \textbf{L'équilibre des puissances :} La conférence d'Algésiras (1906) et les crises marocaines (1905, 1911) montrent que les colonies sont des \emph{monnaies d'échange} diplomatiques en Europe.
\end{itemize}

\begin{attention}[Erreur fréquente : Réduire l'impérialisme à l'économie seule]
Ne pas écrire : « L'impérialisme est \emph{uniquement} la recherche de profits capitalistes. » \textbf{Nuance exigée au Bac :} La rivalité franco-allemande en Europe (Alsace-Lorraine) explique en grande partie la course aux colonies en Afrique (compensation, encerclement, prestige). L'Allemagne de Bismarck, réticente au colonialisme (1884 : « Ma carte de l'Afrique est en Europe »), cède à la pression des lobbies industriels \emph{et} à la nécessité de ne pas laisser la France et la Grande-Bretagne s'accaparer le continent sans contrepartie. L'impérialisme est \textbf{surdéterminé} : économique \emph{et} géopolitique.
\end{attention}

\paragraph{Causes idéologiques et culturelles : la « mission civilisatrice ».}
L'opinion publique métropolitaine doit être convaincue. Trois piliers :
\begin{itemize}
    \item \textbf{Le darwinisme social :} Application abusive de la sélection naturelle aux sociétés humaines. « Les races supérieures » (Européens) ont le devoir biologique de dominer « les races inférieures » (Africains, Asiatiques) pour le progrès de l'humanité. Herbert Spencer, Gustave Le Bon.
    \item \textbf{Le « Fardeau de l'homme blanc » (Rudyard Kipling, 1899) :} Poème adressé aux USA pour les Philippines, devient l'hymne de l'impérialisme : charge morale de civiliser les « peuples à moitié diables, à moitié enfants ».
    \item \textbf{Le discours républicain français (Jules Ferry, 1885) :} « Les races supérieures ont un droit vis-à-vis des races inférieures... le devoir de les civiliser. » Laïcité, école, médecine, abolition de l'esclavage (traite arabo-swahilie en Afrique de l'Est) servent de paravent humanitaire à la conquête militaire.
    \item \textbf{L'exploration géographique :} Livingstone (missionnaire, anti-esclavagiste, sources du Nil), Stanley (journaliste, agent de Léopold II, Congo), Brazza (explorateur français, Gabon-Congo). Ils cartographient, signent des traités, préparent l'occupation effective.
\end{itemize}

\begin{savaistu}[L'exploration, avant-garde de la conquête]
Pierre Savorgnan de Brazza (1852-1905), naturalisé français, explore l'Ogooué et le Congo (1875-1882). Contrairement à Stanley (au service de Léopold II de Belgique), il signe des traités d'amitié et de protection avec le roi Makoko (Téké) en 1880, fondant \textbf{Brazzaville}. Son humanisme relatif (refus de la force, respect des coutumes) ne l'empêche pas d'être l'agent de l'expansion française. Il meurt épuisé en 1905 en enquêtant sur les abus de la concessionnaire \emph{Compagnie du Haut-Ogooué} : symbole des contradictions de la « mission civilisatrice ».
\end{savaistu}

\courssub{L'accélération brutale : de la conférence de Berlin au partage effectif (1880-1885)}

Jusqu'en 1880, la présence européenne est côtière (comptoirs, traite). En 1885, l'intérieur est partagé sur le papier. Trois événements déclenchent la ruée :

\begin{exemplebox}[Les trois détonateurs de la ruée (1882-1884)]
\begin{enumerate}
    \item \textbf{Occupation britannique de l'Égypte (1882) :} Protection du canal de Suez (vérité stratégique vitale pour l'Inde). Provoque la méfiance française (exclue du condominium) $\rightarrow$ compensation en Afrique de l'Ouest (Tunisie déjà prise 1881, poussée vers Niger/Tchad).
    \item \textbf{La « Crise du Congo » (1879-1884) :} Stanley (Léopold II) vs Brazza (France) vs Portugal (revendications anciennes). Léopold II crée l'\emph{Association Internationale du Congo} (AIC), façade philanthropique pour s'approprier le bassin du fleuve. La France et le Portugal protestent. Bismarck (Allemagne) propose une conférence internationale pour régler les différends et garantir le libre-échange.
    \item \textbf{Revendications allemandes en Afrique du Sud-Ouest, Togo, Cameroun, Afrique de l'Est (1884) :} Gustav Nachtigal, commissaire impérial, signe des traités de « protection » (Douala, 12 juillet 1884). L'Allemagne entre dans le club colonial, forçant les autres à officialiser leurs possessions.
\end{enumerate}
\end{exemplebox}

\begin{definition}[Conférence de Berlin (15 nov. 1884 -- 26 fév. 1885) -- Acte général]
Réunissant 14 puissances (dont USA, Ottoman, mais \textbf{aucun Africain}), elle pose les \textbf{règles du jeu} du partage :
\begin{enumerate}
    \item \textbf{Libre-échange :} Bassin du Congo (État Indépendant du Congo de Léopold II) et bassin du Niger : liberté de commerce pour tous les signataires (pas de monopole commercial).
    \item \textbf{Liberté de navigation :} Fleuves Congo, Niger, Zambèze ouverts au trafic de toutes les nations.
    \item \textbf{Principale innovation juridique : Notification et Occupation effective (Art. 34-35).} Toute nouvelle prise de possession sur la côte africaine doit être notifiée aux autres puissances. \textbf{Surtout :} Elle n'est valide que si elle est suivie d'une \textbf{occupation effective} (administration, troupes, traités avec chefs locaux). Fin des revendications purement « papier » sur l'intérieur.
    \item \textbf{Abolition de la traite des Noirs :} Engagement moral (peu respecté sur le terrain par les concessionnaires privés, ex: Congo belge).
\end{enumerate}
\textbf{Conséquence majeure :} Le partage devient une course de vitesse : il faut signer des traités, planter des drapeaux, installer des postes administratifs \emph{avant} les rivaux. L'Afrique est partagée \emph{sur la carte de Berlin} (parallèles, méridiens, fleuves) sans connaître les réalités ethniques.
\end{definition}

\begin{popfigure}
\begin{tikzpicture}[scale=0.65]
    % Simplified Africa outline coordinates (approximate)
    \def\africa{ (0,0) .. controls (2,-1) and (5,-2) .. (7,-1) .. controls (9,0) and (10,2) .. (9,5) .. controls (8,7) and (6,8) .. (4,9) .. controls (2,8) and (0,7) .. (-1,5) .. controls (-2,3) and (-2,1) .. (0,0) }
    % 1880 Map (Left)
    \begin{scope}[xshift=-10cm]
        \draw[thick, popdark] \africa;
        \node[align=center, font=\bfseries\large, popdark] at (-3, 10) {Afrique vers 1880};
        \node[font=\small, align=center] at (-3, -2) {Présence côtière seulement.\\Intérieur : États africains souverains.};
        % Coastal dots
        \foreach \pos/\name in {{(-1,1)}/Saint-Louis, {(1,-1)}/Gorée, {(3,0)}/Lagos, {(5,1)}/Douala, {(6,3)}/Libreville, {(7,5)}/Pointe-Noire, {(6,7)}/Luanda, {(3,8)}/Le Cap, {(0,6)}/Mombasa, {(-1,4)}/Zanzibar, {(-2,2)}/Alger, {(-3,0)}/Tunis} {
            \fill[popblue] \pos circle (2pt);
            \node[font=\tiny, anchor=west, xshift=3pt, popdark] at \pos {\name};
        }
        \draw[->, very thick, poporange] (-6, 4) -- (-8, 4) node[left, font=\small, popdark] {Traite / Échanges côtiers};
    \end{scope}
    % 1914 Map (Right)
    \begin{scope}
        \draw[thick, popdark] \africa;
        \node[align=center, font=\bfseries\large, popdark] at (7, 10) {Afrique en 1914};
        % Colonial blocks (simplified polygons with color)
        % France (Blue) - West, North, Equatorial, Madagascar
        \fill[popblue!30] (-1,1) -- (2,-1) -- (5,0) -- (5,3) -- (3,5) -- (0,6) -- (-2,4) -- (-3,1) -- cycle;
        % UK (Pink/Red) - Egypt, Sudan, East, South, Nigeria, Gold Coast, Sierra Leone, Gambia
        \fill[poppink!40] (-3,1) -- (-3,-1) -- (0,-1) -- (1,0) -- (1,3) -- (3,5) -- (4,7) -- (2,8) -- (0,7) -- (-1,5) -- cycle;
        % Germany (Grey) - Togo, Cameroon, SW Africa, East Africa
        \fill[popmuted!50] (1,0) -- (3,1) -- (4,3) -- (3,4) -- (2,3) -- cycle; % Cameroon approx
        \fill[popmuted!50] (5,5) -- (7,5) -- (7,7) -- (5,7) -- cycle; % East Africa approx
        \fill[popmuted!50] (4,-1) -- (6,-1) -- (6,1) -- (4,1) -- cycle; % SW Africa approx
        % Belgium (Green) - Congo
        \fill[popgreen!30] (2,3) -- (5,3) -- (5,6) -- (2,6) -- cycle;
        % Portugal (Purple) - Angola, Mozambique
        \fill[poppurple!30] (5,5) -- (7,5) -- (7,2) -- (5,2) -- cycle; % Mozambique
        \fill[poppurple!30] (4,1) -- (6,1) -- (6,4) -- (4,4) -- cycle; % Angola
        % Italy (Orange) - Libya, Eritrea, Somalia
        \fill[poporange!30] (-2,5) -- (0,5) -- (0,3) -- (-2,3) -- cycle; % Libya
        % Spain (Yellow) - Rio de Oro, Morocco N, Guinea
        \fill[popgold!40] (-3,2) -- (-2,2) -- (-2,1) -- (-3,1) -- cycle;
        % Legend
        \begin{scope}[yshift=-11cm, xshift=-2cm]
            \foreach \col/\lab/\i in {popblue!30/France/0, poppink!40/Royaume-Uni/1, popmuted!50/Allemagne/2, popgreen!30/Belgique (État indép. Congo)/3, poppurple!30/Portugal/4, poporange!30/Italie/5, popgold!40/Espagne/6} {
                \fill[\col] (0,{\i*0.7}) rectangle (0.5,{\i*0.7+0.5});
                \node[anchor=west, font=\small, popdark] at (0.7,{\i*0.7+0.25}) {\lab};
            }
        \end{scope}
        % Flux arrows
        \draw[->, very thick, popgreen!70!popdark] (8, 2) -- (11, 2) node[right, font=\small, popdark, align=left] {Matières premières\\(Caoutchouc, Cuivre,\\Cacao, Coton, Huile palme)};
        \draw[->, very thick, popblue!70!popdark] (11, -1) -- (8, -1) node[left, font=\small, popdark, align=right] {Biens manufacturés\\(Textile, Armes, Outils,\\Métallurgie, Alcool)};
        \draw[->, very thick, poporange] (8, 5) -- (11, 5) node[right, font=\small, popdark] {Capitaux (Investissements)};
    \end{scope}
\end{tikzpicture}
\legende{Figure 1 : Le partage de l'Afrique. À gauche, la situation vers 1880 : présence européenne cantonnée aux côtes. À droite, le partage achevé en 1914 : l'intérieur est découpé en colonies administratives aux frontières géométriques. Les flèches illustrent l'intégration forcée de l'Afrique dans l'économie mondiale industrialisée : extraction de matières premières vers l'Europe, écoulement de biens manufacturés et de capitaux vers l'Afrique.}
\end{popfigure}

\courssub{Le cas camerounais : du traité de Douala à la mise en valeur coloniale}

Le Cameroun (Kamerun) illustre parfaitement la mécanique impériale : traité inégal, résistance/collaboration, mise en valeur extractive.

\begin{methode}[Analyser un traité inégal : le Traité de Douala (12 juillet 1884)]
\textbf{Objectif :} Comprendre comment le droit international européen est utilisé pour légaliser la spoliation.
\begin{enumerate}
    \item \textbf{Identifier les parties :} D'un côté, \textbf{Gustav Nachtigal}, commissaire impérial allemand, représentant le Reich (puissance industrielle, militaire, juridique). De l'autre, les \textbf{chefs Douala} (King Bell, King Akwa, King Deido), autorités locales souveraines sur le plan coutumier, mais en position de faiblesse militaire et diplomatique.
    \item \textbf{Analyser la traduction et la médiation :} Le traité est rédigé en \textbf{allemand}. La version en \textbf{douala} (langue locale) est une traduction orale faite par des interprètes (souvent des missionnaires ou des commerçants africains acculturés). Les concepts juridiques européens (\emph{Schutzgebiet} = protectorat, \emph{Souveränität} = souveraineté, \emph{Land} = terre) n'ont pas d'équivalents exacts en douala (notion de \emph{maison}, de \emph{peuple}, d'\emph{usage} de la terre, pas de propriété foncière absolue).
    \item \textbf{Décortiquer les clauses clés :}
        \begin{itemize}
            \item \emph{Clause de protection :} L'Allemagne étend sa « protection » sur le territoire. Pour les chefs : alliance militaire contre rivaux locaux (Bakweri, Bakossi) ou concurrents européens (Britanniques). Pour l'Allemagne : base juridique de la souveraineté.
            \item \emph{Clause foncière :} « Les terres inoccupées appartiennent à l'Empire. » Notion européenne de \emph{terra nullius} (terre sans maître) ignorant le droit coutumier de jachère, de forêt sacrée, de parcours.
            \item \emph{Clause commerciale :} Libre-échange, monopole de fait pour les firmes allemandes (Woermann, Jantzen \& Thormählen).
        \end{itemize}
    \item \textbf{Observer la signature :} Les chefs apposent leur \textbf{croix} (signe distinctif personnel, marque d'oralité). Nachtigal appose le \textbf{sceau impérial} (symbole de l'État de droit écrit). Asymétrie scripturaire majeure.
    \item \textbf{Conclure :} Le traité transforme une relation commerciale inégale en une relation politique de suzeraineté. C'est l'acte de naissance juridique du \emph{Kamerun}.
\end{enumerate}
\end{methode}

\paragraph{Résistance et collaboration : une société en tension.}
L'installation effective (postes militaires, administration, justice indigène, impôt de capitation 1896) provoque des réactions variées :
\begin{itemize}
    \item \textbf{Résistance armée :} Peuple Bakweri (mont Cameroun, 1891-1894), Bamoun (Njoya, d'abord résistant puis collaborateur avisé), Bassa, Bamiléké. Répression brutale (politique de la « terre brûlée », déportations).
    \item \textbf{Résistance politique et juridique :} \textbf{Rudolf Duala Manga Bell} (roi Bell, éduqué en Allemagne, avocat). Utilise le droit allemand (pétitions au Reichstag, presse) contre l'expropriation des terres Douala (projet de ségrégation urbaine 1910). Exécuté pour haute trahison (8 août 1914) avec son allié \textbf{Martin Paul Samba} (chef Bulu, ancien élève de l'école militaire de Berlin, organisateur de la résistance armée dans le Sud). \emph{Exemple type de l'élite « évoluée » retournée contre le système colonial.}
    \item \textbf{Collaboration :} Chefs intermédiaires (warrant chiefs), interprètes, commis, porteurs, soldats (askaris). Nécessaires à l'administration indirecte (Indirect Rule à l'allemande : \emph{Bezirkshauptmann} + chefs coutumiers instrumentalisés).
\end{itemize}

\begin{exemplebox}[Mise en valeur du Kamerun : l'économie de plantation et de rail (1884-1914)]
L'Allemagne ne cherche peuplement (peu de colons : ~1 500 en 1913) mais \textbf{rentabilité immédiate}.
\begin{itemize}
    \item \textbf{Infrastructures :} 3 lignes de chemin de fer (Nord : Douala-Nkongsamba 1911, extension Ngaoundéré ; Centre : Douala-Yaoundé 1912 ; Sud : Douala-Buea). Port en eau profonde de Douala (1909). Réseau télégraphique.
    \item \textbf{Plantations (Sud-Ouest, côte volcanique fertile) :} Cacao (variété Trinitario), Café, Caoutchouc, Palmiers à huile, Bananes. Main-d'œuvre : recrutement forcé (corvée), puis travailleurs sous contrat (engagés) du Nord (Adamawa) et du Nigeria britannique.
    \item \textbf{Chiffres clés 1913 (Exemple chiffré) :}
\end{itemize}
\end{exemplebox}

\begin{exemplebox}[Exemple chiffré : Le Kamerun, grenier tropical de l'Allemagne (1913)]
\begin{itemize}
    \item \textbf{60 \%} de l'huile de palme consommée en Allemagne provient du Kamerun.
    \item \textbf{100 \%} du cacao allemand (colonial) vient du Kamerun.
    \item \textbf{Investissements allemands :} $\sim$ \textbf{400 millions de Marks-or} (chemins de fer, ports, administration, plantations).
    \item \textbf{Commerce extérieur du Kamerun (1913) :} Exportations $\approx$ 60 millions Marks (cacao, caoutchouc, bois, ivoire) ; Importations $\approx$ 40 millions Marks (textiles, alcool, quincaillerie, armes). Excédent commercial massif pour la métropole.
    \item \textbf{Main-d'œuvre :} $\sim$ 25 000 travailleurs engagés sur les plantations (conditions très dures, mortalité élevée).
\end{itemize}
\end{exemplebox}

\begin{popfigure}
\begin{tikzpicture}[scale=0.7]
    % Simplified Kamerun outline (approx 1884-1911 borders + Neukamerun)
    % Coast line
    \draw[thick, popdark] (0,0) .. controls (1,1) and (3,2) .. (5,2) .. controls (6,1.5) and (7,0) .. (7,-1) .. controls (6,-2) and (4,-2) .. (2,-1) .. controls (1,-0.5) and (0,0) .. (0,0);
    % Neukamerun (1911 extension East/South) - dashed
    \draw[thick, dashed, popmuted] (5,2) .. controls (7,3) and (9,2) .. (9,0) .. controls (9,-1) and (7,-2) .. (5,-1) .. controls (4,-1.5) and (3,-1) .. (2,-1);
    % Rivers
    \draw[popblue, thick] (3.5,0.5) .. controls (3,0) .. (2.5,-0.5) node[below, font=\tiny, popdark] {Wouri};
    \draw[popblue, thick] (4,1) .. controls (5,1.5) .. (6,2) node[right, font=\tiny, popdark] {Sanaga};
    \draw[popblue, thick] (4,-0.5) .. controls (5,-1) .. (6,-1.5) node[right, font=\tiny, popdark] {Nyong};
    % Railways
    \draw[very thick, poporange] (2,0) -- (3.5,0) node[midway, above, font=\tiny, popdark] {Nord} -- (5,1);
    \draw[very thick, poporange] (2,0) -- (3.5,-0.5) node[midway, below, font=\tiny, popdark] {Centre} -- (5,-1);
    \draw[very thick, poporange] (2,0) -- (3,0.5) node[midway, above, font=\tiny, popdark] {Sud} -- (4,1.5);
    % Cities / Sites
    \fill[popred] (2,0) circle (3pt); \node[above, font=\small\bfseries, popdark] at (2,0.4) {Douala (Port, Capitale)};
    \fill[popred] (4,1.5) circle (3pt); \node[above, font=\small, popdark] at (4,1.9) {Buea (Gouv., Climat)};
    \fill[popred] (3.5,0.5) circle (3pt); \node[left, font=\small, popdark] at (3,0.5) {Nkongsamba};
    \fill[popred] (5,1) circle (3pt); \node[right, font=\small, popdark] at (5.5,1) {Ngaoundéré (Nord)};
    \fill[popred] (5,-1) circle (3pt); \node[right, font=\small, popdark] at (5.5,-1) {Yaoundé (Centre)};
    \fill[popgreen] (3,1.8) circle (3pt); \node[above, font=\small, popdark] at (3,2.2) {Dschang (Climat, Élevage)};
    % Plantation zones (South West)
    \fill[popgreen!20, pattern=north east lines, pattern color=popgreen!50] (1.5,0.5) .. controls (2,1) and (3,1) .. (3.5,0.5) .. controls (3,0) and (2,0) .. (1.5,-0.2) .. controls (1,0) .. cycle;
    \node[font=\small, popgreen!70!popdark, align=center] at (2.5,0.7) {Zone de Plantations\\(Cacao, Café, Palmier)};
    % Legend
    \begin{scope}[xshift=10cm, yshift=-2cm]
        \draw[thick, popdark] (0,0) rectangle (0.3,0.3); \node[anchor=west, font=\small] at (0.5,0.15) {Ville / Poste admin};
        \draw[very thick, poporange] (0, -0.5) -- (0.3, -0.5); \node[anchor=west, font=\small] at (0.5,-0.5) {Chemins de fer};
        \draw[thick, dashed, popmuted] (0, -1) -- (0.3, -1); \node[anchor=west, font=\small] at (0.5,-1) {Frontières 1911 (Neukamerun)};
        \fill[popgreen!20, pattern=north east lines, pattern color=popgreen!50] (0,-1.5) rectangle (0.3,-1.2); \node[anchor=west, font=\small] at (0.5,-1.35) {Plantations};
        \draw[popblue, thick] (0, -2) -- (0.3, -2); \node[anchor=west, font=\small] at (0.5,-2) {Fleuves navigables};
    \end{scope}
    % North Arrow
    \node[font=\small\bfseries, popdark] at (0.5, 2.5) {N};
    \draw[->, thick, popdark] (0.5, 2.2) -- (0.5, 2.4);
    % Scale bar (approx)
    \draw[thick, popdark] (0.5, -2.5) -- (2.5, -2.5) node[midway, below, font=\tiny] {~200 km};
\end{tikzpicture}
\legende{Figure 2 : Croquis schématique du Kamerun allemand (1884-1916). Structure centre-périphérie : Douala (port, nœud ferroviaire), axes ferroviaires vers l'intérieur (Nord : bétail/coton ; Centre : administration/bois ; Sud : plantations), zone de plantations côtières (sol volcanique, climat), Dschang (altitude, climat tempéré, élevage). L'extension \emph{Neukamerun} (1911, traits pointillés) double la surface mais reste peu mise en valeur.}
\end{popfigure}

\courssub{Bilan et héritage : l'Europe, atelier du monde et gendarme colonial}

En 1914, l'Europe contrôle \textbf{90 \%} de la surface africaine (contre 10 \% en 1870). Le partage est achevé. Les frontières tracées à Berlin (parallèles, méridiens, fleuves) découpent les espaces ethniques (Bamiléké coupés entre Cameroun britannique et français, Fulbé entre Nigeria, Cameroun, Tchad, RCA, Douala entre Cameroun et Gabon), source de conflits post-indépendance.

\begin{popfigure}
\begin{tikzpicture}
    % Bar chart for French and British investments 1913 using pgfplots
    \begin{scope}[xshift=-7cm]
        \begin{axis}[
            ybar stacked,
            bar width=12pt,
            width=6cm, height=5cm,
            xlabel={Part des investissements (1913)},
            ylabel={\%},
            symbolic x coords={France, Royaume-Uni},
            xtick=data,
            nodes near coords={\pgfmathprintnumber\pgfplotspointmeta\%},
            nodes near coords align={vertical},
            legend style={at={(0.5,-0.3)}, anchor=north, legend columns=-1, font=\tiny},
            ymin=0, ymax=100,
            tick label style={font=\tiny},
            label style={font=\small},
        ]
        % France data (approx %): Russie 50, Balkans 20, AfrNord 15, AfrNoire 15
        \addplot[popblue] coordinates {(France,50) (Royaume-Uni,0)};
        \addplot[popblue!60!poppurple] coordinates {(France,20) (Royaume-Uni,0)};
        \addplot[poporange] coordinates {(France,15) (Royaume-Uni,0)};
        \addplot[popgreen] coordinates {(France,15) (Royaume-Uni,0)};
        % UK data (approx %): Dominions 35, AmLat 25, Inde 20, Afrique 10, Autres 10
        \addplot[poppink] coordinates {(France,0) (Royaume-Uni,35)};
        \addplot[poporange] coordinates {(France,0) (Royaume-Uni,25)};
        \addplot[popblue] coordinates {(France,0) (Royaume-Uni,20)};
        \addplot[popgreen] coordinates {(France,0) (Royaume-Uni,10)};
        \addplot[poppurple] coordinates {(France,0) (Royaume-Uni,10)};
        \legend{Russie/Alliés, Balkans, Afrique du Nord, Afrique Noire, Dominions, Am. Latine, Inde, Afrique, Autres}
        \end{axis}
        \node[align=center, font=\bfseries\small, popdark] at (0, -3.5) {Répartition des capitaux\ d'exportation (1913)};
        \node[align=center, font=\tiny, popmuted] at (0, -4.2) {France : $\sim$ 44 Md Frs-or | UK : $\sim$ 3,8 Md £ ($\sim$ 95 Md Frs)};
    \end{scope}
\end{tikzpicture}
\legende{Figure 3 : Destinations des capitaux d'exportation (1913). La France place massivement en Europe (Russie, alliée stratégique) et en Afrique du Nord (colonisation de peuplement/protectorats). L'Angleterre, plus libérale, diversifie vers les Dominions (peuplement britannique), l'Amérique latine (influence informelle) et l'Inde (joyau de la couronne). L'Afrique noire reste marginale en volume de capitaux pour les deux, mais stratégique pour les matières premières. Sources : \emph{Annuaire statistique de la France}, \emph{Statistical Abstract for the UK}.}
\end{popfigure}

\begin{attention}[Piège de rédaction : « L'Afrique a été colonisée pour ses matières premières. »]
Cette phrase est \textbf{trop courte} pour le Bac. Elle occulte :
\begin{enumerate}
    \item La \textbf{rivalité inter-européenne} (Alsace-Lorraine $\rightarrow$ Congo, Maroc).
    \item Le rôle des \textbf{moyens techniques} nés de la 2ème révolution industrielle (vapeur, quinine, fusil à répétition, télégraphe, chemin de fer) qui \emph{rendent possible} la conquête de l'intérieur.
    \item La \textbf{dimension idéologique} (opinion publique, presse, géographes, missionnaires) qui \emph{légitime} l'entreprise.
\end{enumerate}
\textbf{Au Bac :} Toujours articuler \textbf{Besoins (Pourquoi ?)}, \textbf{Moyens (Comment ?)}, \textbf{Acteurs/Rivalités (Qui ? Contre qui ?)}.
\end{attention}

\begin{exoresolu}[Sujet type Bac : « Montrez que la deuxième révolution industrielle est le moteur principal du partage de l'Afrique. »]
\textbf{Plan dialectique :}
\textbf{Introduction :} Accroche (partage 1885 vs 1914). Définition 2ème RI (acier, électricité, pétrole, chimie, concentration capitaliste). Définition partage Afrique (Conférence Berlin, occupation effective). Problématique : La 2ème RI est-elle la cause unique ou le condition nécessaire mais non suffisant ? Annonce du plan.
\textbf{I. Les besoins nouveaux créés par la 2ème RI (Moteur économique).}
\begin{itemize}
    \item Matières premières \emph{nouvelles} : Caoutchouc (pneumatiques Dunlop/Michelin, isolation câbles électriques), Cuivre/Étain (électrotechnique), Pétrole (marine de guerre, automobile), Coton (textile mécanisé).
    \item Excédents de capitaux (capital financier Hilferding/Lénine) $\rightarrow$ placement infrastructures coloniales (rails, ports) à rendement garanti.
    \item Surproduction industrielle $\rightarrow$ besoin de marchés captifs (tarifs préférentiels).
\end{itemize}
\textbf{II. Les moyens nouveaux fournis par la 2ème RI (Condition technique).}
\begin{itemize}
    \item Transport : Vapeur (navigation fluviale Congo/Niger), Chemin de fer (pénétration intérieure, évacuation produits), Télégraphe sans fil (liaison métropole-colonies).
    \item Arme : Fusil à répétition (Mauser, Lebel), Mitrailleuse (Maxim) $\rightarrow$ supériorité feu écrasante (Omdurman 1898).
    \item Médecine : Quinine (prophylaxie paludisme $\rightarrow$ survie Européens), Sérothérapies, Hygiène tropicale.
\end{itemize}
\textbf{III. Les rivalités politiques exacerbées par la puissance industrielle (Moteur géopolitique).}
\begin{itemize}
    \item Puissance industrielle = Puissance militaire (flotte cuirassée, artillerie). L'Allemagne unifiée (1871) devient 1ère puissance industrielle continentale $\rightarrow$ demande « place au soleil » (colonies = bases navales, prestige).
    \item France (2ème puissance) : Revanche 1870 + compensation coloniale.
    \item GB : Défense de l'hégémonie industrielle et navale menacée $\rightarrow$ contrôle routes (Suez, Cap).
    \item Crises diplomatiques (Tanger 1905, Agadir 1911) : Colonies comme enjeux/monnaies d'échange entre puissances industrielles rivales.
\end{itemize}
\textbf{Conclusion :} La 2ème RI fournit le \textbf{moteur} (besoins), le \textbf{véhicule} (moyens) et \textbf{l'essence} (rivalités de puissance). Mais le \textbf{volant} reste la décision politique des États (Bismarck, Ferry, Chamberlain, Léopold II) et l'idéologie (mission civilisatrice) qui transforment la potentialité technique en réalité coloniale. Partage = \emph{surdétermination} technique, économique, politique.
\tcblower
\textbf{Éléments de correction attendus :} Structure dialectique respectée. Vocabulaire précis (capital financier, occupation effective, quinine, mitrailleuse, Weltpolitik). Exemples datés et localisés (Omdurman, Berlin 1885, Agadir 1911, Kamerun). Nuance en conclusion (non-déterminisme technique pur).
\end{exoresolu}

\begin{savaistu}[La Conférence de Berlin : le mythe de la « salle du Reichstag »]
La conférence siégea dans le \textbf{Reichstag} (parlement allemand) à Berlin, dans la salle des séances du \emph{Bundesrat}. \textbf{Aucun Africain n'était invité.} Les 14 plénipotentiaires (Allemagne, Autriche-Hongrie, Belgique, Danemark, Espagne, États-Unis, France, Grande-Bretagne, Italie, Pays-Bas, Portugal, Russie, Suède-Norvège, Turquie) tracèrent les frontières à la règle sur de grandes cartes murales. Le roi des Belges Léopold II obtint l'État Indépendant du Congo (76 fois la Belgique) comme propriété \emph{privée}, pas comme colonie belge. L'acte de Berlin est le \textbf{certificat de naissance juridique de l'Afrique contemporaine} : ses frontières sont encore, à 90 \%, celles tracées en 1885 par des diplomates européens ignorant les réalités humaines du terrain.
\end{savaistu}

\begin{exercice}[Entraînement Bac - Document : Extrait du discours de Jules Ferry, Chambre des députés, 28 juillet 1885]
\textbf{Document :} « Messieurs, il faut le dire ouvertement : en effet, les races supérieures ont un droit vis-à-vis des races inférieures... Elles ont le devoir de les civiliser. »
\textbf{Questions :}
\begin{enumerate}
    \item Replacez cet extrait dans son contexte historique (date, auteur, enjeu parlementaire).
    \item Identifiez les deux arguments principaux développés par Jules Ferry pour justifier l'expansion coloniale.
    \item Montrez en quoi ce discours illustre l'articulation entre \emph{idéologie} et \emph{intérêts économiques} dans l'impérialisme français des années 1880.
\end{enumerate}
\end{exercice}

\begin{corrige}[Exercice n°1]
\begin{enumerate}
    \item \textbf{Contexte :} Juillet 1885. Jules Ferry (président du Conseil, ministre des Affaires étrangères) défend sa politique coloniale (prise de Tonkin, Tunisie, Congo, Madagascar) attaquée par la gauche (Clemenceau : « Ferry Tonkinois ») et la droite (coût, risque guerrier). Parlementarisme instable (IIIe République naissante).
    \item \textbf{Arguments :} 1) Argument \emph{humanitaire/civilisateur} : « Devoir de civiliser les races inférieures » (Darwinisme social, mission civilisatrice, abolition traite). 2) Argument \emph{économique/stratégique} (implicite dans le discours complet) : « Débouchés pour l'industrie », « Ports d'abri pour la marine », « Matières premières ».
    \item \textbf{Articulation :} L'idéologie (droit/devoir des races supérieures) fournit la \emph{légitimité morale} auprès de l'opinion et des députés réticents (républicains anti-militaristes). Elle masque et justifie les \emph{intérêts matériels} (lobbys industriels, banques, marine, armée) qui sont le \emph{moteur réel} de l'expansion. Le discours « civilisateur » rend l'impérialisme compatible avec les valeurs républicaines (Liberté, Égalité, Fraternité... pour les « civilisés »).
\end{enumerate}
\end{corrige}

\coursec{V. Synthèse : L'Europe, atelier du monde et gendarme colonial (Bilan 1914)}

\courssub{Transition : de la conquête à l'hégémonie planétaire}

Au terme de la section précédente, nous avons vu comment l'Europe, forte de sa supériorité technique, démographique et financière, a imposé son contrôle sur l'Afrique lors du « partage » (Conférence de Berlin, 1885). Mais l'impérialisme n'est pas une fin en soi : c'est le prolongement violent de la révolution industrielle. En 1914, l'Europe ne domine pas seulement l'Afrique ; elle structure le monde entier. Cette synthèse montre comment l'industrialisation a fait de l'Europe l'« atelier du monde » et le « gendarme colonial », tout en accumulant les contradictions qui la mèneront au suicide de 1914-1918.

\begin{aretenir}[Problématique de synthèse]
Comment l'industrialisation européenne (1780-1914) a-t-elle engendré une hégémonie mondiale à la fois économique, démographique, militaire et coloniale, porteuse en son sein des germes de sa propre destruction (1914) et d'héritages structurels durables pour les territoires colonisés comme le Cameroun ?
\end{aretenir}

\courssub{L'Europe en 1914 : le centre du monde par les chiffres}

L'année 1914 marque l'apogée de la puissance européenne. Quelques indicateurs suffisent à mesurer l'écart abyssal creusé entre l'Europe et le reste du monde en un siècle d'industrialisation.

\begin{exemplebox}[L'Europe en 1914 : quelques repères chiffrés clés]
\begin{itemize}
    \item \textbf{Population :} L'Europe (avec la Russie) compte environ 450 millions d'habitants, soit \textbf{1/4 de l'humanité}.
    \item \textbf{Production industrielle :} L'Europe (y compris la Russie) assure \textbf{90\% de la production industrielle mondiale} (acier, charbon, textile, chimie).
    \item \textbf{Commerce international :} Les puissances européennes contrôlent \textbf{2/3 des échanges mondiaux} (exportations + importations).
    \item \textbf{Domination territoriale :} \textbf{84\% de la surface des terres émergées} est sous domination européenne directe (colonies, protectorats) ou indirecte (influence économique, traités inégaux comme en Chine ou dans l'Empire ottoman).
    \item \textbf{Finance :} Londres, Paris, Berlin sont les places boursières mondiales ; l'Europe détient la quasi-totalité de l'or monétaire mondial et exporte des capitaux massifs (chemins de fer, emprunts d'État, mines).
\end{itemize}
\end{exemplebox}

\begin{popfigure}
\begin{tikzpicture}[scale=0.85, every node/.style={font=\small}]
% Radar chart axes
\def\angles{90,30,-30,-90,-150,150}
\def\labels{Population,PIB,Acier,Marine,Colonies,Armée}
% Grid (20% steps)
\foreach \r in {20,40,60,80,100} {
    \draw[popline,thin] plot[smooth cycle] coordinates {
        (90:\r) (30:\r) (-30:\r) (-90:\r) (-150:\r) (150:\r)
    };
    \node[popmuted,font=\tiny] at (90:\r+3) {\r\%};
}
% Axes lines and labels
\foreach \angle/\label [count=\i] in {90/Population,30/PIB,-30/Acier,-90/Marine,-150/Colonies,150/Armée} {
    \draw[popdark,thin] (0,0) -- (\angle:105);
    \node[font=\small,popdark] at (\angle:112) {\label};
}
% Data: Normalized Index (Max among the 6 = 100) for 1913/1914
% UK: Pop 27, PIB 43, Steel 33, Navy 100, Colonies 100, Army 18
% Germany: Pop 39, PIB 46, Steel 55, Navy 59, Colonies 8, Army 57
% France: Pop 23, PIB 28, Steel 14, Navy 18, Colonies 33, Army 57
% USA: Pop 58, PIB 100, Steel 100, Navy 41, Colonies 1, Army 7
% Russia: Pop 100, PIB 45, Steel 15, Navy 9, Colonies 5, Army 100
% Japan: Pop 31, PIB 14, Steel 1, Navy 14, Colonies 1, Army 14
% UK (popblue)
\draw[popblue, thick, fill=popblue!15] plot[smooth cycle, tension=0.2] coordinates {
    (90:27) (30:43) (-30:33) (-90:100) (-150:100) (150:18)
};
% Germany (poporange)
\draw[poporange, thick, fill=poporange!15] plot[smooth cycle, tension=0.2] coordinates {
    (90:39) (30:46) (-30:55) (-90:59) (-150:8) (150:57)
};
% France (poppurple)
\draw[poppurple, thick, fill=poppurple!15] plot[smooth cycle, tension=0.2] coordinates {
    (90:23) (30:28) (-30:14) (-90:18) (-150:33) (150:57)
};
% USA (popgreen)
\draw[popgreen, thick, fill=popgreen!15] plot[smooth cycle, tension=0.2] coordinates {
    (90:58) (30:100) (-30:100) (-90:41) (-150:1) (150:7)
};
% Russia (poppink)
\draw[poppink, thick, fill=poppink!15] plot[smooth cycle, tension=0.2] coordinates {
    (90:100) (30:45) (-30:15) (-90:9) (-150:5) (150:100)
};
% Japan (popgold)
\draw[popgold, thick, fill=popgold!15] plot[smooth cycle, tension=0.2] coordinates {
    (90:31) (30:14) (-30:1) (-90:14) (-150:1) (150:14)
};
% Legend
\begin{scope}[shift={(8,-6)}]
    \foreach \c/\n [count=\i] in {popblue/Royaume-Uni, poporange/Allemagne, poppurple/France, popgreen/USA, poppink/Russie, popgold/Japon} {
        \draw[\c, thick, fill=\c!20] (0,-\i*0.7) rectangle (0.5,-\i*0.7+0.4);
        \node[anchor=west, font=\small] at (0.7,-\i*0.7+0.2) {\n};
    }
\end{scope}
\end{tikzpicture}
\legende{Graphique radar (araignée) comparant les six grandes puissances en 1914 sur six indicateurs de puissance (indices normalisés : maximum observé parmi les six = 100). Sources : estimations historiques (Bairoch, Kennedy, Maddison).}
\end{popfigure}

\begin{attention}[Piège de lecture : la puissance globale vs la puissance par habitant]
Ne confondez pas \textbf{puissance totale} (agrégats nationaux) et \textbf{niveau de vie / productivité} (par habitant). En 1914, les USA dépassent déjà l'Europe en PIB total et production d'acier, mais leur population est plus faible. La Russie est une « géante aux pieds d'argile » : immense population et armée, mais PIB/habitant et industrialisation très faibles. Le graphique radar montre la \emph{capacité de projection de puissance} (guerre, colonisation), pas le bien-être des populations.
\end{attention}

\courssub{Le système « Centre-Périphérie » : l'interdépendance asymétrique}

L'hégémonie européenne repose sur une division internationale du travail (DIT) hiérarchisée, théorisée plus tard par l'école de la dépendance (A. G. Frank, I. Wallerstein). L'Europe industrialisée (le \textbf{Centre}) organise le monde en \textbf{Périphéries} spécialisées.

\begin{popfigure}
\begin{tikzpicture}[scale=1.1, every node/.style={font=\small}]
% Styles
\tikzset{
    centre/.style={circle, draw=popblue, thick, fill=popblue!20, minimum width=2.5cm, align=center},
    periph/.style={rectangle, rounded corners, draw=poporange, thick, fill=poporange!15, minimum width=2.2cm, minimum height=1.2cm, align=center},
    flux/.style={->, thick, >=Stealth},
    fluxmat/.style={flux, popred, line width=2pt},
    fluxman/.style={flux, popblue, line width=2pt},
    fluxcap/.style={flux, popgold, line width=2pt, dashed},
    label/.style={font=\tiny, fill=white, inner sep=1pt, midway, sloped}
}
% Nodes
\node[centre, align=center] (E) at (0,0){\textbf{CENTRE}\\\textbf{EUROPE}\\(Industrie, Finance,\\Technologie, Armée)};
\node[periph, align=center] (A) at (-4.5, 3){AFRIQUE\\Colonies\\ (Mat. 1ères, Main-d'œuvre)};
\node[periph, align=center] (B) at (4.5, 3){ASIE\\ (Indes, Indochine,\\ Chine, Japon)};
\node[periph, align=center] (C) at (-4.5, -3){AMÉRIQUE LATINE\\ (Produits tropicaux,\\ Minerais, Viande, Blé)};
\node[periph, align=center] (D) at (4.5, -3){DOMINIONS / USA\\ (Produits tempérés,\\ Investissements, Marchés)};
% Flux Matières premières (Rouge) : Periph -> Centre
\draw[fluxmat] (A) -- node[label, above, align=center]{Coton, Caoutchouc,\\Minéraux, Huiles,\\Main-d'œuvre forcée} (E);
\draw[fluxmat] (B) -- node[label, above, align=center]{Jute, Thé, Riz,\\Étain, Caoutchouc,\\Opium} (E);
\draw[fluxmat] (C) -- node[label, below, align=center]{Cuivre, Nitrates,\\Viande, Blé, Café,\\Sucre} (E);
\draw[fluxman] (D) -- node[label, below, align=center]{Blé, Laine, Viande,\\Métaux, Pétrole} (E);
% Flux Biens manufacturés (Bleu) : Centre -> Periph
\draw[fluxman] (E) -- node[label, below, align=center]{Textiles, Machines,\\Armes, Outils,\\Médicaments} (A);
\draw[fluxman] (E) -- node[label, below, align=center]{Textiles, Chemins de fer,\\Armes, Produits chimiques} (B);
\draw[fluxman] (E) -- node[label, above, align=center]{Textiles, Machines,\\Outils, Biens de luxe} (C);
\draw[fluxman] (E) -- node[label, above, align=center]{Capitaux, Machines,\\Technologies, Banques} (D);
% Flux Capitaux (Or) : Centre -> Periph (Infrastructures extractives)
\draw[fluxcap] (E) -- node[label, left, pos=0.3, align=center]{Capitaux (Emprunts,\\Actions, Concessions)\\→ Rails, Ports, Mines} (A);
\draw[fluxcap] (E) -- node[label, right, pos=0.3, align=center]{Capitaux →\\Infrastructures} (B);
\draw[fluxcap] (E) -- node[label, left, pos=0.7, align=center]{Capitaux →\\Rails, Ports, Banques} (C);
\draw[fluxcap] (E) -- node[label, right, pos=0.7, align=center]{Investissements\\massifs} (D);
% Legend
\begin{scope}[shift={(-7,-5)}]
    \draw[fluxmat] (0,0) -- (1,0) node[label, right, pos=0.5, sloped=false] {Matières premières / Main-d'œuvre};
    \draw[fluxman] (0,-0.6) -- (1,-0.6) node[label, right, pos=0.5, sloped=false] {Biens manufacturés / Services};
    \draw[fluxcap] (0,-1.2) -- (1,-1.2) node[label, right, pos=0.5, sloped=false] {Capitaux / Infrastructures extractives};
\end{scope}
\end{tikzpicture}
\legende{Schéma synthèse du système Centre-Périphérie vers 1914. L'Europe (Centre) importe matières premières et main-d'œuvre (flèches rouges), exporte biens manufacturés (flèches bleues) et exporte des capitaux pour construire les infrastructures d'extraction (rails, ports) plutôt que de développement intégré (flèches or pointillées).}
\end{popfigure}

\begin{definition}[Économie d'enclave]
Structure coloniale où l'infrastructure de transport (chemin de fer, port) relie \textbf{uniquement} la zone de production (plantation, mine, forêt) au port d'exportation, sans maillage du territoire ni intégration du marché intérieur local. L'objectif est l'extraction à moindre coût, non le développement du territoire colonisé.
\end{definition}

\begin{exoresolu}[Analyse d'un flux type : le coton égyptien vers Manchester (1913)]
\textbf{Énoncé :} En 1913, l'Égypte (protectorat britannique) exporte 1,2 million de tonnes de coton brut vers le Royaume-Uni. Manchester transforme ce coton en tissus. Une partie revient en Égypte (tissus bon marché), l'autre part vers l'Inde et l'Afrique. Analysez ce flux au prisme du modèle Centre-Périphérie.

\textbf{Solution détaillée :}
\begin{enumerate}
    \item \textbf{Spécialisation contrainte (Périphérie) :} L'Égypte est contrainte par l'administration britannique (Lord Cromer) à la monoculture du coton (irrigation barrages, crédit forcé). Elle perd sa souveraineté alimentaire (importation de blé).
    \item \textbf{Capture de la valeur ajoutée (Centre) :} Le filage, le tissage, la teinture, le commerce de gros, le transport maritime (flotte britannique), l'assurance (Lloyd's) et la finance (Banque d'Angleterre) restent en Europe. La valeur ajoutée (salaires, profits, impôts, innovation technique) enrichit le Centre.
    \item \textbf{Désindustrialisation de la Périphérie :} L'afflux de tissus industriels britanniques (moins chers, standardisés) ruine l'artisanat textile local (tisseurs coptes, villageois).
    \item \textbf{Infrastructure d'enclave :} Le réseau ferré égyptien (Delta - Le Caire - Alexandrie - Port-Saïd) dessert les zones cotonnières et le canal de Suez (artère impériale), pas l'intérieur du pays (Assouan, oasis).
    \item \textbf{Conclusion :} Relation d'échange inégal : la Périphérie vend du travail incorporé brut (valeur faible) et achète du travail incorporé qualifié (valeur forte). Le terme de l'échange se dégrade pour la Périphérie.
\end{enumerate}
\end{exoresolu}

\courssub{Les failles de l'hégémonie : les contradictions internes de 1914}

L'apogée européenne est instable. Quatre failles majeures traversent le système, rendant la guerre quasi inévitable.

\begin{aretenir}[Les quatre failles de l'Europe en 1914]
\begin{enumerate}
    \item \textbf{Rivalités impérialistes (Diplomatie de la canonnière) :} Le monde est « partagé » (Berlin 1885), mais les parts sont inégales. L'Allemagne (arrivée tardive) veut sa « place au soleil » (Crises marocaines 1905, 1911 : Tanger, Agadir). La France et le Royaume-Uni se rapprochent (Entente cordiale 1904). L'Empire ottoman se disloque (Guerres balkaniques 1912-1913), attisant les appétits autrichiens, russes, italiens.
    \item \textbf{Course aux armements (Technologie guerrière) :} Révolution industrielle appliquée à la guerre. \textbf{Dreadnought} (cuirassé tout-gros-calibres, turbine, 1906) rend la flotte mondiale obsolète → course navale GB/Allemagne. Artillerie à tir rapide, mitrailleuses, gaz, chemin de fer stratégique (plan Schlieffen). Budgets militaires explosent (5\% du PIB en France/Allemagne).
    \item \textbf{Question sociale et mouvements ouvriers :} La IIe Internationale (1889) unit les partis socialistes (SPD en Allemagne, SFIO en France, Labour au UK). Grèves massives (1910-1913). Peur des élites : « guerre sociale » ou « guerre extérieure » pour canaliser le mécontentement (théorie de la diversion patriotique).
    \item \textbf{Nationalismes non résolus :} \emph{Alsace-Lorraine} (France vs Allemagne), \emph{Balkans} (Serbie, Bulgarie, Grèce, Roumanie vs Autriche-Hongrie, Empire ottoman ; « poudrière »), \emph{Irlande} (Home Rule), \emph{Pologne} (partagée entre 3 empires), minorités dans l'Empire austro-hongrois (Tchèques, Slovaques, Hongrois, Roumains, Serbes, Croates, Slovènes).
\end{enumerate}
\end{aretenir}

\begin{attention}[Erreur fréquente : déterminisme de la guerre]
Ne pas écrire : « La guerre était inévitable à cause de l'impérialisme. »
\emph{Nuance :} L'impérialisme crée les \textbf{conditions structurelles} (rivalités, alliances, armement), mais le déclenchement d'août 1914 résulte de \textbf{choix politiques précis} (ultimatum autrichien à la Serbie, mobilisation russe, plan Schlieffen allemand, violation neutralité belge, entrée en guerre britannique). L'histoire n'est pas un destin mécanique.
\end{attention}

\courssub{Héritage pour le Cameroun : structures coloniales et contentieux durables}

Le Cameroun, colonie allemande (Kamerun, 1884-1916) puis mandat franco-britannique (SDN, 1919-1960), illustre parfaitement la transition de l'Europe « atelier/gendarme » vers l'ordre post-1918.

\begin{savaistu}[Le chemin de fer Douala-Yaoundé (Mittelkammerun) : l'enclave par excellence]
Commencé en 1907, achevé en 1911 par l'administration allemande. \textbf{Objectif unique :} évacuer le cacao, le caoutchouc, l'huile de palme de la forêt équatoriale vers le port de Douala.
\begin{itemize}
    \item \textbf{Coût humain :} Travail forcé (corvée, \emph{Zwangsarbeit}), recrutement violent dans le Nord (Adamawa) et le Centre. Estimations : \textbf{des milliers de morts} (maladies, accidents, épuisement, répression).
    \item \textbf{Structure persistante :} Ce rail (ligne Nord-Sud) structure encore aujourd'hui l'axe économique principal du Cameroun (Douala-Yaoundé-Ngaoundéré). Il n'a jamais été prolongé vers l'Ouest (Bamenda, Bafoussam) ni l'Extrême-Nord (Maroua) de manière efficace avant l'indépendance. C'est l'archétype de l'\textbf{économie d'enclave}.
\end{itemize}
\end{savaistu}

\begin{exemplebox}[Bilan commercial du Kamerun en 1913 (année record avant-guerre)]
\begin{center}
\begin{tabularx}{\linewidth}{|l|X|X|}\hline
\textbf{Flux} & \textbf{Produits principaux} & \textbf{Observations} \\ \hline
\textbf{Exportations} (vers Allemagne) & Cacao : 4 500 t \newline Caoutchouc : 1 200 t \newline Huile de palme / Amandes : 15 000 t \newline Bois, ivoire, peaux & Monocultures de rente imposées par l'administration et les compagnies (Woermann, Jantzen \& Thormählen). Main-d'œuvre villageoise contrainte (impôt de capitation payable en numéraire ou en travail). \\ \hline
\textbf{Importations} (depuis Allemagne) & Tissus (cotonnades), Alcool (rhum, schnaps), Quincaillerie (outils, armes), Produits manufacturés divers & Échange inégal : matières premières brutes vs biens finis. L'alcool sert de monnaie d'échange et de paiement de la main-d'œuvre (ravages sanitaires/sociaux). \\ \hline
\textbf{Bilan} & \textbf{Excédent commercial massif pour l'Allemagne} & Le Kamerun « s'autofinance » (budget colonial excédentaire) grâce aux excédents commerciaux et à l'impôt de capitation. L'Allemagne n'investit pas un mark métropolitain net ; la colonie paie son administration, son armée (Schutztruppe) et ses infrastructures. \\ \hline
\end{tabularx}
\end{center}
\legende{Données issues des rapports annuels du Gouvernement impérial du Kamerun (1913).}
\end{exemplebox}

\begin{definition}[Héritage bilingue et administratif]
La partition du Kamerun en 1919 (Traité de Versailles, mandats SDN Art. 22) entre France (4/5 Est) et Royaume-Uni (1/5 Ouest, adjoint au Nigeria) crée :
\begin{itemize}
    \item \textbf{Bilinguisme officiel} (Français/Anglais) et dualité juridique (Droit civil français / Common law britannique).
    \item \textbf{Administration centralisée} (modèle français jacobin) vs \textbf{Indirect Rule} (modèle britannique : chefs traditionnels comme auxiliaires).
    \item \textbf{Contentieux foncier} : Spoliation des terres des plantations (Sud-Ouest, Littoral) au profit de colons/entreprises ; lois foncières post-indépendance (1974) ne résolvent pas totalement les droits coutumiers.
\end{itemize}
\end{definition}

\courssub{Ouverture : 1914-1918, le suicide de l'Europe industrielle et l'accélération de l'histoire coloniale}

La Première Guerre mondiale n'est pas un accident : c'est l'aboutissement logique des contradictions exposées ci-dessus. Elle agit comme un \textbf{révélateur} et un \textbf{accélérateur}.

\begin{methode}[Rédiger une conclusion de synthèse (type Bac/Probatoire)]
Suivez ces 3 étapes pour conclure une dissertation ou une analyse de documents sur cette période :
\begin{enumerate}
    \item \textbf{Rappeler la problématique} : Reformulez la question centrale (ex: « Dans quelle mesure l'industrialisation a-t-elle fondé la domination européenne et ses limites ? »).
    \item \textbf{Résumer les 2-3 arguments majeurs} :
        \begin{itemize}
            \item Argument 1 : L'industrialisation (charbon/fer/vapeur puis acier/électricité/chimie) donne à l'Europe une supériorité productive, militaire, démographique \emph{quantitative} (1914 = apogée).
            \item Argument 2 : Cette puissance se traduit par un système mondial Centre-Périphérie (impérialisme, économie d'enclave) qui structure durablement les colonies (ex: Cameroun).
            \item Argument 3 : Mais le système porte ses contradictions : rivalités inter-impérialistes, course aux armements, nationalismes, tensions sociales.
        \end{itemize}
    \item \textbf{Ouvrir sur la conséquence immédiate ou durable} :
        \begin{itemize}
            \item \emph{Immédiate :} Août 1914 = suicide de l'Europe (guerre totale, industrielle, 10 millions de morts, ruine financière, hégémonie passée aux USA/Japon).
            \item \emph{Durable :} Traité de Versailles + Article 22 SDN (Mandats) → légitimation internationale de la colonisation mais principe du « bien-être et développement » des populations → germes du droit des peuples à disposer d'eux-mêmes → décolonisation post-1945. Au Cameroun : passage du Kamerun allemand aux mandats français/britanniques, héritage bilingue, frontières actuelles.
        \end{itemize}
\end{enumerate}
\end{methode}

\begin{exoresolu}[Sujet de dissertation type : « L'Europe, atelier du monde et gendarme colonial (1815-1914) ». Problématique et plan détaillé]
\textbf{Problématique :} Comment la double révolution industrielle (1780-1914) a-t-elle permis à l'Europe d'imposer son hégémonie planétaire (économique, coloniale, militaire) tout en engendrant les conflits qui feront basculer le monde en 1914 ?

\textbf{Plan en 3 parties (avec arguments chiffrés) :}

\textbf{I. L'Europe, atelier du monde : la puissance par l'industrie (1780-1914)}
\begin{itemize}
    \item A. Deux révolutions successives : vapeur/charbon/fer (UK, 1780-1850) → acier/électricité/pétrole/chimie (Allemagne/USA/France, 1850-1914). Gains de productivité exponentiels.
    \item B. Dominations chiffrées 1914 : 90\% production industrielle, 2/3 commerce, 1/4 population, 84\% territoire. Graphique radar (Fig. 1).
    \item C. Puissance financière : Exportation de capitaux (rails, mines, emprunts d'État), étalon-or, City de Londres = banque du monde.
\end{itemize}

\textbf{II. L'Europe, gendarme colonial : l'ordre impérialiste (1870-1914)}
\begin{itemize}
    \item A. Causes structurelles : Besoin de matières premières, débouchés, placement de capitaux excédentaires (Hobson, Lénine), démographie excédentaire, idéologies (racisme scientifique, mission civilisatrice).
    \item B. Mise en œuvre : Conférence de Berlin (1885), partage de l'Afrique, protectorats en Asie, traités inégaux (Chine, Ottoman). Cameroun : Kamerun allemand (1884-1916).
    \item C. Système Centre-Périphérie : Échange inégal, économie d'enclave (rail Douala-Yaoundé), désindustrialisation des périphéries, infrastructures extractives only (Fig. 2).
\end{itemize}

\textbf{III. L'apogée porteuse de rupture : les contradictions fatales (1914)}
\begin{itemize}
    \item A. Rivalités inter-impérialistes : Monde partagé mais parts inégales → Crises (Maroc, Balkans), alliances rigides (Triple Alliance vs Triple Entente), course aux armements (Dreadnought).
    \item B. Tensions internes : Question sociale (IIe Internationale, grèves), nationalismes opprimés (Balkans, Alsace-Lorraine, Irlande, minorités Autriche-Hongrie).
    \item C. 1914 : Basculement. Guerre industrielle totale (chars, avions, gaz, sous-marins, mobilisation totale). Suicide démographique/financier de l'Europe. Hégémonie transférée aux USA/Japon.
    \item D. Héritage colonial : Mandats SDN (Art. 22) → Cameroun franco-britannique → Indépendance (1960/1961) avec structures d'enclave, bilinguisme, contentieux foncier non résolus.
\end{itemize}

\textbf{Conclusion :} L'Europe de 1914 est le produit abouti de l'industrialisation : puissance sans précédent, mais fragilisée par ses propres succès (inégalités, rivalités, résistances). La Grande Guerre en est l'effondrement et le révélateur. Les structures coloniales mises en place (économie d'enclave, administration, frontières) survivent à l'Europe coloniale et pèsent encore sur le Cameroun contemporain.
\end{exoresolu}

\begin{aretenir}[Bilan final : 3 dates-clés pour ancrer la synthèse]
\begin{itemize}
    \item \textbf{1885} : Conférence de Berlin → Partage légal de l'Afrique (apogée diplomatique).
    \item \textbf{1913/1914} : Apogée statistique (Production, Commerce, Territoire, Armement) → « Belle Époque » pour les élites, poudrière pour les peuples.
    \item \textbf{1919} : Traité de Versailles + SDN (Art. 22) → Fin des empires centraux, mandats (Cameroun), début de la fin de l'hégémonie européenne (USA/Japon sortent renforcés, nationalismes coloniaux naissants).
\end{itemize}
\end{aretenir}

\newpage

\coursec{Activité d'intégration}

\coursec{La reconstruction/renaissance de l'Europe et les révolutions industrielles}

\courssub{Objectifs de l'activité}
Cette activité d'intégration vise à mobiliser l'ensemble des savoirs et savoir-faire de la leçon sur un cas d'étude concret et national : le Cameroun sous domination allemande (1884-1914). Elle permet à l'élève de :
\begin{itemize}
    \item \textbf{Recontextualiser} les révolutions industrielles européennes (besoins en matières premières, capitaux, débouchés) dans la logique impérialiste en Afrique.
    \item \textbf{Analyser} un corpus documentaire hétérogène (traité, statistiques, témoignage, carte) en croisant les sources.
    \item \textbf{Produire} un croquis commenté (compétence graphique) représentant l'espace colonial utile (\emph{Kamerun utile}).
    \item \textbf{Rédiger} une réponse argumentée, nuancée et structurée (20 lignes), distinguant court terme (contraintes, violence) et long terme (héritage infrastructurel), conformément aux attendus du Probatoire.
\end{itemize}

\courssub{Situation : Dossier — Le Cameroun face à la révolution industrielle allemande (1884-1914) : Exploitation ou Modernisation ?}
\medskip
\textbf{Contexte historique.} Au sortir de la Conférence de Berlin (1884-1885), l'Allemagne de Bismarck, puissance industrielle montante (\cle{Deuxième Révolution Industrielle} : acier, chimie, électricité), acquiert le \emph{Kamerun}. Le territoire devient une \cle{colonie d'exploitation} intégrée à l'économie mondiale. L'administration coloniale, appuyée sur des compagnies concessionnaires (Woermann, Jantzen \& Thormählen), met en place un système extractiviste. La question centrale est de savoir si l'industrialisation allemande a constitué un levier de développement pour les sociétés camerounaises ou un mécanisme de prédation violente.

\medskip
\textbf{Document 1 : Extrait du Traité de Douala (12 juillet 1884), signé par les rois Bell et Akwa et le Dr Nachtigal.}
\og Les rois et chefs du pays de Douala cèdent à Sa Majesté l'Empereur d'Allemagne tous leurs droits souverains sur le territoire... en échange de la protection allemande. Les Allemands s'engagent à respecter les droits de propriété et les coutumes locales. Le commerce reste libre pour toutes les nations, mais l'administration allemande perçoit les droits de douane. \fg
\emph{Source : Archives nationales du Cameroun, Série Traité de protectorat.}

\medskip
\textbf{Document 2 : Infrastructures de transport — Le chemin de fer Nord (1901-1911) et Sud (1906-1914).}
\begin{center}
\begin{tabularx}{\linewidth}{|l|X|X|}\hline
\rowcolor{popblueL}
\textbf{Ligne} & \textbf{Caractéristiques techniques} & \textbf{Flux principaux (année 1913)} \\ \hline
Nord (Douala -- Nkongsamba, 160 km) & Voie métrique, 12 tunnels, 40 ponts. Coût : 45 millions de Marks-or. Main-d'œuvre forcée (corvée) massive. & Café, cacao, huile de palme vers Douala. Matériel ferroviaire, textiles, alcools vers l'intérieur. \\ \hline
Sud (Douala -- Yaoundé, 260 km) & Voie métrique, franchissement de la Sanaga. Coût : 85 millions de Marks-or. & Caoutchouc, ivoire, puis cacao. Ravitaillement postes administratifs/militaires. \\ \hline
Port de Douala (aménagé 1890-1914) & Quais, entrepôts, grue à vapeur, dragage chenal Wouri. & 95\% du commerce extérieur. Export 1913 : 42\,000 t ; Import 1913 : 38\,000 t. \\ \hline
\end{tabularx}
\end{center}
\emph{Source : Rapports annuels du Gouvernement du Kamerun, 1913.}

\medskip
\textbf{Document 3 : Statistiques des exportations camerounaises vers l'Allemagne (en tonnes).}
\begin{center}
\begin{tabular}{|c|cccc|}\hline
\rowcolor{popblueL}
\textbf{Produit} & \textbf{1900} & \textbf{1905} & \textbf{1910} & \textbf{1913} \\ \hline
Caoutchouc sauvage & 1\,200 & 2\,800 & 1\,500 & 800 \\
Cacao & 200 & 1\,500 & 6\,200 & 11\,500 \\
Café & -- & 50 & 300 & 1\,200 \\
Huile de palme / Amandes & 8\,500 & 12\,000 & 18\,000 & 22\,000 \\
Ivoire & 45 & 30 & 15 & 8 \\ \hline
\textbf{Total exportations} & \textbf{10\,500} & \textbf{16\,500} & \textbf{26\,500} & \textbf{36\,000} \\ \hline
\end{tabular}
\end{center}
\emph{Note : Le caoutchouc sauvage décline (épuisement, concurrence asiatique) au profit des cultures de plantation (cacao, café) imposées par l'administration. Source : Statistiques du Reichskolonialamt.}

\medskip
\textbf{Document 4 : Témoignage oral attribué à Rudolf Duala Manga Bell (1910), roi des Bell, avant son exécution (1914).}
\og Nous avons signé le traité pour être des partenaires, pas des sujets sans droits. Aujourd'hui, nos terres sont arrachées pour les plantations allemandes (Société des Plantations de la Sanaga), nos hommes meurent sur les chantiers du chemin de fer (taux de mortalité 15-20\% par an), l'impôt de capitation nous ruine. L'école allemande forme des commis, pas des chefs. Si la modernité est ce train qui nous écrase, nous la refusons. \fg
\emph{Source : Tradition orale Duala, recueillie par H. N. Mougoue, \emph{Histoire du Cameroun}, 1984.}

\medskip
\textbf{Document 5 : Description de la carte schématique du \emph{Kamerun} (1911) — \emph{Le Kamerun utile}.}
La carte montre une concentration extrême des infrastructures et de l'administration sur la façade atlantique et le long des deux axes ferroviaires (Nord vers Nkongsamba, Sud vers Yaoundé). La \emph{zone de peuplement dense et d'économie de plantation} (cacao, café, hévéa) ne dépasse pas 100 km de large de part et d'autre des rails. Le Nord (Adamaoua, Lac Tchad) et l'Est (forêt dense) sont quadrillés par des postes militaires et des pistes, sans rails ni écoles, organisés en \emph{chasses à l'homme} pour le recrutement forcé de porteurs et travailleurs. Douala concentre 80\% de la population européenne, les banques, l'hôpital, l'imprimerie, le siège du gouvernement.

\courssub{Consignes}
En vous appuyant \emph{uniquement} sur les documents du dossier et les connaissances de la leçon, traitez les questions suivantes :

\begin{enumerate}
    \item \textbf{Identification des besoins industriels (Savoir) :} À partir des documents 2 et 3, citez trois matières premières camerounaises essentielles pour l'industrie allemande (\cle{Deuxième Révolution Industrielle}) et précisez pour chacune son usage industriel précis (ex. : chimie, alimentation, isolation).
    \item \textbf{Analyse des infrastructures (Savoir-faire : analyse doc) :} Les documents 2 et 5 montrent que le chemin de fer ne dessert pas l'ensemble du territoire. Expliquez \textbf{pour qui} et \textbf{pour quoi} ces infrastructures ont été construites. Montrez en quoi le tracé ferroviaire structure l'espace colonial (\emph{Kamerun utile} vs \emph{Kamerun inutile}).
    \item \textbf{Production graphique (Savoir-faire : croquis) :} Réalisez un croquis commenté et légendé du \emph{Kamerun utile} (format A4 paysage). Y figureront : le trait de côte, le fleuve Wouri, Douala, les deux lignes de chemin de fer (Nord/Sud) avec leurs terminus, la zone de plantations (hachures), les postes administratifs principaux, une flèche d'exportation vers l'Europe, l'échelle et l'orientation. Titre obligatoire.
    \item \textbf{Rédaction argumentée (Savoir-faire : rédaction) :} Rédigez une réponse structurée (introduction, deux parties équilibrées, conclusion) d'environ \textbf{20 lignes} à la question : \og L'industrialisation allemande a-t-elle profité au Cameroun ? \fg Vous distinguerez explicitement le \textbf{court terme} (contraintes, violence, prédation) et le \textbf{long terme} (héritage infrastructurel, insertion dans l'économie monétaire, formation d'une élite).
\end{enumerate}

\courssub{Ressources à mobiliser}
\begin{definition}[Révolutions industrielles et besoins]
\textbf{Première (vers 1780-1850) :} Charbon, Fer, Vapeur $\rightarrow$ Textile, Métallurgie lourde, Chemins de fer.
\textbf{Deuxième (1850-1914) :} Acier, Électricité, Pétrole, Chimie $\rightarrow$ Chimie organique (caoutchouc synthétique plus tard, mais besoin immédiat de caoutchouc naturel pour isolation câbles, pneus), Industries alimentaires (cacao, huile de palme pour savonnerie/glycérine), Armement (ivoire, caoutchouc).
\end{definition}

\begin{definition}[Impérialisme : fondements (Hobson, Lénine, programme MINESEC)]
\begin{itemize}
    \item \textbf{Économiques :} Surproduction $\rightarrow$ recherche débouchés ; Suraccumulation capitaux $\rightarrow$ placement rentable (infrastructures coloniales) ; Besoin matières premières bon marché.
    \item \textbf{Politiques/Stratégiques :} Prestige national, équilibre des puissances, bases navales.
    \item \textbf{Idéologiques :} Mission civilisatrice, darwinisme social, racisme scientifique.
\end{itemize}
\end{definition}

\begin{propriete}[Économie coloniale type \emph{colonie d'exploitation}]
Structure extravertie : Production de biens primaires (non transformés sur place) $\rightarrow$ Exportation vers la métropole $\rightarrow$ Importation de biens manufacturés. Infrastructures \emph{en éventail} (port $\rightarrow$ intérieur) et non \emph{en maillage} (intégration régionale). Main-d'œuvre contrainte (travail forcé, impôt de capitation, recrutement).
\end{propriete}

\begin{methode}[Rédaction argumentée type Probatoire (20 lignes)]
\begin{enumerate}
    \item \textbf{Introduction (2-3 lignes) :} Accroche (chiffre/clé), définition des termes, thèse nuancée (double bilan), annonce du plan (Court terme / Long terme).
    \item \textbf{Développement Partie 1 -- Court terme : L'exploitation prédatrice (8-10 lignes) :} Violence conquête, travail forcé (chemins de fer), accaparement terres (plantations), fiscalité lourde, destruction chefferies résistantes (Manga Bell). Appui docs 1, 2, 4.
    \item \textbf{Développement Partie 2 -- Long terme : Un héritage ambigu (8-10 lignes) :} Infrastructures (rail, port) réutilisées post-1916, initiation cultures rentes (cacao/café = base économie actuelle), émergence élites lettrées (écoles missionnaires/allemandes), unification administrative embryonnaire.
    \item \textbf{Conclusion (2 lignes) :} Bilan synthétique : profit quasi exclusif pour l'Allemagne (1884-1914), externalités positives récupérées plus tard par les Camerounais.
\end{enumerate}
\end{methode}

\courssub{Démarche et corrigé commenté}
\begin{exoresolu}[Corrigé détaillé de l'activité d'intégration]
\textbf{Consignes 1 à 4 : Résolution guidée.}
\tcblower
\textbf{1. Identification des besoins industriels allemands (Doc 2, 3 + Cours).}
\begin{itemize}
    \item \textbf{Caoutchouc naturel (Doc 3 : 1\,200 t en 1900) :} Indispensable à l'industrie chimique et électrique naissante (\cle{Deuxième Révolution Industrielle}) pour l'\textbf{isolation des câbles électriques} (développement réseaux urbains, tramways, téléphone) et la fabrication des \textbf{pneumatiques} (essor automobile : Daimler, Benz ; vélos). L'Allemagne n'a pas encore le caoutchouc synthétique (Bayer, 1910-1915).
    \item \textbf{Cacao (Doc 3 : 11\,500 t en 1913) :} Matière première de l'\textbf{industrie agro-alimentaire} (chocolaterie : Stollwerck, Sarotti) et de la \textbf{chimie fine} (beurre de cacao pour pharmaceutique, cosmétique). Répond à la demande de masse urbaine allemande.
    \item \textbf{Huile de palme / Amandes de palme (Doc 3 : 22\,000 t en 1913) :} Base de l'\textbf{industrie du savon} (Lever Brothers, usines allemandes type Henkel) et de la \textbf{glycérine} (explosifs, nitroglycérine pour l'armement et les mines). Produit de la Première Révolution Industrielle (savonnerie) amplifié par la Seconde (chimie de guerre).
\end{itemize}
\emph{Justification :} Ces trois produits couvrent les piliers de la Deuxième Révolution : Électricité (caoutchouc), Chimie organique (huile de palme/glycérine), Alimentation de masse (cacao). Le fer/acier (Doc 2 : rails, locomotives) est importé d'Allemagne (Krupp, Thyssen), confirmant la division internationale du travail.

\medskip
\textbf{2. Analyse des infrastructures : Pour qui ? Pour quoi ? (Doc 2, 5).}
\begin{itemize}
    \item \textbf{Pour qui ?} Prioritairement pour les \textbf{compagnies concessionnaires et l'État allemand} (rentabilité capitaux investis : 45 puis 85 millions Marks-or). Secondairement pour l'\textbf{administration militaire} (contrôle territoire, Doc 5 : pistes Nord/Est). Les populations locales sont \textbf{usagers contraints} (portage, corvée, impôt payé en nature/travail).
    \item \textbf{Pour quoi ?} \textbf{Évacuer la production} (cacao, café, caoutchouc) vers le port de Douala $\rightarrow$ Hambourg/Bremen. \textbf{Ravitiller l'intérieur} (postes militaires, personnels européens, matériel). \textbf{Ne pas intégrer} les régions entre elles (pas de liaison Est-Ouest, pas de rail vers le Nord/Est).
    \item \textbf{Structuration de l'espace (Doc 5) :} Le tracé \emph{en Y} (Douala $\rightarrow$ Nkongsamba / Douala $\rightarrow$ Yaoundé) définit le \emph{Kamerun utile} (zone de rente, 15\% du territoire, 60\% de la population, 95\% des recettes). Le reste (\emph{Kamerun inutile}) est zone de main-d'œuvre (chasse aux porteurs) et de pacification militaire, sans investissement productif. C'est l'archétype de l'infrastructure \emph{en éventail} coloniale.
\end{itemize}

\medskip
\textbf{3. Croquis commenté : \emph{Le Kamerun utile} (1911) -- Guide de réalisation.}
\begin{enumerate}
    \item \textbf{Support :} Feuille A4 paysage. Cadre au crayon gris.
    \item \textbf{Fond de carte :} Contour simplifié du Kamerun (polygone approximatif : base côtière large, élargissement Nord, pointe Est). Fleuve Wouri (Douala) + Sanaga (vers Yaoundé) + Bénoué (Nord) en traits bleus fins.
    \item \textbf{Repérage :} Douala (carré rouge + nom), Nkongsamba (triangle nord), Yaoundé (triangle sud), Garoua (point Nord), Ngaoundéré (point Nord-Est).
    \item \textbf{Infrastructures :} \textbf{Ligne Nord} (trait continu épais noir, Douala-Nkongsamba). \textbf{Ligne Sud} (trait continu épais noir, Douala-Yaoundé). \textbf{Port de Douala} (symbole ancre/navire).
    \item \textbf{Zone de plantations :} Hachures obliques \textcolor{popgreen}{vert} de part et d'autre des rails (largeur $\approx$ 1 cm de part et d'autre = 100 km réel). Légende : \og Zone de cultures d'exportation (cacao, café, hévéa) \fg.
    \item \textbf{Administration :} Carrés noirs petits aux postes clés (Douala, Buea, Yaoundé, Ngaoundéré, Garoua, Mora).
    \item \textbf{Flux :} Grande flèche \textcolor{poporange}{orange} épaisse Douala $\rightarrow$ bord droit carte (Europe) : \og Export matières premières \fg. Flèche fine retour : \og Import manufacturés \fg.
    \item \textbf{Éléments obligatoires :} Titre centré : \og \textbf{Le Kamerun utile : espace ferroviaire et plantation (1911)} \fg. Échelle graphique (0 - 100 - 200 km). Flèche Nord (haut de carte). Légende organisée (Signes ponctuels, Linéaires, Surfaciques, Flux).
\end{enumerate}
\emph{Attendu graphique :} Propreté, lisibilité, respect du code graphique (ponctuel/linéaire/surfacique), légende hiérarchisée.

\medskip
\textbf{4. Rédaction argumentée : \og L'industrialisation allemande a-t-elle profité au Cameroun ? \fg (Proposition de corrigé type 20 lignes).}

\textbf{Introduction.} L'industrialisation allemande (\cle{Deuxième Révolution Industrielle}) a fait du Kamerun une \cle{colonie d'exploitation} type (1884-1914). Si elle a doté le territoire d'infrastructures lourdes, le bilan immédiat est celui d'une prédation violente ; les bénéfices structurels ne seront récupérés qu'ultérieurement par les Camerounais.

\textbf{1. Court terme (1884-1914) : Une modernité imposée par la contrainte et la violence.} L'industrialisation allemande s'est traduite par une extraction intensive de matières premières (caoutchouc, cacao, huile de palme : +240\% exportations 1900-1913, Doc 3) au profit exclusif de la métropole (chimie, alimentation, armement). Cette prédation s'est appuyée sur un appareil coercitif : travail forcé pour les chemins de fer (Doc 2 : mortalité 15-20\%/an), accaparement des terres Duala pour plantations (Sanaga), impôt de capitation contraignant à la culture de rente. La résistance politique a été brisée (exécution de Rudolf Duala Manga Bell, 1914, Doc 4), prouvant que le \og partenariat \fg du traité de 1884 (Doc 1) a laissé place à la spoliation. Les infrastructures (rail, port) sont des \textbf{outils d'évacuation} (éventail vers Douala), non d'intégration nationale.

\textbf{2. Long terme (post-1916) : Un héritage infrastructurel et institutionnel réapproprié.} Paradoxalement, le réseau ferré (Nord/Sud) et le port de Douala, bien que conçus pour l'export, sont devenus l'\textbf{ossature logistique} de l'État camerounais post-indépendance (transfert cacao/café, désenclavement Ouest). L'introduction forcée des cultures pérennes (cacao, café, hévéa) a créé la \textbf{base de l'économie monétaire} actuelle (premières recettes d'exportation). L'école allemande puis missionnaire a formé les premiers cadres (fonctionnaires, pasteurs, syndicalistes) qui porteront le nationalisme. L'unification administrative embryonnaire (un territoire, une capitale Yaoundé, une langue administrative) a préfiguré l'unité nationale.

\textbf{Conclusion.} L'industrialisation allemande n'a \textbf{pas profité} aux Camerounais de 1884 à 1914 : elle a coûté des vies, des terres et la souveraineté. En revanche, les \textbf{externalités positives} (rails, port, cultures rentes, élites) ont constitué un \textbf{capital initial} que les Camerounais ont su réinvestir après 1916, transformant l'outil du colonisateur en levier de leur souveraineté économique et politique.
\end{exoresolu}

\courssub{Bilan de l'activité}
Cette activité a permis de vérifier l'acquisition des compétences cibles du Probatoire :
\begin{itemize}
    \item \textbf{Mobilisation des savoirs :} L'élève relie concrètement les besoins de la \cle{Deuxième Révolution Industrielle} (caoutchouc/électricité, cacao/alimentation, huile/chimie) à la politique coloniale allemande au Kamerun.
    \item \textbf{Analyse critique de sources :} Le croisement du traité (juridique), des stats (quantitatif), du témoignage (subjectif/résistant) et de la carte (spatial) permet de déconstruire le mythe de la \og mission civilisatrice \fg pour révéler la rationalité capitaliste (investissement $\rightarrow$ rentabilité).
    \item \textbf{Compétence graphique :} Le croquis du \emph{Kamerun utile} synthétise spatialement la dichotomie centre/périphérie inhérente à l'impérialisme économique.
    \item \textbf{Écriture argumentée :} La distinction explicite \emph{court terme / long terme} structure la pensée historique, évite le manichéisme et répond à l'exigence de nuance du jury. L'élève comprend que l'histoire ne se juge pas seulement à l'aune des intentions, mais des dynamiques longues.
\end{itemize}
\emph{Prolongement possible :} Comparer avec le mandat français (1919-1960) : continuité des infrastructures, rupture sur le travail forcé (code du travail indigène 1925), extension du rail vers le Nord (Ngaoundéré 1974 seulement).

\newpage

\coursec{Méthodes et exercices résolus}

\coursec{Leçon 01 : La reconstruction de l'Europe et les révolutions industrielles (1815-1914)}

\courssub{Synthèse des savoirs indispensables : L'Europe de 1815 à 1914}

\begin{aretenir}[I. L'Europe sort des guerres : reconstruction politique et ordre nouveau (1815-1848)]
\textbf{1. Le Congrès de Vienne (1814-1815) et le « Concert européen » :} Les vainqueurs de Napoléon (Autriche \cle{Metternich}, Royaume-Uni \cle{Castlereagh}, Prusse \cle{Hardenberg}, Russie \cle{Alexandre Ier}, France \cle{Tallyrand}) redessinent la carte pour encercler la France (Royaume des Pays-Bas, Royaume de Sardaigne renforcé, Confédération germanique) et restaurent les monarchies légitimes (\cle{Principes de légitimité} et d'\cle{équilibre}). Création de la \cle{Sainte-Alliance} (Autriche, Prusse, Russie) pour défendre l'ordre monarchique et religieux contre le libéralisme et le nationalisme.
\textbf{2. L'Europe des libéraux et des nations (1815-1848) :} Vague révolutionnaire : 1820 (Espagne, Naples), 1830 (France \cle{Trois Glorieuses}, Belgique indépendance, Pologne), 1848 (\cle{Printemps des peuples} : Paris, Vienne, Berlin, Francfort, Milan, Venise). Revendications : constitutions, suffrage, unité nationale (Allemagne, Italie), indépendance (Hongrie, Tchèques).
\textbf{3. Bilan 1848 :} Échec immédiat des révolutions (répression, divisions libéraux/populaires), mais avènement du suffrage universel masculin (France), abolition définitive de l'esclavage (France 1848), et prise de conscience nationale. L'Allemagne et l'Italie restent divisées ; l'Autriche survit.
\end{aretenir}

\begin{aretenir}[II. La Première Révolution Industrielle : le charbon, le fer et la vapeur (vers 1780-1850)]
\textbf{1. Berceau : le Royaume-Uni.} Facteurs : capital commercial (triangulaire), main-d'œuvre (enclosures, démographie), charbon/fer abondants, stabilité politique (1688), empire colonial (coton, marchés).
\textbf{2. Ruptures techniques :} \cle{Coke} (Darby 1709) $\rightarrow$ fer ; \cle{Machine à vapeur} à double effet (Watt 1769) $\rightarrow$ énergie mobile/continue ; \cle{Mécanisation textile} (mule-jenny, métier Jacquard) $\rightarrow$ usine.
\textbf{3. Transport :} Locomotive (Stephenson 1829), bateau à vapeur $\rightarrow$ unification marchés, baisse coûts.
\textbf{4. Conséquences sociales :} \cle{Prolétarisation}, paupérisme, travail femmes/enfants, luddisme, premières lois (UK 1833, 1842), naissance conscience de classe (Chartisme, Marx/Engels \emph{Manifeste} 1848).
\textbf{5. Diffusion :} Belgique (premier pays continental, 1830), France (Second Empire : rail, banques), Prusse (Zollverein 1834).
\end{aretenir}

\begin{aretenir}[III. La Deuxième Révolution Industrielle : acier, électricité, pétrole, chimie (1850-1914)]
\textbf{1. Nouvelles énergies/matières :} \cle{Acier} (Bessemer 1855, Thomas 1878 pour minerai phosphoreux, Martin) $\rightarrow$ rails, coques, armement. \cle{Électricité} (Gramme, Edison, Tesla) $\rightarrow$ éclairage, force motrice décentralisée, communications (télégraphe, téléphone). \cle{Pétrole} (moteur Diesel, essence) $\rightarrow$ automobile, aviation. \cle{Chimie} (teintures, engrais Haber-Bosch, explosifs, plastiques).
\textbf{2. Nouveaux leaders :} \cle{USA} (ressources, marché intérieur, Fordisme 1913), \cle{Allemagne} (chimie, électricité, banques d'affaires, protectionnisme 1879), France, Japon (Meiji), Russie (Witte).
\textbf{3. Organisation :} \cle{Taylorisme} (organisation scientifique, chronométrage), \cle{Fordisme} (chaîne, standardisation, salaires \emph{\$5/jour}), concentration capitalistique (\cle{Trusts}, \cle{Cartels}, \cle{Holdings}), fusion banques/industries = \cle{Capital financier} (Hilferding, Lénine).
\textbf{4. Société :} Amélioration relative niveau de vie (salaires réels $\uparrow$, consommation de masse), grèves de masse, \cle{Deuxième Internationale} (1889), législation sociale (Allemagne 1883-1889 : assurance maladie, accidents, vieillesse).
\end{aretenir}

\begin{aretenir}[IV. Les fondements de l'impérialisme européen : pourquoi l'Afrique ? (1870-1914)]
\textbf{1. Causes économiques (Thèse Hobson/Lénine) :} Surproduction, suraccumulation de capital $\rightarrow$ recherche \cle{débouchés}, \cle{matières premières} (caoutchouc, cuivre, étain, pétrole, coton, huile de palme), placement capitaux (chemins de fer, mines). \cle{Division internationale du travail (DIT)} : Centre (industrie) / Périphérie (matières premières).
\textbf{2. Causes politiques/stratégiques :} Rivalités puissances (prestige, \cle{balance of power}), points stratégiques (Suez, Cap, Détroits), sécurité des métropoles (Bismarck : colonies « pour distraire l'opinion » ou protéger commerçants). \cle{Conférence de Berlin} (1884-1885) : règles du partage (effectivité, libre-échange bassins Congo/Niger, abolition traite).
\textbf{3. Causes idéologiques :} \cle{Darwinisme social} (lutte pour la vie entre races/nations), \cle{Mission civilisatrice} (J. Ferry, Kipling « fardeau de l'homme blanc »), racisme scientifique, rôle presse/ligues coloniales/geographes.
\textbf{4. Le partage :} 1880 : 80\% Afrique libre. 1914 : 2 États indépendants (Libéria, Éthiopie). Modes : protectorats, concessions, colonies de peuplement (Algérie, Afrique du Sud, Kenya, Algérie).
\end{aretenir}

\begin{aretenir}[V. Synthèse : L'Europe, atelier du monde et gendarme colonial (Bilan 1914)]
\textbf{1. Hégémonie économique :} Europe + USA + Japon = 80\% production industrielle mondiale. Londres = place financière mondiale (Livre sterling étalon-or). Contrôle flux commerciaux (flottes, câbles sous-marins \emph{All Red Line}, canaux Suez/Panama).
\textbf{2. Domination politique :} Partage formel (Afrique, Asie, Pacifique) + impérialisme informel (Chine, Empire ottoman, Amérique latine). \cle{Diplomatie du canon} / \cle{Gunboat diplomacy}.
\textbf{3. Contradictions fatales :} Dépendance aux importations (blé, coton, caoutchouc, pétrole, métaux). Tensions sociales internes (grèves, antimilitarisme vs nationalisme). Rivalités inter-impérialistes (Maroc 1905/1911, Balkans 1908-1913, course navale Dreadnought) $\rightarrow$ Système d'alliances (Triplice / Triple Entente) $\rightarrow$ \cle{Guerre totale 1914-1918}.
\end{aretenir}

% --- MÉTHODES ET EXERCICES CORRIGÉS ---

\courssub{Méthode 1 : Analyser un document iconographique (caricature, affiche, photographie)}

\begin{methode}[Analyse d'un document iconographique : la démarche en 4 temps]
\textbf{Objectif :} Appliquer les notions du chapitre (industrialisation, impérialisme) à l'analyse d'une source visuelle pour en extraire le sens historique.
\begin{enumerate}
    \item \textbf{Identification (La « fiche d'identité ») :} Nature (caricature, affiche de propagande, photo de reportage, tableau), titre, auteur, date, source, contexte historique précis (à quel événement ou période du chapitre cela se rapporte-t-il ?).
    \item \textbf{Description objective (Ce que je vois) :} Décrire les éléments visuels sans interpréter : personnages (physionomie, vêtements, attitude, symboles), objets, décor, composition (premier plan / arrière-plan), couleurs, texte/légende intégrée. Utiliser le vocabulaire technique (\cle{allégorie}, \cle{stéréotype}, \cle{mise en scène}).
    \item \textbf{Analyse et Interprétation (Ce que cela veut dire) :} Relier la description aux connaissances du cours.
    \begin{itemize}
        \item Quel est le \cle{message} de l'auteur ? (Critique, justification, glorification, dénonciation).
        \item Quels \cle{arguments} visuels appuient ce message ? (Exagération, contraste, symboles).
        \item Quel lien avec les notions : \cle{révolution industrielle} (progrès technique, condition ouvrière, pollution), \cle{impérialisme} (mission civilisatrice, exploitation, conférence de Berlin, résistance africaine) ?
    \end{itemize}
    \item \textbf{Synthèse critique (Conclusion) :} Rédiger une phrase de conclusion qui replace le document dans son contexte, en évalue la portée (témoignage, propagande) et ses limites (subjectivité, public visé).
\end{enumerate}
\textbf{Reflexe Bac :} Toujours citer le document explicitement : « Le document montre que... », « Dans le coin supérieur droit, on observe... ». Ne jamais faire d'histoire sans s'appuyer sur le document.
\end{methode}

\begin{exoresolu}[Analyse d'une caricature : « Le partage de l'Afrique » (vers 1885)]
\textbf{Énoncé :} Vous présentez ci-après une caricature parue dans la presse européenne en 1885. Elle représente Otto von Bismarck coupant un gâteau en forme de carte de l'Afrique avec un grand couteau, entouré de représentants des puissances européennes (France, Grande-Bretagne, Portugal, Belgique) qui tendent leurs assiettes.
\begin{enumerate}
    \item Identifiez le document (nature, date, auteur probable, contexte).
    \item Décrivez la scène et les symboles utilisés.
    \item Expliquez le message de l'auteur en le reliant aux \cle{fondements de l'impérialisme européen} étudiés en cours (économiques, politiques, idéologiques).
    \item Concluez sur la portée historique de ce document pour comprendre la \cle{Conférence de Berlin}.
\end{enumerate}
\tcblower
\textbf{Solution détaillée :}
\begin{enumerate}
    \item \textbf{Identification :}
    \begin{itemize}
        \item \textbf{Nature :} Caricature de presse (dessin satirique politique).
        \item \textbf{Date :} 1885 (année de l'Acte final de la Conférence de Berlin).
        \item \textbf{Auteur :} Inconnu (souvent anonyme dans la presse d'époque), publié dans un journal européen (ex: \emph{Punch} à Londres ou \emph{Le Charivari} à Paris).
        \item \textbf{Contexte :} Fin de la Conférence de Berlin (nov. 1884 – fév. 1885) qui a officialisé le partage de l'Afrique entre puissances européennes sans aucune représentation africaine. Bismarck en est l'hôte et l'arbitre.
    \end{itemize}
    \item \textbf{Description :}
    \begin{itemize}
        \item \textbf{Personnage central :} Bismarck, vêtu d'un uniforme militaire ou costume sombre, tenant un grand couteau (symbole de l'arbitraire et de la violence du découpage). Il est actif, il « tranche ».
        \item \textbf{Objet central :} Un gâteau géant dont la forme est le continent africain (reconnaissable à sa géographie : golfe de Guinée, Corne de l'Afrique, Cap de Bonne-Espérance). Le gâteau symbolise une proie appétissante, une ressource à consommer.
        \item \textbf{Personnages périphériques :} 4 à 5 hommes en costumes civils ou uniformes (représentants de la France, UK, Portugal, Belgique, Italie/Allemagne). Ils sont passifs, assis, tendent des assiettes vides, l'air gourmand ou impatient. Ils attendent leur part.
        \item \textbf{Absence notable :} Aucun Africain n'est représenté à table. Le continent est traité comme une chose (\cle{res nullius}).
    \end{itemize}
    \item \textbf{Analyse / Interprétation (Lien avec les fondements de l'impérialisme) :}
    \begin{itemize}
        \item \textbf{Fondement politique / diplomatique :} La scène illustre la \cle{diplomatie du « partage pacifique »} voulue par Bismarck pour éviter les guerres entre Européens (ex: crise du Congo, Fachoda évitée). Le couteau « légalise » le découpage par le droit international (Acte de Berlin).
        \item \textbf{Fondement économique :} Le gâteau = richesses (caoutchouc, ivoire, minerais, huile de palme, main-d'œuvre). Les assiettes = appétits capitalistes, recherche de \cle{débouchés} et de \cle{matières premières} pour la \cle{Deuxième Révolution Industrielle} (acier, chimie, pneus).
        \item \textbf{Fondement idéologique :} L'absence d'Africains et la métaphore culinaire (manger un gâteau) traduisent le mépris racial et la théorie de la \cle{mission civilisatrice} (ou plutôt son hypocrisie) : l'Afrique est un vide à remplir, une « terre sans maître » offerte à la « civilisation » européenne. C'est la vision du \cle{darwinisme social} appliqué aux nations.
    \end{itemize}
    \item \textbf{Synthèse critique :}
    Ce document est une \cle{source critique} précieuse : il dévoile la réalité brutale du partage (violence, appétit, exclusion des Africains) que le langage diplomatique officiel (Acte de Berlin : « libre commerce », « civilisation », « abolition de l'esclavage ») masque. Il témoigne de la vision eurocentrée du monde en 1885. Sa limite : c'est une vision européenne (regard du prédateur), pas le vécu des populations colonisées.
\end{enumerate}
\textbf{Conclusion type Bac :} Cette caricature résume par l'image la conférence de Berlin : un partage arbitraire de l'Afrique motivé par les intérêts économiques de la Deuxième Révolution Industrielle et la rivalité politique entre puissances, légitimé par une idéologie raciste, sans consentement des peuples africains.
\end{exoresolu}

\courssub{Méthode 2 : Rédiger un paragraphe argumenté structuré (Causes / Fait / Conséquences)}

\begin{methode}[Rédaction d'un paragraphe argumenté historique : la structure C-F-C]
\textbf{Objectif :} Rédiger et justifier une réponse historique cohérente en mobilisant des connaissances précises (dates, acteurs, concepts) pour expliquer un phénomène (ex: Pourquoi l'industrialisation ? Pourquoi le partage de l'Afrique ?).
\begin{enumerate}
    \item \textbf{Analyser le sujet :} Souligner les mots-clés (verbes : expliquer, montrer, analyser ; notions : \cle{révolution industrielle}, \cle{impérialisme}, \cle{prolétariat} ; bornes chronologiques).
    \item \textbf{Construire le plan (Brouillon) :} Organiser les arguments en 2 ou 3 parties logiques.
    \begin{itemize}
        \item \textbf{Modèle Causalité :} \textbf{Causes} (structurelles, conjoncturelles, déclenchantes) $\rightarrow$ \textbf{Fait majeur} (définition, chronologie) $\rightarrow$ \textbf{Conséquences} (immédiates, à long terme, positives/négatives).
        \item \textbf{Modèle Thématique :} \textbf{Argument 1} (économique) / \textbf{Argument 2} (social/politique) / \textbf{Argument 3} (technologique/idéologique).
    \end{itemize}
    \item \textbf{Rédiger le paragraphe :}
    \begin{itemize}
        \item \textbf{Phrase d'amorce (Thèse) :} Répond directement à la question, annonce le plan. Ex: « L'industrialisation de l'Europe résulte de la convergence de facteurs économiques, techniques et démographiques entre 1780 et 1850. »
        \item \textbf{Développement (Arguments + Exemples) :} Un argument = une phrase thèse + une explication + un \cle{exemple précis} (date, chiffre, nom, invention, loi). Utiliser les connecteurs logiques : \emph{En effet, par conséquent, de plus, toutefois, ainsi}.
        \item \textbf{Phrase de conclusion (Ouverture) :} Résume l'essentiel et ouvre sur la période suivante. Ex: « Ainsi, la Première Révolution Industrielle transforme profondément les structures sociales, préparant l'avènement de la Deuxième Révolution Industrielle et l'expansion coloniale. »
    \end{itemize}
    \item \textbf{Relire :} Vérifier l'orthographe des noms propres, les dates, la cohérence chronologique, l'absence de jugements de valeur anachroniques.
\end{enumerate}
\textbf{Formule à retenir :} \textbf{1 Idée directrice = 1 Paragraphe = 1 Argument + 1 Preuve (Exemple/Chiffre/Date) + 1 Lien logique.}
\end{methode}

\begin{exoresolu}[Rédaction : Les causes de la Première Révolution Industrielle en Angleterre]
\textbf{Énoncé :} « En vous appuyant sur vos connaissances, rédigez un paragraphe argumenté d'une quinzaine de lignes pour expliquer pourquoi la Première Révolution Industrielle démarre en Angleterre à la fin du XVIIIe siècle. »
\tcblower
\textbf{Solution détaillée :}
\textbf{1. Analyse \& Plan (Brouillon) :}
\begin{itemize}
    \item \textbf{Sujet :} Causes du démarrage de la 1ère RI en Angleterre (vers 1780-1820).
    \item \textbf{Plan thématique (3 piliers) :}
    \begin{enumerate}
        \item \textbf{Conditions préalables (Capital \& Travail) :} Accumulation capital (commerce triangulaire, agriculture), main-d'œuvre disponible (enclosures, démographie).
        \item \textbf{Innovations techniques \& Énergétiques :} Machine à vapeur (Watt 1769), coke (Darby), machinisme textile (mule-jenny, métier à tisser).
        \item \textbf{Cadre institutionnel \& Marché :} Stabilité politique (après 1688), libre-échange intérieur, empire colonial (débouchés, coton), réseau bancaire/transport (canaux).
    \end{enumerate}
\end{itemize}
\textbf{2. Rédaction finale :}
La Première Révolution Industrielle s'amorce en Angleterre dès les années 1780 grâce à la convergence exceptionnelle de facteurs structurels. \textbf{En premier lieu}, le pays dispose d'un capital abondant issu du commerce maritime et du commerce triangulaire, ainsi que d'une main-d'œuvre nombreuse libérée par la révolution agraire (lois des \cle{enclosures} au XVIIIe siècle) et la croissance démographique. \textbf{En second lieu}, une série d'innovations techniques majeures transforme l'appareil productif : l'utilisation du \cle{coke} par Abraham Darby (1709) pour la sidérurgie, l'invention de la \cle{machine à vapeur} à double effet par James \cle{Watt} (brevet 1769) qui libère l'industrie de la contrainte hydraulique, et la mécanisation du textile (mule-jenny de Crompton, 1779). \textbf{Enfin}, le cadre politique et commercial est favorable : la stabilité issue de la \cle{Glorieuse Révolution} (1688) garantit la propriété privée, le marché intérieur est unifié (pas de douanes intérieures), et l'Empire colonial assure l'approvisionnement en \cle{coton} brut et l'écoulement des produits manufacturés. \textbf{Par conséquent}, l'Angleterre devient dès 1820 « l'atelier du monde », dominant la production mondiale de charbon, de fer et de textiles, avant de diffuser son modèle en Belgique, France et Allemagne.
\textbf{3. Vérification :} Dates clés (1769, 1709, 1688) présentes. Concepts (\cle{coke}, \cle{enclosures}, \cle{machine à vapeur}, \cle{coton}) utilisés. Structure C-F-C respectée (Conditions $\rightarrow$ Innovations $\rightarrow$ Cadre $\rightarrow$ Résultat). Connecteurs logiques utilisés.
\end{exoresolu}

\courssub{Méthode 3 : Construire et exploiter un tableau comparatif}

\begin{methode}[Le tableau comparatif : outil pour structurer la comparaison de deux périodes]
\textbf{Objectif :} Appliquer les notions pour mettre en évidence les ruptures et continuités entre la Première et la Deuxième Révolution Industrielle (exigence fréquente au Bac : « Montrez les différences... »).
\begin{enumerate}
    \item \textbf{Choisir les critères de comparaison pertinents (colonnes) :} Énergie, Matières premières / Produits phares, Secteurs moteurs, Innovations majeures, Organisation du travail / Entreprise, Acteurs / Pays leaders, Conséquences sociales.
    \item \textbf{Remplir les cases avec des faits précis (lignes) :} Un fait par case (date, nom, chiffre, terme technique). Éviter les phrases vagues.
    \item \textbf{Analyser le tableau (La « valeur ajoutée ») :} Ne pas se contenter de le recopier. Rédiger 3 phrases de synthèse :
    \begin{itemize}
        \item \textbf{Rupture majeure :} Ex: Passage du charbon/fer à l'électricité/pétrole/acier ; passage de l'artisanat/usine à la grande usine tayloriste/fordiste.
        \item \textbf{Continuité :} Ex: Rôle central de l'État, recherche de productivité, division internationale du travail.
        \item \textbf{Conséquence directe :} Ex: Accélération de l'impérialisme (besoin de pétrole, caoutchouc, marchés).
    \end{itemize}
\end{enumerate}
\textbf{Reflexe Bac :} Si le sujet demande « Comparez... », le tableau est ton brouillon obligatoire. La copie finale doit contenir le tableau (si demandé) OU un paragraphe synthétique issu du tableau.
\end{methode}

\begin{exoresolu}[Comparaison : 1ère vs 2ème Révolution Industrielle]
\textbf{Énoncé :} Réalisez un tableau comparatif structuré opposant la Première Révolution Industrielle (vers 1780-1850) et la Deuxième Révolution Industrielle (1850-1914). Sur la base de ce tableau, rédigez une synthèse de 10 lignes mettant en évidence les deux ruptures majeures.
\tcblower
\textbf{Solution détaillée :}
\textbf{1. Tableau comparatif (Brouillon / Copie) :}
\begin{center}
\begin{tabularx}{\linewidth}{|l|X|X|}
\hline
\rowcolor{popblueL}
\textbf{Critères} & \textbf{Première Révolution Industrielle (1780-1850)} & \textbf{Deuxième Révolution Industrielle (1850-1914)} \\
\hline
\textbf{Énergie dominante} & Charbon (machine à vapeur Watt) & Électricité, Pétrole, Charbon (turbines, moteurs thermiques) \\
\hline
\textbf{Matières / Produits phares} & Fer, Coton, Laine & \cle{Acier} (Bessemer 1855, Martin 1864), Chimie (teintures, engrais, explosifs), Aluminium, Caoutchouc \\
\hline
\textbf{Secteurs moteurs} & Textile, Sidérurgie (fer), Houille, Chemin de fer (1ère génération) & Sidérurgie (acier), Chimie, Électricité, Automobile, Pétrole, Chemin de fer (réseaux denses) \\
\hline
\textbf{Innovations symboliques} & Machine à vapeur, Métier à tisser mécanique, Locomotive (Stephenson 1829) & Convertisseur Bessemer, Dynamite (Nobel), Téléphone, Ampoule (Edison), Moteur Diesel/Explosion, Ligne de montage (Ford 1913) \\
\hline
\textbf{Organisation du travail} & Usine de manufacture, Division du travail (Smith), Travail aux pièces, Prolétarisation & \cle{Taylorisme} (organisation scientifique), \cle{Fordisme} (production de masse), Concentration capitalistique (Trusts, Cartels, Holdings) \\
\hline
\textbf{Pays leaders} & \cle{Royaume-Uni} (monopole initial) & \cle{USA}, \cle{Allemagne} (dépassent UK), France, Japon, Russie (rattrapage dirigé par l'État) \\
\hline
\textbf{Conséquences sociales} & Paupérisme, Luddisme, Premières lois sociales (UK 1833/1842), Naissance syndicats & Amélioration relative niveau de vie (salaires réels $\uparrow$), Grèves de masse, Deuxième Internationale (1889), Législation sociale avancée (Allemagne Bismarck 1883) \\
\hline
\end{tabularx}
\end{center}

\textbf{2. Synthèse argumentée (Ruptures majeures) :}
\textbf{Rupture 1 : Le changement de paradigme techno-énergétique.} On passe d'une industrie \cle{« charbon-fer-vapeur »} (énergie thermique discontinue, matière première fer malléable mais fragile) à une industrie \cle{« électricité-pétrole-acier-chimie »} (énergies flexibles, transportables à distance, acier résistant et standardisé, produits chimiques de synthèse). Cela permet la \cle{deuxième révolution des transports} (automobile, aviation, métro) et la \cle{révolution des communications} (télégraphe, téléphone, radio).
\textbf{Rupture 2 : La mutation du capitalisme et de l'entreprise.} Le capitalisme de concurrence (petites usines, patron propriétaire) laisse place au \cle{capitalisme de monopole} (concentration : trusts, cartels, holdings). L'entreprise devient une \cle{grande entreprise} bureaucratisée, financée par les banques d'affaires (crédit mobilier), appliquant le \cle{taylorisme} (séparation conception/exécution) puis le \cle{fordisme} (standardisation, productivité, consommation de masse). Cette concentration économique explique l'urgence de l'expansion coloniale (recherche de matières premières spécifiques : caoutchouc, cuivre, pétrole et de marchés captifs).
\end{exoresolu}

\courssub{Méthode 4 : Exploiter des données statistiques (calculs de taux, graphiques)}

\begin{methode}[Analyse quantitative : calculer et interpréter des taux d'accroissement]
\textbf{Objectif :} Appliquer des outils mathématiques simples à des séries statistiques historiques (production, démographie, commerce) pour objectiver les ruptures (industrialisation, croissance).
\begin{enumerate}
    \item \textbf{Lire le tableau / graphique :} Identifier la grandeur mesurée (tonnes, km, \%), l'unité, la période, la source, la zone géographique.
    \item \textbf{Calculer le taux d'accroissement (variation relative) :}
    \[ \text{Taux d'accroissement (\%)} = \frac{\text{Valeur finale} - \text{Valeur initiale}}{\text{Valeur initiale}} \times 100 \]
    \item \textbf{Calculer le taux moyen annuel (si période longue) :}
    \[ \text{Taux moyen annuel (\%/an)} \approx \frac{\text{Taux d'accroissement total}}{\text{Nombre d'années}} \]
    (Ou formule exacte composée : $V_f = V_i (1+t)^n \Rightarrow t = \sqrt[n]{V_f/V_i} - 1$).
    \item \textbf{Interpréter historiquement :} Relier le chiffre au contexte. Une hausse de 300\% de la production d'acier = passage au procédé Bessemer/Thomas + demande ferroviaire + réarmement. Une baisse = crise, guerre, épuisement.
    \item \textbf{Présenter les résultats :} Tableau de synthèse avec colonnes : Période, Valeur initiale, Valeur finale, Variation absolue, Taux d'accroissement (\%), Commentaire historique.
\end{enumerate}
\textbf{Attention :} Respecter les unités (millions de tonnes vs tonnes). Arrondir au dixième de pourcent. Ne pas confondre « part dans le total » (structure) et « taux de croissance » (dynamique).
\end{methode}

\begin{exoresolu}[Analyse de la production mondiale d'acier (1870-1913)]
\textbf{Énoncé :} Le tableau ci-dessous donne la production d'acier (en millions de tonnes) des principales puissances en 1870 et 1913.
\begin{center}
\begin{tabular}{|c|c|c|}
\hline
\textbf{Pays} & \textbf{1870} & \textbf{1913} \\
\hline
Royaume-Uni & 0,5 & 7,7 \\
Allemagne & 0,3 & 17,6 \\
USA & 0,3 & 31,8 \\
France & 0,4 & 4,6 \\
Russie & 0,2 & 4,8 \\
\hline
\textbf{Monde (estim.)} & \textbf{0,7} & \textbf{76,0} \\
\hline
\end{tabular}
\end{center}
\begin{enumerate}
    \item Calculez le taux d'accroissement de la production mondiale d'acier entre 1870 et 1913.
    \item Calculez le taux d'accroissement moyen annuel sur la période.
    \item Comparez les dynamiques de l'Allemagne et du Royaume-Uni. Quelle rupture cela illustre-t-il ?
    \item Reliez ces chiffres aux fondements de l'impérialisme.
\end{enumerate}
\tcblower
\textbf{Solution détaillée :}
\textbf{1. Taux d'accroissement mondial (1870-1913) :}
\[ \frac{76{,}0 - 0{,}7}{0{,}7} \times 100 = \frac{75{,}3}{0{,}7} \times 100 \approx \textbf{10 757 \%} \]
(La production a été multipliée par environ 108,5).
\textbf{2. Taux moyen annuel (approximation linéaire) :}
Période = 43 ans.
\[ \frac{10\,757\%}{43} \approx \textbf{250 \% par an} \]
(Calcul exact composé : $t = (76{,}0/0{,}7)^{1/43} - 1 \approx 11{,}2\%$ par an. La croissance est exponentielle).
\textbf{3. Comparaison Allemagne vs Royaume-Uni :}
\begin{itemize}
    \item \textbf{UK :} $\frac{7{,}7-0{,}5}{0{,}5} \times 100 = \textbf{1 440 \%}$. Multiplié par 15,4.
    \item \textbf{Allemagne :} $\frac{17{,}6-0{,}3}{0{,}3} \times 100 = \textbf{5 767 \%}$. Multiplié par 58,7.
\end{itemize}
\textbf{Interprétation :} L'Allemagne rattrape et double le Royaume-Uni (17,6 vs 7,7 Mt en 1913). Cela illustre la \cle{Deuxième Révolution Industrielle} : l'Allemagne, entrée tardivement dans l'industrialisation, adopte directement les technologies les plus modernes (procédé Thomas pour le minerai phosphoreux lorrain, chimie de pointe, électricité), bénéficie d'un État interventionniste (Zollverein, politique ferroviaire, protectionnisme 1879) et de banques d'investissement puissantes. Le UK souffre de l'\"handicap de la première heure\" (usines vétustes, conservatisme patronal).
\textbf{4. Lien avec l'impérialisme :}
Cette explosion de la production d'acier (+10 000\% mondial) crée une demande insatiable en \cle{minerais de fer} (Lorraine, Suède, Afrique), \cle{charbon}, \cle{manganèse}, \cle{nickel}, \cle{cuivre} (électricité) et \cle{pétrole}. Les puissances industrialisées (Allemagne, USA, France, Japon) n'ont pas ces ressources sur leur sol métropolitain. La \cle{Conférence de Berlin (1885)} et le \cle{partage de l'Afrique} répondent directement à cette \cle{quête de matières premières} stratégiques pour l'industrie de l'acier et de l'armement (marine de guerre Dreadnought).
\end{exoresolu}

\courssub{Méthode 5 : Réaliser un croquis commenté / Carte mentale (Géographie historique)}

\begin{methode}[Croquis commenté : Localiser les bassins industriels et les flux impérialistes]
\textbf{Objectif :} Appliquer les notions à une situation concrète (carte) pour visualiser l'organisation de l'espace par l'industrialisation et l'impérialisme.
\begin{enumerate}
    \item \textbf{Analyser le fond de carte :} Échelle, projection, limites (Europe, Monde), repères (méridiens, parallèles, capitales, fleuves).
    \item \textbf{Construire la légende (Organisée par thèmes) :} Utiliser la sémiologie graphique standardisée.
    \begin{itemize}
        \item \textbf{Surfaces (Zones) :} Bassins houillers/industriels (hachures couleur \cle{poporange} / \cle{popblue}), Zones colonisées (couleurs par puissance : UK \cle{poppink}, France \cle{popblue}, Allemagne \cle{poporange}, Belgique \cle{popgold}, Portugal \cle{popgreen}).
        \item \textbf{Points :} Villes industrielles ($\bullet$), Ports charbonniers ($\blacktriangle$), Capitales ($\star$), Lieux de conférences (Berlin $\star$).
        \item \textbf{Lignes / Flèches :} Flux matières premières (flèche fine \cle{popmuted} vers l'Europe), Flux produits manufacturés (flèche épaisse \cle{popblue} vers colonies), Lignes de chemin de fer transcontinentaux (Transsibérien, Cap-Le Caire, Bagdad-Bahn) (trait \cle{popdark}).
    \end{itemize}
    \item \textbf{Réaliser le croquis (Propreté, précision) :} Crayons de couleur, règle, écriture lisible (stylo feutre fin). Titre en haut, Nord, Échelle, Légende encadrée à droite.
    \item \textbf{Rédiger le commentaire (5-10 lignes) :} Structure : 1. Organisation de l'espace industriel européen (pôle UK/Rhénanie/USA). 2. Hiérarchisation mondiale (Centre/Périphérie). 3. Flux d'échanges inégaux (matières premières vs produits finis). 4. Rivalités géopolitiques (points de friction : Fachoda, Maroc, Balkans).
\end{enumerate}
\textbf{Reflexe Bac :} Un croquis sans commentaire vaut 0 point. Un commentaire sans croquis vaut 0 point. Les deux sont indissociables. La légende doit être \textbf{hiérarchisée} (Titre 1, Titre 2, symboles).
\end{methode}

\begin{exoresolu}[Croquis : L'Europe industrielle et l'expansion coloniale vers 1914]
\textbf{Énoncé :} À l'aide du fond de carte monde fourni (Europe + Afrique + Amériques + Asie schématiques), réalisez un croquis commenté intitulé : \emph{« L'Europe, atelier du monde et gendarme colonial en 1914 »}.
\tcblower
\textbf{Solution détaillée :}
\textbf{1. Légende structurée (à reproduire sur la copie) :}
\begin{center}
\begin{tabular}{|l|l|l|}
\hline
\rowcolor{popblueL}
\textbf{Thème} & \textbf{Symbole} & \textbf{Information} \\
\hline
\textbf{EUROPE INDUSTRIELLE} & & \\
Bassins houillers / sidérurgiques majeurs & \textcolor{poporange}{\rule{1cm}{2pt}} Hachures orange & UK (Newcastle, Galles, Glasgow), Nord/France, Ruhr/Sarre, Haute-Silésie, Pennsylvanie (USA), Donetz (Russie) \\
Villes industrielles majeures ($\geq$ 500k hab.) & $\bullet$ Noir & Londres, Paris, Berlin, Ruhr (Essen/Dortmund), Manchester, Birmingham, Pittsburgh, Chicago, Moscou, Saint-Pétersbourg \\
Ports charbonniers / sidérurgiques & $\blacktriangle$ Noir & Cardiff, Newcastle, Rotterdam, Anvers, Hambourg, Marseille, New Orleans \\
Réseau ferroviaire dense & \textcolor{popdark}{\rule{1cm}{2pt}} Trait plein fin & Europe de l'Ouest, USA Nord-Est, Japon \\
\hline
\textbf{EXPANSION COLONIALE (1914)} & & \\
Empire Britannique & \textcolor{poppink}{\rule{1cm}{2pt}} Plein rose & Afrique (Égypte, Soudan, Afrique du Sud, Nigéria, Kenya), Inde, Canada, Australie \\
Empire Français & \textcolor{popblue}{\rule{1cm}{2pt}} Plein bleu & Afrique (AOF, AEF, Maghreb, Madagascar), Indochine \\
Empire Allemand & \textcolor{poporange}{\rule{1cm}{2pt}} Plein orange & Afrique (Cameroun, Togo, Tanzanie, Namibie), Pacifique \\
Empire Belge & \textcolor{popgold}{\rule{1cm}{2pt}} Plein or & Congo belge \\
Autres (Portugal, Italie, Espagne, Japon, USA) & \textcolor{popgreen}{\rule{1cm}{2pt}} Plein vert / Hachures & Zones respectives \\
\hline
\textbf{FLUX ET RIVALITÉS} & & \\
Flux Matières premières (Coton, Caoutchouc, Minerais, Pétrole) & $\xrightarrow{\text{---}}$ Flèche fine \textcolor{popmuted}{grise} & Colonies $\rightarrow$ Métropoles (Europe, USA, Japon) \\
Flux Produits manufacturés / Capitaux / Migrants & $\xrightarrow{\text{===}}$ Flèche épaisse \textcolor{popblue}{bleue} & Métropoles $\rightarrow$ Colonies / Mondes nouveaux \\
Projets ferroviaires impériaux inachevés & \textcolor{popdark}{\rule{1cm}{2pt}} Trait pointillé & Transafricain (Cap-Le Caire), Bagdad-Bahn, Transsibérien \\
Conférence de Berlin (1885) & $\star$ Rouge & Berlin \\
Crises diplomatiques (Fachoda 1898, Maroc 1905/11) & $\otimes$ Rouge & Lieux des crises \\
\hline
\end{tabular}
\end{center}

\textbf{2. Commentaire type (10 lignes) :}
En 1914, l'espace mondial est structuré par la domination industrielle de l'Europe et de ses prolongements (USA, Japon). Le \cle{cœur industriel} (hachures orange) forme une « mégalopole » discontinue allant de l'Écosse à la Silésie et de la Pennsylvanie au Donbass, concentrant 80\% de la production mondiale d'acier et de charbon. Cette puissance industrielle repose sur une \cle{périphérie coloniale} colorée par empires, soumise à une \cle{division internationale du travail} inégale : les flèches grises (matières premières : coton égyptien/indien, caoutchouc congolais/amazonien, minerais africains, pétrole roumain/caucasien) convergent vers les métropoles ; les flèches bleues (produits finis, capitaux, migrants, armes) en repartent. Le partage de l'Afrique (légende 1885) achève l'encerclement. Mais cette carte révèle les \cle{tensions} : rivalités franco-britanniques (Fachoda), germano-britanniques (Bagdad-Bahn, marine), austro-russes (Balkans). L'Europe est à la fois l'atelier du monde et le foyer d'une guerre imminente pour le repartage.
\end{exoresolu}

\courssub{Méthode 6 : Élaborer une problématique et un plan de dissertation (Synthèse)}

\begin{methode}[De la question de cours à la dissertation : Problématiser et Planifier]
\textbf{Objectif :} Rédiger et justifier une démarche intellectuelle complète pour traiter un sujet de synthèse type Bac (ex: « L'industrialisation de l'Europe, fondement de sa puissance mondiale (1815-1914) »).
\begin{enumerate}
    \item \textbf{Analyser le sujet (Mots-clés, Bornes, Verbe) :}
    \begin{itemize}
        \item \textbf{Mots-clés :} Industrialisation, Europe, Puissance mondiale, Fondement.
        \item \textbf{Bornes :} 1815 (Congrès de Vienne / Paix) -- 1914 (Guerre mondiale / Apogée).
        \item \textbf{Verbe :} « Montrer que », « Expliquer comment », « Dans quelle mesure ». $\rightarrow$ Exige une démonstration structurée, nuancée.
    \end{itemize}
    \item \textbf{Problématiser (La question centrale) :} Transformer le sujet en question ouverte qui contient une tension.
    \begin{itemize}
        \item \emph{Mauvaise :} « Quelles sont les révolutions industrielles ? » (Trop descriptif).
        \item \emph{Bonne :} « \textbf{En quoi la double révolution industrielle (1780-1914) a-t-elle constitué le socle matériel, technique et financier de l'hégémonie européenne, tout en engendrant les contradictions (sociales, impérialistes) qui mènent à 1914 ?} »
    \end{itemize}
    \item \textbf{Construire le plan détaillé (3 parties, 2-3 sous-parties) :} Utiliser la logique \textbf{Chronologique-Thématique} ou \textbf{Analytique (Causes/Modalités/Limites)}.
    \begin{itemize}
        \item \textbf{I. L'industrialisation, moteur de la puissance européenne (1815-1870 env.) :} A. 1ère RI (UK) : hégémonie commerciale/navale. B. Diffusion continentale (Belgique, France, Allemagne) : unification politique (Zollverein, Unité italienne/allemande) par le rail/fer.
        \item \textbf{II. La 2ème RI : accélération, concentration et nouvel impérialisme (1870-1914) :} A. Nouveaux leaders (USA, Allemagne) : acier, chimie, électricité. B. Capitalisme monopoliste $\rightarrow$ Partage de l'Afrique/Asie (matières premières, marchés). C. Puissance militaire (flottes Dreadnought, artillerie).
        \item \textbf{III. Les limites et contradictions d'une puissance fragile :} A. Dépendance aux importations (blé, coton, pétrole, caoutchouc). B. Question sociale interne (grèves, internationalisme ouvrier vs nationalisme). C. Rivalités inter-impérialistes (course aux armements, alliances) $\rightarrow$ 1914.
    \end{itemize}
    \item \textbf{Rédiger l'introduction (Modèle « Entonnoir ») :}
    \begin{enumerate}
        \item \textbf{Accroche :} Citation, chiffre clé, événement marquant (Ex: « En 1914, l'Europe contrôle 84\% de la surface du globe »).
        \item \textbf{Définition des termes / Contexte :} Définir industrialisation, hégémonie. Situer 1815-1914.
        \item \textbf{Problématique :} La question formulée à l'étape 2.
        \item \textbf{Annonce du plan :} « Nous verrons d'abord que... Puis nous analyserons... Enfin nous montrerons que... »
    \end{enumerate}
\end{enumerate}
\textbf{Formule Bac :} \textbf{Introduction = Accroche + Définitions + Problématique + Annonce du plan.} (Obligatoire, notée sur 4-5 pts souvent).
\end{methode}

\begin{exoresolu}[Sujet de dissertation type Bac : « L'Europe, atelier du monde et gendarme colonial (1815-1914) »]
\textbf{Énoncé :} Sujet : « Montrez comment l'industrialisation de l'Europe a fondé sa domination mondiale entre 1815 et 1914. » Rédigez l'introduction complète (accroche, définitions, problématique, annonce de plan) et le plan détaillé des 3 parties principales avec arguments et exemples pour chacune.
\tcblower
\textbf{Solution détaillée :}
\textbf{INTRODUCTION}
\textbf{[Accroche]} En 1815, au Congrès de Vienne, l'Europe sort exsangue de 25 ans de guerres révolutionnaires et napoléoniennes. Un siècle plus tard, en 1914, elle domine politiquement, économiquement et culturellement près de 85\% de la surface du globe, contrôlant des empires coloniaux immenses.
\textbf{[Définitions / Contexte]} L'\cle{industrialisation} est le passage d'une économie agraire et artisanale à une économie mécanisée, utilisant l'énergie fossile (charbon, puis pétrole/électricité) et la machine-outil, organisée en usines. La \cle{domination mondiale} s'exerce par le contrôle des échanges (libre-échange imposé), la colonisation (partage de l'Afrique, Asie), la puissance navale et la diffusion des normes techniques/juridiques. La période 1815-1914 couvre le « long XIXe siècle », de la restauration à la Grande Guerre.
\textbf{[Problématique]} \textbf{Comment le double processus d'industrialisation (Première puis Deuxième Révolution Industrielle) a-t-il fourni à l'Europe les moyens matériels, techniques, financiers et idéologiques de son hégémonie planétaire, tout en semant les germes de son autodestruction ?}
\textbf{[Annonce du plan]} Nous analyserons dans un premier temps comment la Première Révolution Industrielle assure la prééminence britannique et initie la pénétration informelle (I). Dans un second temps, nous montrerons que la Deuxième Révolution Industrielle renforce la puissance européenne mais la déplace vers de nouveaux acteurs (Allemagne, USA) et impose le partage formel du monde (II). Enfin, nous verrons que cette domination industrielle porte en elle ses contradictions : dépendance aux ressources, tensions sociales et rivalités guerrières (III).

\textbf{PLAN DÉTAILLÉ AVEC ARGUMENTS ET EXEMPLES}

\textbf{I. 1815-1870 : La Première Révolution Industrielle, fondement de l'hégémonie britannique et de l'expansion informelle}
\begin{itemize}
    \item \textbf{A. Le « atelier du monde » britannique :} Monopole charbon/fer/textile/vapeur. Supériorité navale (marine à vapeur, cuirassés). Livre sterling = monnaie de référence. \emph{Ex: Great Exhibition 1851 (Crystal Palace)}.
    \item \textbf{B. Diffusion continentale et construction des États-nations :} Belgique (1er pays continental industrialisé), France (Second Empire : réseaux ferroviaires, banques type Crédit Mobilier), Prusse (Zollverein 1834 $\rightarrow$ unification 1871). Le rail structure l'espace national et militaire.
    \item \textbf{C. Impérialisme informel et libre-échange imposé :} Guerres de l'Opium (1839-1842, 1856-1860) ouverture Chine/Japon (Traité de Kanagawa 1854). Traité de commerce Cobden-Chevalier (1860). Domination économique sans administration directe (sauf Algérie, Inde).
\end{itemize}

\textbf{II. 1870-1914 : La Deuxième Révolution Industrielle, nouvel impérialisme et multipolarité}
\begin{itemize}
    \item \textbf{A. Nouveaux leaders, nouvelles technologies :} Acier (Bessemer/Thomas), Chimie (Allemagne), Électricité/Pétrole (USA). Dépassement du UK (Production acier : USA 31,8 Mt, All 17,6 Mt, UK 7,7 Mt en 1913).
    \item \textbf{B. Capitalisme monopoliste et partage de l'Afrique :} Concentration (Trusts, Cartels, Krupp, Schneider, Rockefeller). Excédents de capitaux $\rightarrow$ exportation de capitaux (Hobson, Lénine). Conférence de Berlin (1885) : partage formel. \emph{Ex: Congo (Caoutchouc/Ivoire), Rand (Or), Minette lorraine}.
    \item \textbf{C. Puissance militaire et contrôle des mers :} Révolution navale (Dreadnought 1906, turbines, pétrole). Course aux armements. Canaux stratégiques (Suez 1869, Panama 1914). Contrôle des câbles sous-marins (télégraphie : « All Red Line »).
\end{itemize}

\textbf{III. Une domination fragilisée par ses contradictions internes et externes}
\begin{itemize}
    \item \textbf{A. Dépendance structurelle aux matières premières :} L'Europe ne produit pas : Coton (USA/Égypte/Inde), Caoutchouc (Congo/Amazonie/Asie), Pétrole (USA/Roumanie/Russie/Persie), Cuivre/Étain (Afrique/Amérique du Sud), Blé (USA/Argentine/Russie/Canada). Vulnérabilité stratégique (blocus 1914-18).
    \item \textbf{B. La question sociale : le revers de la médaille :} Prolétarisation massive, misère ouvrière (premier âge), puis luttes syndicales, grèves générales, montée du socialisme (IIe Internationale 1889). Législations sociales (Bismarck 1883, France 1910) mais tensions persistantes.
    \item \textbf{C. Rivalités inter-impérialistes : la course à l'abîme :} Alliances croisées (Triplice 1882, Triple Entente 1907). Crises marocaines (1905, 1911), crises balkaniques (1908, 1912-13). Course aux armements (budgets militaires $\times$ 4 entre 1870-1914). L'industrialisation a produit les armes de la destruction mutuelle (1914-1918).
\end{itemize}

\textbf{CONCLUSION (Ébauche) :} L'industrialisation a bien été le \cle{socle matériel} de la domination européenne (production, armes, transports, médecine, communication). Mais elle a créé un système mondial interdépendant et instable, où la rivalité pour les ressources et les marchés, amplifiée par la puissance industrielle elle-même, a conduit au suicide de 1914.
\end{exoresolu}

\courssub{Méthode 7 : Argumenter à partir d'une citation d'historien (Exercice de contextualisation)}

\begin{methode}[Exploiter une citation d'historien : La méthode « Auteur - Thèse - Contexte - Débat »]
\textbf{Objectif :} Rédiger et justifier une analyse critique d'un extrait historiographique (fréquent aux épreuves de composition ou commentaire).
\begin{enumerate}
    \item \textbf{Identifier l'auteur et l'œuvre :} Qui ? (Historien, époque, école : Annales, Marxiste, Libéral, Post-colonial). Quand ? (Date publication = historiographie contextuelle).
    \item \textbf{Repérer la thèse centrale (L'argument principal) :} Résumer en une phrase : « Pour X, l'impérialisme s'explique principalement par... »
    \item \textbf{Contextualiser (Le « pourquoi » de l'argument) :} À quel débat historiographique l'auteur répond-il ? (Ex: Débat sur les causes de l'impérialisme : Économique \emph{Hobson/Lénine} vs Politique \emph{Schumpeter/Gallagher-Robinson} vs Culturel \emph{Said}).
    \item \textbf{Confronter / Illustrer par le cours :} Valider, nuancer ou contester la thèse à l'aide des connaissances précises du programme (Dates, Chiffres, Exemples : Conférence de Berlin, Spécificité allemande, Cas du Congo, Guerre des Boers).
    \item \textbf{Conclure sur la portée :} Apport de l'historien, limites de son approche, actualité du débat.
\end{enumerate}
\textbf{Reflexe Bac :} Ne jamais laisser une citation « pendue ». Elle doit être \textbf{analysée}, pas seulement paraphrasée. Utiliser : « Cette affirmation de X illustre la thèse \emph{économique} de l'impérialisme... Toutefois, elle sous-estime le facteur \emph{stratégique} comme le montre l'exemple de... »
\end{methode}

\begin{exoresolu}[Analyse de citation : Lénine, \emph{L'Impérialisme, stade suprême du capitalisme} (1917)]
\textbf{Énoncé :} « L'impérialisme est le stade suprême du capitalisme... La caractéristique essentielle de l'impérialisme n'est pas la domination industrielle, mais la domination du capital financier... Le partage du monde est achevé. » (Lénine, 1917).
\textbf{Consigne :} Présentez l'auteur et sa thèse. Expliquez comment cette citation s'éclaire par l'étude de la Deuxième Révolution Industrielle et du partage de l'Afrique. Nuancez cette analyse économique par d'autres facteurs (politiques, stratégiques).
\tcblower
\textbf{Solution détaillée :}
\textbf{1. Identification \& Thèse :}
\begin{itemize}
    \item \textbf{Auteur :} Vladimir \cle{Lénine} (1870-1924), dirigeant bolchévique, écrit en exil à Zurich (1916-1917).
    \item \textbf{Œuvre :} \emph{L'Impérialisme, stade suprême du capitalisme} (publié 1917). École \cle{marxiste}.
    \item \textbf{Thèse centrale :} L'impérialisme n'est pas une politique choisie, mais une \cle{nécessité économique} structurelle du capitalisme arrivé à son stade monopoliste (fusion capital bancaire + capital industriel = \cle{capital financier}). La suraccumulation de capital dans les métropoles impose l'exportation de capitaux (plus que de marchandises) et le partage territorial du monde pour garantir des sphères d'influence exclusives (marchés, matières premières, placement de capitaux). « Le partage est achevé » $\rightarrow$ inévitablement, les puissances rivales devront se faire la guerre pour un \emph{repartage} (Guerre 14-18).
\end{itemize}
\textbf{2. Éclairage par le cours (Validation) :}
\begin{itemize}
    \item \textbf{Capital financier :} Correspond à la réalité de la \cle{Deuxième Révolution Industrielle} : fusion banques/industries (Deutsche Bank + Krupp/Siemens en Allemagne ; Crédit Lyonnais + Schneider en France ; Morgan + US Steel aux USA). Holdings, Trusts, Cartels (Rhin-Westphalie).
    \item \textbf{Exportation des capitaux :} Massive 1880-1914 (France $\rightarrow$ Russie, Empire ottoman, Amérique latine ; UK $\rightarrow$ Dominions, USA, Argentine ; Allemagne $\rightarrow$ Turquie, Roumanie). Chemins de fer (Transsibérien, Bagdad-Bahn) financés par capitaux français/allemands.
    \item \textbf{Partage achevé (1914) :} Carte de l'Afrique : 2 États indépendants (Libéria, Éthiopie) vs 1880 (80\% indépendant). Conférence de Berlin (1885) acte la fin du « vide » juridique.
    \item \textbf{Guerre pour le repartage :} Lénine prédit 1914. Crises marocaines (Allemagne vs France/UK), Balkans (Autriche vs Russie) = conflits pour le repartage des zones d'influence.
\end{itemize}
\textbf{3. Nuances / Limites (Critique historiographique) :}
\begin{itemize}
    \item \textbf{Facteur politique/stratégique (Gallagher \& Robinson, 1953) :} L'impérialisme n'est pas seulement économique. Ex: Occupation de l'Égypte (1882) pour le \cle{Canal de Suez} (route des Indes, stratégique), pas pour le coton seulement. \cle{Fachoda} (1898) : enjeu nilotique/géopolitique, pas rentable immédiatement.
    \item \textbf{Facteur idéologique / « Mission civilisatrice » :} Rôle de l'opinion publique, presse, géographes, missions, ligues coloniales (Ferry, Chamberlain, Rhodés). \cle{Darwinisme social}, racisme scientifique.
    \item \textbf{Contre-exemple économique :} Colonies souvent \emph{peu rentables} pour le capital métropolitain (coûts admin/militaires > profits). Profits concentrés dans quelques mines (Rand, Katanga, Kédougou). L'Allemagne industrielle puissante a peu de colonies ; la France a un grand empire mais capitalisme moins concentré.
    \item \textbf{Chronologie :} Le « partage » n'est pas figé en 1914 ; mandats SDN 1919 = nouveau repartage.
\end{itemize}
\textbf{4. Conclusion :}
L'analyse de Lénine reste \textbf{indispensable} pour comprendre la \cle{logique systémique} du capitalisme monopoliste (concentration, exportation de capitaux, rivalités inter-impérialistes) qui structure la période 1870-1914. Cependant, elle est \textbf{insuffisante} seule : elle néglige l'autonomie du politique (État, diplomatie, armée), les logiques stratégiques (Suez, Détroits) et la dimension culturelle/idéologique du consentement colonial. Une explication complète articule \textbf{déterminisme économique} (besoins de la 2ème RI) et \textbf{volonté politique} (puissance, prestige, sécurité).
\end{exoresolu}

\begin{attention}[Les 5 erreurs qui coûtent des points au Bac (Histoire 1re Tech)]
\begin{enumerate}
    \item \textbf{Confondre chronologie et causalité :} Écrire « La Révolution industrielle a causé la Révolution française » (inversion) ou mélanger les dates des deux RI (ex: mettre l'électricité en 1820). \textbf{Reflexe :} \cle{Repères chronologiques} obligatoires : 1769 (Watt), 1855 (Bessemer), 1885 (Berlin), 1914 (Guerre).
    \item \textbf{Le « catalogue » sans argumentation :} Lister des inventions (machine à vapeur, chemin de fer, téléphone...) sans expliquer \emph{en quoi} cela change la société, l'économie ou la géopolitique. \textbf{Reflexe :} Chaque fait cité doit répondre à la question : \emph{« Donc quoi ? »} (Conséquence).
    \item \textbf{Oublier l'échelle mondiale / La vision eurocentrée :} Traiter l'industrialisation comme un phénomène purement européen sans mentionner \cle{coton} (USA/Inde/Égypte), \cle{caoutchouc} (Congo/Amazonie), \cle{mains-d'œuvre} (coolies, esclaves, migrants), \cle{marchés} (Chine, Afrique). \textbf{Reflexe :} Systématiser \cle{Division Internationale du Travail (DIT)} : Centre (Industrie) / Périphérie (Matières premières).
    \item \textbf{Anachronisme terminologique :} Dire « Prolétariat » pour 1780 (on dit « ouvriers », « artisans »), « Capitalisme monopoliste » pour 1830 (concurrence), « Tiers Monde » pour 1900 (terme 1952, Sauvy), « Allemagne » avant 1871 (Confédération germanique / Prusse). \textbf{Reflexe :} Vocabulaire \textbf{daté} et \textbf{contextualisé}.
    \item \textbf{Négliger la rédaction / La structure :} Pas d'introduction (accroche, définitions, problématique, plan), pas de conclusion, paragraphes sans idée directrice, pas de connecteurs logiques (\emph{donc, cependant, par conséquent, en outre}). \textbf{Reflexe :} \cle{Une idée = Un paragraphe = Une preuve}. Soigner l'orthographe des noms propres (Bismarck, Bessemer, Taylor, Ford, Lénine, Hobson, Ferry, Rhodés).
\end{enumerate}
\end{attention}

\newpage

\coursec{Exercices}

\courssub{Je vérifie mes connaissances}

\begin{exercice}[QCM : repères de la leçon]
Pour chaque question, entoure la (ou les) bonne(s) réponse(s).
\begin{enumerate}
\item Le Congrès de Vienne, qui réorganise l'Europe après la chute de Napoléon, se tient en :
\begin{enumerate}
\item 1789 \item 1815 \item 1848 \item 1871
\end{enumerate}
\item La première machine à vapeur perfectionnée, utilisée dans les mines et les usines, est mise au point par :
\begin{enumerate}
\item James Watt \item Thomas Edison \item Henry Ford \item Alfred Nobel
\end{enumerate}
\item Le procédé Thomas (1878) permet de :
\begin{enumerate}
\item produire de l'électricité \item déphosphorer la fonte pour obtenir de l'acier \item fabriquer du ciment \item extraire le pétrole
\end{enumerate}
\item La conférence de Berlin, qui fixe les règles du partage de l'Afrique, se tient en :
\begin{enumerate}
\item 1815 \item 1848 \item 1884-1885 \item 1914
\end{enumerate}
\item Le traité de Douala (12 juillet 1884) fait du Cameroun un protectorat :
\begin{enumerate}
\item britannique \item français \item allemand \item belge
\end{enumerate}
\end{enumerate}
\end{exercice}

\begin{exercice}[Vrai ou faux, à justifier]
Réponds par vrai ou faux, puis justifie ta réponse en une ou deux phrases.
\begin{enumerate}
\item La Première Révolution industrielle repose essentiellement sur l'électricité et le pétrole.
\item Le chemin de fer est à la fois un produit et un moteur de l'industrialisation au XIX\textsuperscript{e} siècle.
\item Le « concert européen » désigne l'ensemble des grandes puissances qui règlent leurs différends par des conférences internationales après 1815.
\item L'Allemagne, unifiée en 1871, reste une puissance industrielle secondaire jusqu'en 1914.
\item L'impérialisme européen en Afrique ne s'explique que par des raisons économiques.
\end{enumerate}
\end{exercice}

\begin{exercice}[Définitions]
Définis chacune des notions suivantes en deux ou trois phrases, en donnant un exemple précis :
\begin{enumerate}
\item le protectionnisme ;
\item le traité de Douala ;
\item le procédé Thomas ;
\item le concert européen ;
\item l'économie d'enclave.
\end{enumerate}
\end{exercice}

\begin{exercice}[Calculs sur données industrielles]
Voici la production d'acier de l'Allemagne et de la Grande-Bretagne (en millions de tonnes) :
\begin{center}
\begin{tabularx}{\linewidth}{|l|X|X|X|}\hline
\rowcolor{popblueL} Année & 1870 & 1890 & 1913 \\ \hline
Allemagne & 0,3 & 2,2 & 17,6 \\ \hline
Grande-Bretagne & 0,3 & 3,6 & 7,7 \\ \hline
\end{tabularx}
\end{center}
\begin{enumerate}
\item Calcule le taux de variation de la production d'acier allemande entre 1870 et 1913 (en \%).
\item Calcule le taux de variation de la production britannique sur la même période.
\item En quelle année la production allemande dépasse-t-elle la production britannique ? Quelle conclusion historique peux-tu en tirer sur le rapport de forces industriel en Europe à la veille de 1914 ?
\end{enumerate}
\end{exercice}

\courssub{Je m'entraîne}

\begin{exercice}[Frise chronologique]
Sur une frise allant de 1780 à 1914, place au moins dix événements étudiés dans la leçon (inventions, traités, conférences, unifications, colonisations). Classe-les ensuite en deux couleurs : événements industriels d'un côté, événements politiques et coloniaux de l'autre. Que remarques-tu sur leur répartition dans le temps ?
\end{exercice}

\begin{exercice}[Tableau comparatif des deux révolutions industrielles]
Recopie et complète le tableau suivant :
\begin{center}
\begin{tabularx}{\linewidth}{|l|X|X|}\hline
\rowcolor{popblueL} Critères & Première Révolution industrielle & Deuxième Révolution industrielle \\ \hline
Période & & \\ \hline
Sources d'énergie & & \\ \hline
Industries dominantes & & \\ \hline
Pays leaders & & \\ \hline
Conséquences sociales & & \\ \hline
\end{tabularx}
\end{center}
\end{exercice}

\begin{exercice}[Analyse de document : production et empire]
On donne l'évolution de la production de charbon en Allemagne : 34 millions de tonnes en 1870, 190 millions en 1913. Sur la même période, l'Allemagne acquiert un empire colonial (Togo, Cameroun, Afrique du Sud-Ouest, Afrique orientale) entre 1884 et 1890.
\begin{enumerate}
\item Calcule le taux de variation de la production charbonnière allemande entre 1870 et 1913.
\item Rappelle pourquoi le charbon est une ressource stratégique à cette époque.
\item Relie l'évolution de la production sidérurgique et charbonnière allemande à l'extension de son empire colonial : quels besoins nouveaux la croissance industrielle crée-t-elle ?
\end{enumerate}
\end{exercice}

\begin{exercice}[Étude de cas comparée : Cameroun et Congo belge]
Compare la mise en valeur du Cameroun (protectorat allemand de 1884 à 1916, puis mandats franco-britanniques) et du Congo (État indépendant du Congo, propriété personnelle de Léopold II de 1885 à 1908, puis colonie belge).
\begin{enumerate}
\item Dégage deux similitudes (pense à l'économie d'enclave et au travail forcé).
\item Dégage deux différences (pense au statut administratif et aux investissements réalisés).
\item Conclus en une phrase sur la logique commune de la colonisation européenne en Afrique centrale.
\end{enumerate}
\end{exercice}

\courssub{Je me prépare au Bac}

\begin{exercice}[Dissertation structurée (type Probatoire/Bac technique)]
\textbf{Sujet :} À l'aide d'exemples précis pris en Europe et en Afrique (Cameroun), montrez comment les deux révolutions industrielles ont transformé les rapports entre l'Europe et le reste du monde entre 1780 et 1914.
\par\medskip
Consignes : rédige une introduction (accroche, définition des termes, cadrage spatio-temporel, problématique, annonce du plan), un développement en deux ou trois parties équilibrées et une conclusion. Tu mobiliseras notamment : la Première Révolution industrielle britannique, la Deuxième Révolution industrielle (acier, électricité, chimie), la conférence de Berlin, le traité de Douala et la mise en valeur du Cameroun.
\end{exercice}

\begin{exercice}[Commentaire de documents croisés (type Bac)]
\textbf{Document 1 (carte décrite) :} Carte de l'Afrique en 1914 au 1/50 000 000. On y observe que tout le continent, à l'exception de l'Éthiopie et du Libéria, est partagé entre sept puissances européennes : la France (Afrique du Nord et de l'Ouest), la Grande-Bretagne (Égypte, Soudan, Nigeria, Afrique du Sud), l'Allemagne (Togo, Cameroun, Afrique du Sud-Ouest, Afrique orientale), la Belgique (Congo), le Portugal (Angola, Mozambique), l'Italie (Libye) et l'Espagne (Maroc du Nord).
\par\medskip
\textbf{Document 2 (graphique décrit) :} Production mondiale d'acier (en millions de tonnes) : 0,5 en 1850 ; 7 en 1870 ; 28 en 1890 ; 76 en 1913. Production mondiale de charbon : 82 millions de tonnes en 1850 ; 1 340 millions en 1913.
\begin{enumerate}
\item Présente les deux documents (nature, date, objet).
\item Décris l'évolution de la production d'acier et de charbon entre 1850 et 1913 à l'aide de deux calculs.
\item Montre, à partir du document 1, que le partage de l'Afrique est achevé en 1914.
\item Explique le lien entre l'essor industriel européen (document 2) et la conquête coloniale (document 1) : matières premières, débouchés, moyens techniques (bateaux à vapeur, chemins de fer, armes).
\item Conclus : en quoi l'année 1914 marque-t-elle l'apogée de la domination européenne sur le monde ?
\end{enumerate}
\end{exercice}

\begin{exercice}[Croquis noté : l'Europe industrielle en 1914]
Réalise un croquis organisé et légendé de l'Europe industrielle en 1914.
\par\medskip
\textbf{Éléments attendus (barème indicatif sur 10 points) :}
\begin{enumerate}
\item le contour simplifié de l'Europe et le titre du croquis (1 point) ;
\item au moins trois bassins houillers : bassin de la Ruhr, Silésie, bassins britanniques (1,5 point) ;
\item au moins deux régions sidérurgiques : Ruhr-Westphalie, Lorraine (1 point) ;
\item au moins deux régions textiles : Lancashire, Flandre (1 point) ;
\item au moins trois grands ports : Londres, Hambourg, Rotterdam, Marseille (1,5 point) ;
\item au moins deux capitales financières : Londres, Paris, Berlin (1 point) ;
\item les grands axes ferroviaires reliant les bassins industriels aux ports et aux capitales (1 point) ;
\item une légende organisée et numérotée, une flèche du nord (1 point) ;
\item propreté et lisibilité (1 point).
\end{enumerate}
\end{exercice}

\newpage

\coursec{Corrigés des exercices}

\begin{corrige}[Exercice 1 — QCM : repères de la leçon]
\begin{enumerate}
\item \textbf{Réponse b : 1815.} Le Congrès de Vienne (septembre 1814 -- juin 1815) réorganise l'Europe après la chute de Napoléon : restauration des dynasties légitimes, redécoupage des frontières et création de la Sainte-Alliance.
\item \textbf{Réponse a : James Watt.} L'Écossais James Watt perfectionne la machine à vapeur (brevet de 1769, diffusion industrielle à partir des années 1780), qui devient le moteur de la Première Révolution industrielle.
\item \textbf{Réponse b : déphosphorer la fonte pour obtenir de l'acier.} Le procédé Thomas (1878) permet d'utiliser les minerais de fer phosphoreux (comme ceux de Lorraine) pour produire de l'acier en grande quantité.
\item \textbf{Réponse c : 1884-1885.} La conférence de Berlin (novembre 1884 -- février 1885), réunie à l'initiative de Bismarck, fixe les règles du partage de l'Afrique : occupation effective, libre navigation sur le Congo et le Niger.
\item \textbf{Réponse c : allemand.} Le traité de Douala du 12 juillet 1884, signé entre les chefs Duala (Bell, Akwa, Deïdo) et les représentants allemands (Nachtigal), fait du Cameroun un protectorat allemand.
\end{enumerate}
\end{corrige}

\begin{corrige}[Exercice 2 — Vrai ou faux, à justifier]
\begin{enumerate}
\item \textbf{Faux.} La Première Révolution industrielle (vers 1780-1850) repose sur le charbon, le fer et la machine à vapeur ; l'électricité et le pétrole sont les énergies de la Deuxième Révolution industrielle (1850-1914).
\item \textbf{Vrai.} Le chemin de fer est un produit de l'industrie (rails en fer puis en acier, locomotives) et un moteur de l'industrialisation : il transporte massivement charbon, matières premières et produits finis, et stimule la métallurgie par ses commandes.
\item \textbf{Vrai.} Après 1815, les grandes puissances (Grande-Bretagne, Russie, Autriche, Prusse, puis France) règlent leurs différends par des conférences internationales (congrès d'Aix-la-Chapelle, de Vérone...) : c'est le « concert européen ».
\item \textbf{Faux.} Unifiée en 1871, l'Allemagne devient la première puissance industrielle d'Europe continentale : en 1913, elle produit 17,6 millions de tonnes d'acier, plus du double de la Grande-Bretagne (7,7 millions).
\item \textbf{Faux.} L'impérialisme a des causes économiques (matières premières, débouchés, capitaux à placer), mais aussi politiques (prestige national, rivalités entre puissances), démographiques (émigration) et idéologiques (prétendue « mission civilisatrice », théories racistes).
\end{enumerate}
\end{corrige}

\begin{corrige}[Exercice 3 — Définitions]
\begin{enumerate}
\item \textbf{Le protectionnisme} est une politique économique qui protège la production nationale en taxant les produits étrangers (droits de douane) ou en limitant les importations. \emph{Exemple :} à la fin du XIX\textsuperscript{e} siècle, l'Allemagne et la France élèvent des barrières douanières pour protéger leur agriculture et leur industrie face à la concurrence britannique.
\item \textbf{Le traité de Douala} est l'accord signé le 12 juillet 1884 entre les chefs Duala (Bell, Akwa, Deïdo) et l'Allemagne, représentée par Gustav Nachtigal. Il place le Cameroun sous protectorat allemand : les chefs cèdent leurs droits de souveraineté en échange d'une protection et du maintien de leurs prérogatives locales.
\item \textbf{Le procédé Thomas} est un procédé sidérurgique mis au point en 1878 par Sidney Gilchrist Thomas : il permet de déphosphorer la fonte dans un convertisseur garni de chaux, et donc d'utiliser les minerais phosphoreux (Lorraine, Ruhr) pour produire de l'acier bon marché. Il fait de l'Allemagne la première sidérurgie d'Europe.
\item \textbf{Le concert européen} désigne le système diplomatique mis en place après le Congrès de Vienne (1815) : les grandes puissances se réunissent régulièrement en congrès pour régler leurs différends et maintenir l'équilibre européen sans guerre générale. Il fonctionne, avec des crises, jusqu'en 1914.
\item \textbf{L'économie d'enclave} est un modèle d'exploitation coloniale où la production (mines, plantations) est concentrée dans des zones limitées, reliées directement au port d'exportation par un chemin de fer, sans développer l'arrière-pays. \emph{Exemple :} au Cameroun, la ligne du Nord (Douala--Yaoundé) et les plantations de la côte servent à exporter caoutchouc, cacao et huile de palme vers l'Europe.
\end{enumerate}
\end{corrige}

\begin{corrige}[Exercice 4 — Calculs sur données industrielles]
\begin{enumerate}
\item Taux de variation de la production allemande entre 1870 et 1913 :
\[ t = \frac{17,6 - 0,3}{0,3} \times 100 = \frac{17,3}{0,3} \times 100 \approx \num{5767}\,\%. \]
La production d'acier allemande a été multipliée par environ $17,6 / 0,3 \approx 58,7$, soit une hausse d'environ \num{5767}\,\%.
\item Taux de variation de la production britannique (donnée 1870 corrigée : 0,5 Mt) :
\[ t = \frac{7,7 - 0,5}{0,5} \times 100 = \frac{7,2}{0,5} \times 100 = \num{1440}\,\%. \]
La production britannique a été multipliée par environ $7,7 / 0,5 = 15,4$, soit une hausse de \num{1440}\,\%.
\item En \textbf{1890}, la production allemande est de 2,2 millions de tonnes contre 3,6 pour la Grande-Bretagne : elle reste inférieure. En \textbf{1913}, elle atteint 17,6 millions contre 7,7 : le dépassement a donc lieu \textbf{entre 1890 et 1913} (en réalité dès le début des années 1900). \textbf{Conclusion historique :} à la veille de 1914, l'Allemagne unifiée est devenue la première puissance industrielle d'Europe, détrônant la Grande-Bretagne, pionnière de la première industrialisation. Ce renversement du rapport de forces industriel nourrit les rivalités (course aux armements, rivalité navale, concurrence coloniale) qui conduisent à la Première Guerre mondiale.
\end{enumerate}
\begin{attention}[Point de vigilance]
Le taux de variation se calcule toujours par rapport à la valeur de \emph{départ} (ici 1870) : $t = \dfrac{V_{\text{finale}} - V_{\text{initiale}}}{V_{\text{initiale}}} \times 100$. Ne pas diviser par la valeur finale. Vérifier les valeurs initiales fournies par le sujet (la Grande-Bretagne produit déjà 0,5 Mt d'acier en 1870, contre 0,3 Mt pour l'Allemagne).
\end{attention}
\end{corrige}

\begin{corrige}[Exercice 5 — Frise chronologique]
\textbf{Exemple de frise attendue} (au moins dix événements) :
\begin{itemize}
\item \textbf{Événements industriels} (couleur 1) : 1769/1785 machine à vapeur de Watt diffusée dans les usines ; 1825 première ligne de chemin de fer (Stockton--Darlington) ; 1856 procédé Bessemer (acier) ; 1878 procédé Thomas ; 1879 ampoule électrique d'Edison ; 1885 automobile de Benz ; 1903 premier avion des frères Wright.
\item \textbf{Événements politiques et coloniaux} (couleur 2) : 1815 Congrès de Vienne ; 1848 « Printemps des peuples » ; 1861 unité italienne ; 1871 unité allemande (Empire proclamé à Versailles) ; 1884 traité de Douala (protectorat allemand au Cameroun) ; 1884-1885 conférence de Berlin ; 1885 l'État indépendant du Congo attribué à Léopold II.
\end{itemize}
\textbf{Remarque attendue :} les deux séries d'événements se succèdent et s'accélèrent : la première vague industrielle (1780-1850) précède la grande vague coloniale (1880-1914), qui coïncide avec la Deuxième Révolution industrielle. L'industrialisation apparaît ainsi comme le moteur de l'expansion coloniale : l'Europe, devenue « atelier du monde », cherche matières premières et débouchés en Afrique.
\par\medskip
\textbf{Barème indicatif (sur 10) :} frise graduée et lisible (2 pts) ; au moins dix événements exacts et bien datés (4 pts) ; classement correct en deux catégories (2 pts) ; remarque pertinente sur la répartition (2 pts).
\end{corrige}

\begin{corrige}[Exercice 6 — Tableau comparatif des deux révolutions industrielles]
\begin{center}
\begin{tabularx}{\linewidth}{|l|X|X|}\hline
\rowcolor{popblueL} Critères & Première Révolution industrielle & Deuxième Révolution industrielle \\ \hline
Période & vers 1780-1850 & vers 1850-1914 \\ \hline
Sources d'énergie & charbon (houille), machine à vapeur & électricité, pétrole, charbon toujours dominant \\ \hline
Industries dominantes & textile (coton), métallurgie du fer, chemins de fer & acier (sidérurgie lourde), chimie, électricité, automobile \\ \hline
Pays leaders & Grande-Bretagne, puis Belgique et France & Allemagne et États-Unis, Grande-Bretagne dépassée \\ \hline
Conséquences sociales & naissance du prolétariat ouvrier, urbanisation, misère ouvrière & croissance du salariat, concentration des entreprises (cartels, trusts), mouvement ouvrier organisé (syndicats, partis socialistes) \\ \hline
\end{tabularx}
\end{center}
\textbf{Point de vigilance :} ne pas écrire que la Première Révolution industrielle utilise l'électricité : c'est l'erreur classique. Le charbon reste important dans les deux phases, mais la deuxième ajoute l'électricité et le pétrole.
\end{corrige}

\begin{corrige}[Exercice 7 — Analyse de document : production et empire]
\begin{enumerate}
\item Taux de variation de la production charbonnière allemande (1871-1913) :
\[ t = \frac{190 - 34}{34} \times 100 = \frac{156}{34} \times 100 \approx 458,8\,\%. \]
La production a été multipliée par environ $190 / 34 \approx 5,6$, soit une hausse d'environ 459\,\%.
\item Le charbon est une ressource stratégique car il est la principale source d'énergie de l'époque : il alimente les machines à vapeur des usines, les locomotives, les navires à vapeur, et il est indispensable à la sidérurgie (hauts fourneaux, production d'acier). Qui contrôle le charbon contrôle la puissance industrielle et militaire.
\item La croissance industrielle allemande crée des besoins nouveaux : des \textbf{matières premières} que l'Europe ne possède pas ou plus en quantité suffisante (caoutchouc, huile de palme, coton, minerais) ; des \textbf{débouchés} pour écouler les produits manufacturés ; des \textbf{territoires} où placer les capitaux excédentaires (chemins de fer, ports, plantations). L'empire colonial allemand (Togo, Cameroun, Afrique du Sud-Ouest, Afrique orientale), acquis entre 1884 et 1890, répond à ces besoins : le Cameroun fournit caoutchouc et cacao et constitue un marché pour les produits allemands. La puissance industrielle et l'expansion coloniale sont donc directement liées.
\end{enumerate}
\end{corrige}

\begin{corrige}[Exercice 8 — Étude de cas comparée : Cameroun et Congo belge]
\begin{enumerate}
\item \textbf{Deux similitudes :}
\begin{itemize}
\item Dans les deux cas, la mise en valeur repose sur une \textbf{économie d'enclave} : exploitation du caoutchouc, de l'huile de palme et des minerais dans des zones côtières ou fluviales, reliées directement aux ports d'exportation (Douala, Matadi) sans développement de l'arrière-pays.
\item Dans les deux cas, la colonisation s'appuie sur le \textbf{travail forcé} des populations africaines (réquisitions pour les plantations et les portages au Cameroun ; corvées et violences pour la récolte du caoutchouc au Congo), provoquant de lourdes pertes humaines.
\end{itemize}
\item \textbf{Deux différences :}
\begin{itemize}
\item \textbf{Statut administratif :} le Cameroun est un protectorat de l'État allemand dès 1884 (puis un mandat franco-britannique après 1916/1919), tandis que le Congo est d'abord la propriété \emph{personnelle} du roi Léopold II (État indépendant du Congo, 1885-1908) avant de devenir une colonie de l'État belge en 1908.
\item \textbf{Investissements :} l'Allemagne réalise au Cameroun des investissements d'infrastructure relativement importants (chemins de fer de Douala--Yaoundé et Bonabéri--Nkongsamba, port de Douala), tandis que le régime léopoldien privilégie le pillage rapide du caoutchouc avec un minimum d'investissements.
\end{itemize}
\item \textbf{Conclusion :} malgré des statuts différents, la colonisation européenne en Afrique centrale obéit partout à la même logique : exploiter les ressources et la main-d'œuvre africaines au profit exclusif des métropoles industrielles.
\end{enumerate}
\end{corrige}

\begin{corrige}[Exercice 9 — Dissertation structurée (type Probatoire/Bac technique)]
\textbf{Introduction attendue :} accroche (en 1914, l'Europe domine le monde industriel et colonial) ; définition des termes (révolutions industrielles : passages à la production mécanisée et à la grande industrie) ; cadrage spatio-temporel (Europe et Afrique, 1780-1914) ; problématique possible : \emph{comment l'industrialisation a-t-elle transformé l'Europe en puissance dominante et l'Afrique en espace dominé ?} ; annonce du plan.
\par\medskip
\textbf{Plan possible en deux parties :}
\begin{itemize}
\item \textbf{I. Les révolutions industrielles font de l'Europe l'atelier du monde.} Première Révolution industrielle britannique (vapeur de Watt, textile, fer, chemin de fer à partir de 1825) ; essor de la Belgique et de la France ; Deuxième Révolution industrielle (acier avec Bessemer 1856 et Thomas 1878, électricité, chimie, automobile) ; montée en puissance de l'Allemagne unifiée (17,6 Mt d'acier en 1913 contre 7,7 pour la Grande-Bretagne) ; concentration industrielle, réseaux ferroviaires, grands ports.
\item \textbf{II. La puissance industrielle débouche sur la domination du monde, notamment de l'Afrique.} Besoins nouveaux : matières premières, débouchés, placements de capitaux ; moyens techniques de la conquête (bateau à vapeur, chemin de fer, armes à feu, quinine) ; conférence de Berlin (1884-1885) et partage de l'Afrique ; cas du Cameroun : traité de Douala (12 juillet 1884), protectorat allemand, économie d'enclave (plantations, chemin de fer Douala--Yaoundé), travail forcé ; en 1914, toute l'Afrique est colonisée sauf l'Éthiopie et le Libéria.
\end{itemize}
\textbf{Conclusion attendue :} en 1914, l'industrie a fait de l'Europe le centre économique et impérial du monde, mais les rivalités qu'elle a créées (concurrence germano-britannique, conflits coloniaux) mènent à la Première Guerre mondiale.
\par\medskip
\textbf{Barème indicatif (sur 20) :} introduction complète (4 pts) ; développement structuré et équilibré (10 pts) ; exemples précis et dates exactes (3 pts) ; conclusion avec ouverture (2 pts) ; présentation et orthographe (1 pt).
\par\medskip
\textbf{Points de vigilance :} annoncer le plan en fin d'introduction ; citer au moins un exemple africain (Cameroun) ; ne pas confondre les deux révolutions industrielles ; dater précisément le traité de Douala et la conférence de Berlin.
\end{corrige}

\begin{corrige}[Exercice 10 — Commentaire de documents croisés (type Bac)]
\begin{enumerate}
\item \textbf{Présentation :} le document 1 est une carte politique de l'Afrique en 1914 (échelle 1/50 000 000) montrant le partage colonial du continent entre sept puissances européennes. Le document 2 est un graphique statistique présentant l'évolution de la production mondiale d'acier et de charbon entre 1850 et 1913.
\item \textbf{Évolution des productions :}
\begin{itemize}
\item Acier : $t = \dfrac{76 - 0,5}{0,5} \times 100 = \num{15100}\,\%$ entre 1850 et 1913 ; la production est multipliée par $76/0,5 = 152$.
\item Charbon : $t = \dfrac{1340 - 82}{82} \times 100 \approx \num{1534}\,\%$ ; la production est multipliée par $1340/82 \approx 16,3$.
\end{itemize}
Les deux productions explosent, signe d'une industrialisation mondiale menée par l'Europe.
\item \textbf{Partage achevé :} le document 1 montre qu'en 1914, seuls deux territoires échappent à la domination européenne : l'Éthiopie (victorieuse de l'Italie à Adoua en 1896) et le Libéria. Tout le reste du continent est réparti entre la France, la Grande-Bretagne, l'Allemagne, la Belgique, le Portugal, l'Italie et l'Espagne : le partage de l'Afrique, lancé par la conférence de Berlin (1884-1885), est achevé.
\item \textbf{Lien entre essor industriel et conquête coloniale :} l'industrie européenne en pleine croissance a besoin de matières premières (caoutchouc, coton, huile de palme, minerais), de débouchés pour ses produits manufacturés et de territoires où investir ses capitaux. Inversement, l'industrie fournit les moyens techniques de la conquête : bateaux à vapeur qui remontent les fleuves, chemins de fer qui pénètrent l'intérieur, armes à feu perfectionnées qui écrasent les résistances africaines. L'industrie est à la fois la cause et l'instrument de la colonisation.
\item \textbf{Conclusion :} en 1914, l'Europe concentre l'essentiel de la production industrielle mondiale et domine politiquement presque toute l'Afrique et une grande partie de l'Asie : c'est l'apogée de sa domination. Mais les rivalités industrielles et coloniales entre puissances déclenchent la même année la Première Guerre mondiale, qui entame ce déclin.
\end{enumerate}
\textbf{Barème indicatif (sur 20) :} présentation des documents (3 pts) ; calculs exacts (4 pts) ; analyse de la carte (3 pts) ; explication du lien industrie-colonisation (6 pts) ; conclusion (2 pts) ; expression (2 pts).
\par\medskip
\textbf{Points de vigilance :} toujours présenter un document avant de l'analyser (nature, date, auteur ou source, objet) ; justifier les affirmations par des chiffres tirés des documents ; vérifier les calculs (diviser par la valeur initiale).
\end{corrige}

\begin{corrige}[Exercice 11 — Croquis noté : l'Europe industrielle en 1914]
\textbf{Corrigé-type (éléments à placer correctement) :}
\begin{itemize}
\item \textbf{Bassins houillers} : bassin de la Ruhr (ouest de l'Allemagne), Silésie (à la frontière germano-polonaise, à l'est), bassins britanniques (Pays de Galles, centre et nord de l'Angleterre, bassin écossais) ; on peut ajouter le bassin du Nord-Pas-de-Calais et le bassin belge (Sambre-Meuse).
\item \textbf{Régions sidérurgiques} : Ruhr-Westphalie, Lorraine (est de la France), Silésie, centre de l'Angleterre (Birmingham, Sheffield).
\item \textbf{Régions textiles} : Lancashire (Manchester, nord-ouest de l'Angleterre), Flandre (Belgique), vallée du Rhin, Nord de la France.
\item \textbf{Grands ports} : Londres, Liverpool, Hambourg, Rotterdam, Anvers, Marseille, Le Havre.
\item \textbf{Capitales financières} : Londres, Paris, Berlin.
\item \textbf{Axes ferroviaires} : lignes reliant la Ruhr à Hambourg et Rotterdam ; Londres aux bassins britanniques ; Paris à la Lorraine et au Nord ; Berlin à la Silésie.
\end{itemize}
\textbf{Exemple de légende numérotée :} 1. bassin houiller ; 2. région sidérurgique ; 3. région textile ; 4. grand port ; 5. capitale financière ; 6. axe ferroviaire. Ne pas oublier la flèche du nord et le titre : « L'Europe industrielle en 1914 ».
\par\medskip
\textbf{Barème :} appliquer le barème de l'énoncé (10 points) : contour et titre (1 pt) ; trois bassins houillers (1,5 pt) ; deux régions sidérurgiques (1 pt) ; deux régions textiles (1 pt) ; trois ports (1,5 pt) ; deux capitales financières (1 pt) ; axes ferroviaires (1 pt) ; légende numérotée et flèche du nord (1 pt) ; propreté (1 pt).
\par\medskip
\textbf{Points de vigilance :} chaque élément dessiné doit figurer dans la légende (et réciproquement) ; vérifier la localisation relative (la Ruhr est à l'\emph{ouest} de l'Allemagne, la Silésie à l'\emph{est} ; le Lancashire au nord-ouest de l'Angleterre) ; un croquis sans titre ni nord est pénalisé.
\end{corrige}

\newpage

\coursec{Fiche bilan}

\begin{popfigure}
\begin{tikzpicture}[
    mindmap,
    grow cyclic,
    every node/.style={concept, execute at begin node=\hskip0pt},
    concept color=popblue,
    level 1/.style={level distance=4.6cm, sibling angle=72},
    level 2/.style={level distance=3.1cm, sibling angle=45},
    texte court/.style={text width=2.8cm, align=center, font=\small\sffamily}
]
\node[concept, texte court] {L'Europe 1815--1914 : reconstruction, industrialisation, impérialisme}
    child[concept color=popblue] {
        node[concept, texte court] {I. Reconstruction politique (1815--1848)}
            child[concept color=popblueL] {node[concept, texte court] {Congrès de Vienne (1815)}}
            child[concept color=popblueL] {node[concept, texte court] {Sainte-Alliance}}
            child[concept color=popblueL] {node[concept, texte court] {Printemps des peuples (1848)}}}
    child[concept color=poporange] {
        node[concept, texte court] {II. 1\iere{} Révolution industrielle (1780--1850)}
            child[concept color=poporangeL] {node[concept, texte court] {Charbon, fer, vapeur}}
            child[concept color=poporangeL] {node[concept, texte court] {Angleterre pionnière}}
            child[concept color=poporangeL] {node[concept, texte court] {Usine et prolétariat}}}
    child[concept color=popgreen] {
        node[concept, texte court] {III. 2\ieme{} Révolution industrielle (1850--1914)}
            child[concept color=popgreenL] {node[concept, texte court] {Acier, électricité, pétrole}}
            child[concept color=popgreenL] {node[concept, texte court] {Chimie, moteur à explosion}}
            child[concept color=popgreenL] {node[concept, texte court] {Taylorisme et trusts}}}
    child[concept color=poppurple] {
        node[concept, texte court] {IV. Fondements de l'impérialisme (1870--1914)}
            child[concept color=poppurpleL] {node[concept, texte court] {Causes éco., pol., idéol.}}
            child[concept color=poppurpleL] {node[concept, texte court] {Conférence de Berlin (1884--1885)}}
            child[concept color=poppurpleL] {node[concept, texte court] {Partage de l'Afrique}}}
    child[concept color=poppink] {
        node[concept, texte court] {V. Bilan 1914 : atelier du monde}
            child[concept color=poppinkL] {node[concept, texte court] {Puissances industrielles}}
            child[concept color=poppinkL] {node[concept, texte court] {Empires coloniaux}}
            child[concept color=poppinkL] {node[concept, texte court] {Tensions et guerre}}};
\end{tikzpicture}
\legende{Carte mentale de synthèse : les cinq piliers de la leçon (chronologie 1815--1914). Lecture : du centre vers l'extérieur, chaque branche correspond à une partie du plan, chaque feuille à une notion clé à retenir.}
\end{popfigure}

\begin{aretenir}[1. Définitions clés à connaître par cœur]
\begin{itemize}
    \item \cle{Congrès de Vienne (1814--1815)} : congrès qui réorganise l'Europe après la chute de Napoléon selon trois principes : la \cle{légitimité} (restauration des dynasties détrônées), l'\cle{équilibre} des puissances et le concert européen (réunions régulières des grandes puissances).
    \item \cle{Sainte-Alliance (1815)} : coalition des monarchies absolutistes (Russie, Autriche, Prusse) destinée à réprimer les révolutions libérales et nationales.
    \item \cle{Révolution industrielle} : passage progressif de l'artisanat à l'usine, de l'énergie animale et hydraulique aux énergies fossiles, entraînant une croissance économique forte et soutenue.
    \item \cle{Première révolution industrielle} : charbon, fer, machine à vapeur, industrie textile, chemin de fer ; elle naît au Royaume-Uni vers 1780 et s'étend à l'Europe continentale jusqu'en 1850.
    \item \cle{Deuxième révolution industrielle} : acier, électricité, pétrole, chimie, automobile, organisation scientifique du travail ; elle est dominée par l'Allemagne et les États-Unis entre 1850 et 1914.
    \item \cle{Taylorisme} : organisation scientifique du travail (OST) fondée sur le travail à la chaîne et la parcellisation des tâches pour augmenter la production.
    \item \cle{Trust} : regroupement de grandes entreprises d'un même secteur pour contrôler la production et les prix (capitalisme de concentration).
    \item \cle{Impérialisme} : politique d'expansion territoriale, économique et politique des puissances industrielles sur l'Afrique et l'Asie à la fin du XIX\ieme{} siècle.
    \item \cle{Conférence de Berlin (1884--1885)} : réunion de 14 puissances qui fixe les règles du partage de l'Afrique : liberté de navigation sur le Congo et le Niger, principe de l'\cle{occupation effective} des territoires revendiqués.
    \item \cle{Partage de l'Afrique} : colonisation de la quasi-totalité du continent entre 1880 et 1914 ; seuls l'Éthiopie (victorieuse à Adoua en 1896) et le Libéria échappent à la domination directe.
\end{itemize}
\end{aretenir}

\begin{aretenir}[2. Repères chronologiques et acteurs majeurs]
\begin{center}
\begin{tabularx}{\linewidth}{|>{\bfseries}c|X|X|}\hline
\rowcolor{popblueL}\textbf{Période / Événement} & \textbf{Acteurs / Inventions clés} & \textbf{Conséquences structurelles}\\\hline
1815 : Congrès de Vienne & Metternich (Autriche), Talleyrand (France), Alexandre I\ier{} (Russie) & Ordre conservateur, Sainte-Alliance, nouvelles frontières\\\hline
1820--1848 : vagues révolutionnaires & Libéraux, nationalistes, ouvriers (1848) & Chute de Metternich, constitutions, suffrage universel masculin en France\\\hline
1780--1850 : 1\iere{} révolution industrielle & Watt (machine à vapeur), Stephenson (locomotive) & Urbanisation, naissance du prolétariat, domination britannique\\\hline
1850--1914 : 2\ieme{} révolution industrielle & Bessemer (acier), Edison et Siemens (électricité), Benz et Daimler (automobile), Taylor (OST) & Production de masse, trusts, montée de l'Allemagne et des États-Unis\\\hline
1870--1914 : impérialisme & Bismarck, Jules Ferry, Léopold II (Belgique), Cecil Rhodes & Partage de l'Afrique et de l'Asie, tensions diplomatiques, système d'alliances\\\hline
1884--1885 : Conférence de Berlin & 14 puissances (Européens + États-Unis) & Acte général : libre navigation sur le Congo et le Niger, occupation effective\\\hline
\end{tabularx}
\end{center}
\end{aretenir}

\begin{attention}[Erreurs fréquentes à éviter]
\begin{itemize}
    \item Ne pas confondre les deux révolutions industrielles : la \textbf{première} repose sur le \cle{charbon}, le \cle{fer} et la \cle{vapeur} (Royaume-Uni) ; la \textbf{seconde} sur l'\cle{acier}, l'\cle{électricité}, le \cle{pétrole} et la \cle{chimie} (Allemagne, États-Unis).
    \item La Conférence de Berlin date de \textbf{1884--1885}, pas de 1885 seul ; elle ne partage pas l'Afrique territoire par territoire, elle fixe les \emph{règles} du partage.
    \item Ne pas écrire que toute l'Afrique est colonisée : l'\cle{Éthiopie} et le \cle{Libéria} restent indépendants.
    \item Le Congrès de Vienne (1815) restaure les monarchies : il ne faut pas le présenter comme une avancée démocratique.
    \item Toujours dater et situer précisément les exemples cités dans une copie : un exemple sans date perd sa valeur de preuve.
\end{itemize}
\end{attention}

\begin{methode}[Commentaire de document historique (texte, image ou carte)]
\begin{enumerate}
    \item \textbf{Identification} : nature du document, auteur, date, source, contexte historique précis.
    \item \textbf{Analyse} : dégager 2 ou 3 idées principales, relever le vocabulaire spécifique, la thèse de l'auteur, la portée et les limites du document.
    \item \textbf{Exploitation} : mobiliser les connaissances du cours (dates, acteurs, mécanismes) pour expliquer, confirmer ou nuancer le document.
    \item \textbf{Conclusion} : intérêt du document pour l'historien, ouverture sur la suite du sujet.
\end{enumerate}
\end{methode}

\begin{methode}[Dissertation / sujet de réflexion]
\begin{enumerate}
    \item \textbf{Analyse du sujet} : définir les termes, fixer les bornes chronologiques et spatiales, dégager le problème central.
    \item \textbf{Plan détaillé} : deux parties équilibrées (deux ou trois sous-parties chacune) : \emph{causes / conséquences} ou \emph{acteurs / mécanismes / bilan}.
    \item \textbf{Rédaction} : introduction (accroche, définitions, problématique, annonce du plan), développement (arguments + exemples précis datés + connecteurs logiques), conclusion (bilan synthétique et ouverture).
\end{enumerate}
\end{methode}

\begin{aretenir}[Checklist avant le Probatoire / Baccalauréat technique]
$\square$ Je sais définir : révolution industrielle, impérialisme, Congrès de Vienne, Conférence de Berlin, taylorisme, trust.\par
$\square$ Je place les 5 grandes dates repères : 1815 (Vienne), 1848 (Printemps des peuples), 1870--1871 (unifications allemande et italienne), 1884--1885 (Berlin), 1914 (début de la Première Guerre mondiale).\par
$\square$ Je distingue les deux révolutions industrielles : énergies, matières premières, pays leaders, organisation du travail.\par
$\square$ J'explique les trois types de causes de l'impérialisme (économiques, politiques, idéologiques) avec des exemples précis.\par
$\square$ Je décris le mécanisme du partage de l'Afrique (Conférence de Berlin, occupation effective, résistances africaines).\par
$\square$ Je sais analyser une carte du partage de l'Afrique (puissances, colonies, zones d'influence).\par
$\square$ Je sais commenter un texte sur les conditions ouvrières (témoignage, lois sociales, syndicalisme).\par
$\square$ Je relie industrialisation et impérialisme (besoins en matières premières, recherche de débouchés, excédents de capitaux).\par
$\square$ Je rédige une introduction de dissertation complète (accroche, définitions, problématique, plan annoncé).\par
$\square$ Je structure un développement en paragraphes argumentés (idée force + exemple précis + lien avec la problématique).\par
$\square$ Je conclus par un bilan synthétique et une ouverture pertinente (ex. : 1914, Société des Nations, décolonisation).\par
\end{aretenir}

\findecours

\end{document}
